\documentclass[10pt,a4paper]{article}
\usepackage[a4paper,margin=22mm,headheight=14pt]{geometry}
\usepackage{kotex}
\usepackage{amsmath,amssymb}
\usepackage{xcolor}
\usepackage[colorlinks=true,linkcolor=blue!55!black,urlcolor=blue!55!black]{hyperref}
\setlength{\parindent}{0pt}
\setlength{\parskip}{0.55em}
\setcounter{tocdepth}{2}
\begin{document}
\begin{titlepage}\centering\vspace*{3cm}
{\Large 경제사\par}\vspace{1.1cm}
{\Huge\bfseries 슬라이드별 학습 가이드\par}\vspace{1.1cm}
{\large 인구, 농업, 공업과 경제사 연구의 방법\par}
\vfill {\large 2026년 2학기\par}\end{titlepage}
\tableofcontents\clearpage

\section{강의의 지도}
경제사를 읽는 단위는 사건의 연도가 아니라 인구, 생산, 교환과 제도의 변화다. 총량과 1인당 성과를 구별하고, 다음 강의에서 사용할 분석 언어를 먼저 정리한다.\par

\subsection{슬라이드 01 · 경제사를 읽는 질문}
경제사는 과거의 사건을 시간순으로 외우는 과목을 넘어, \textbf{사람들이 무엇을 생산하고 어떻게 교환했으며 그 행동을 어떤 제도가 제약했는가}를 묻는다. 같은 흉작이라도 토지 소유, 조세, 시장 접근성에 따라 농민의 대응과 분배 결과가 달라진다. 따라서 사건을 볼 때 생산성·상대가격·권리 구조·집단 간 이해관계를 함께 적어 두면 이후의 인구·농업·공업 이야기가 이어진다.\par

\subsection{슬라이드 02 · 장기 역사와 분류사의 시각}
장기사적 접근에서는 하루의 가격 변동보다 수십 년에 걸친 인구 구조, 기술, 제도와 생활 수준을 살핀다. 분류사는 한 시대의 정치사를 모두 훑기보다 인구·농업·상업처럼 \textbf{같은 질문을 여러 시대에 되풀이}하는 방식이다. 예를 들어 농업 노동생산성의 변화가 도시 노동 공급에 어떤 조건을 마련했는지 비교하면, 산업혁명 이전과 이후를 한 흐름에서 이해할 수 있다.\par

\subsection{슬라이드 03 · 경제학에서 가져올 도구}
미시경제학의 한계생산, 기회비용, 수요·공급은 과거의 선택을 해석하는 출발점이다. 거시경제학의 총생산, 1인당 산출, 성장률은 시대와 지역을 비교할 때 쓰인다. 다만 오늘의 완전경쟁시장이나 계약 집행 능력을 과거에도 당연히 있었다고 가정하면 안 된다. \textbf{같은 유인도 토지 권리와 정보, 이동의 자유가 다르면 다른 결과}를 낳는다.\par

\subsection{슬라이드 04 · 교재를 함께 읽는 순서}
브로델은 일상의 물질생활과 장기 구조를, 클라크는 맬서스 체제와 산업혁명 이후의 소득 도약을, 경제사 개설서는 제도와 사건의 연대기를 각각 비춘다. 서로 다른 책의 시대 구분과 인과 설명이 꼭 일치하지는 않는다. 먼저 해당 강의의 핵심 질문을 잡고, 책에서 그 질문을 뒷받침하는 자료·반례·대안 설명을 찾는 것이 효율적이다.\par

\subsection{슬라이드 05 · 주제별 학습 지도}
앞부분의 방법론은 이후 강의에서 \textbf{무엇이 관찰이고 무엇이 해석인가}를 구분하는 기준이다. 인구는 생산을 나누는 분모이자 노동의 공급, 농업은 식량과 토지의 권리, 공업은 작업 조직과 기술의 문제로 이어진다. 이후 상업·산업혁명·금융·대공황에서도 같은 네 질문—누가 어떤 자원을 통제하는가, 어떤 거래가 가능한가, 누가 위험을 지는가, 산출과 생활 수준은 어떻게 바뀌는가—를 적용할 수 있다.\par

\subsection{슬라이드 06 · 평가에서 요구되는 설명}
경제사 서술형 답안은 연도를 많이 적는 것보다 \textbf{기제와 근거를 연결}해야 한다. “흑사병 뒤 임금이 올랐다”에서 멈추지 말고 노동 공급의 감소 → 노동의 한계생산 및 교섭력 변화 → 임금·지대의 상대가격 변화 → 영주와 농민의 제도적 대응 순으로 써 보자. 수치나 그림을 인용할 때는 분모, 기간, 지역을 밝히면 비교의 오류를 줄일 수 있다.\par

\subsection{슬라이드 07 · 수업 자료의 활용}
슬라이드의 도표에는 단위·축·표본기간이 들어 있다. 이미지를 확대해 이를 확인하고, 화살표가 \textbf{시간 순서}인지 \textbf{인과 주장}인지 구별해야 한다. 예컨대 인구 증가와 공업화가 같이 나타났다는 사실만으로 어느 쪽이 원인인지 정할 수 없다. 노동 수요, 식량 공급, 이주 제약 같은 중간 경로가 필요하다.\par

\subsection{슬라이드 08 · 자료를 읽을 때의 책임}
과거의 제도는 현재의 규범으로 즉시 재단하기보다, 당시 행위자가 실제로 가진 선택지와 강제의 정도를 먼저 확인해야 한다. 그렇다고 노예제·농노제의 폭력을 생산성 지표 하나로 정당화할 수는 없다. \textbf{효율성, 분배, 자유와 권리}는 서로 다른 평가 기준이며, 이후 강의에서 모두 중요하다.\par

\subsection{슬라이드 09 · 앞으로 이어질 핵심 질문}
“과거 경제는 왜 지금과 달랐고, 어떤 조건에서 변화했는가?”라는 질문을 염두에 두자. 기술만으로는 농업혁명과 산업혁명을 설명하기 어렵고, 제도만으로도 생산성과 인구의 변화를 설명하기 어렵다. 다음 강의에서는 사료와 경제 이론을 연결하는 여러 방법을 먼저 비교한다.\par

\section{경제사학: 사실, 이론, 반사실}
역사 자료에 경제 이론을 적용하되 자료의 생성 과정과 제도적 맥락을 함께 읽는다. 발전단계론, 성장사학, 클리오메트릭스와 신제도경제사를 차례로 연결한다.\par

\subsection{슬라이드 01 · 경제사학은 무엇을 설명하는가}
경제사학은 과거의 경제생활을 관찰 가능한 흔적에서 복원하고, 왜 그 결과가 나타났는지 설명한다. 경제학의 이론을 과거에 대입할 때에는 \textbf{그 시대의 권리·기술·정보·시장 범위}를 확인해야 한다. 현대의 가격과 국민계정이 없던 시기의 장부와 토지대장도 일관된 질문을 세우면 경제 자료가 된다.\par

\subsection{슬라이드 02 · 대상과 과제}
연구 대상은 추상적 “경제”만이 아니라 생산·소비·가족·노동·이주를 선택한 사람들이다. 미시적 기록을 모아 장기 추세를 찾아도, 관찰된 반복을 곧 보편 법칙으로 선언해서는 안 된다. 가령 흑사병 뒤 임금 상승이라는 경향은 지역별 토지 소유와 노동 이동의 제약에 따라 달라졌다. 이 차이를 설명할 때 역사상이 풍부해진다.\par

\subsection{슬라이드 03 · 경제학에 주는 경험적 토대}
경제 이론은 관계를 명료하게 만들고, 경제사는 이론의 가정과 예측을 긴 시간과 여러 제도 아래에서 시험한다. 현대 데이터만 보면 드물거나 한 번만 일어난 충격을 관찰하기 어렵다. 역사 기록은 인구 붕괴, 제도 전환, 기술 도입처럼 \textbf{큰 변화가 실제로 어떻게 전개되었는지} 보여 준다. 자료가 있다는 것과 인과관계를 식별했다는 것은 별개다.\par

\subsection{슬라이드 04 · 슘페터의 문제의식}
슘페터가 경제사를 중시한 까닭은 현재의 사실 자체가 과거 선택의 누적 결과이기 때문이다. 같은 제도를 가진 듯한 두 나라가 다른 성장 경로를 보이면 초기 조건, 정치적 타협, 누적 투자까지 거슬러야 한다. 이론사·통계·역사를 분리된 암기 영역으로 보지 말고, \textbf{질문을 만들고 검증하는 서로 다른 도구}로 이해하자.\par

\subsection{슬라이드 05 · 사료와 이론의 왕복}
사료는 기록자가 특정 목적을 위해 만든 것이다. 세금 장부의 경작 면적은 과세를 피하려는 축소 신고, 군인 신장 자료는 선발 기준과 지원자의 구성에 영향을 받는다. 먼저 누가 왜 기록했는지 비판한 뒤 경제 이론으로 변수 간 기제를 세운다. 이론에 맞는 숫자만 고르는 대신, 이론이 예측한 반례도 함께 찾아야 한다.\par

\subsection{슬라이드 06 · 역사를 실험실로 쓸 때}
역사에서는 같은 사회를 두 번 살려 한 변수만 바꿀 수 없다. 하지만 전쟁·전염병·법 개정처럼 집단마다 노출이 달랐던 사건은 비교의 기회를 제공한다. 관찰한 차이를 원인으로 해석하려면 \textbf{충격 전 추세, 동시에 일어난 변화, 비교집단의 적절성}을 검토해야 한다. “자연 실험”은 통제 실험과 동의어가 아니다.\par

\subsection{슬라이드 07 · 두 가지 큰 서술 방식}
발전단계론은 제도와 생산 양식의 질적 전환을, 성장사학은 소득·생산성 같은 연속적인 양을 강조한다. 단계 구분은 전체 구조를 그리는 데 유용하지만 경계가 흐린 사례를 억지로 넣기 쉽다. 성장률은 비교 가능하지만 권리와 분배의 차이를 지우기 쉽다. 두 접근은 경쟁하는 유일한 정답보다 서로의 빈틈을 드러내는 질문으로 읽는 편이 좋다.\par

\subsection{슬라이드 08 · 독일 역사학파의 맥락}
19세기 독일은 영국과 산업화 단계가 달랐다. 보편적 자유무역 처방이 모든 국가에 같은 효과를 낸다는 주장에 역사학파는 의문을 제기했다. 관세동맹과 정치적 통일은 더 넓은 국내 시장을 형성했지만, 제도와 산업의 성장은 동일한 속도로 움직이지 않았다. 정책을 평가할 때 \textbf{현재의 비교우위와 미래의 학습 효과}를 구분해야 한다.\par

\subsection{슬라이드 09 · 리스트의 유치산업 논리}
리스트에게 보호의 목적은 영원한 독점 유지가 아니라 학습 중인 제조업이 규모와 숙련을 얻을 시간을 사는 것이었다. 오늘의 사적 비용이 세계가격보다 높아도 학습이 외부효과를 만들어 미래의 사회적 비용을 낮춘다면 일시적 보호가 논리적으로 가능하다. 그러나 학습 효과를 측정하고 종료 시점을 정하지 못하면 보호는 지대 추구로 굳어질 수 있다.\par

\subsection{슬라이드 10 · 단계론에 대한 검증 질문}
“목축 다음 농업, 다음 상공업”이라는 순서가 관찰되어도 그 순서를 만들어 내는 인과 기제가 증명된 것은 아니다. 해상 무역을 먼저 키운 도시국가나 농업과 공업이 병존한 지역은 단선적 서술의 반례가 된다. 역사적 정책 사례를 오늘의 처방에 쓸 때에는 초기 기술, 국가 역량, 국제질서가 얼마나 비슷한지 먼저 따져야 한다.\par

\subsection{슬라이드 11 · 생산력과 생산관계}
생산력은 노동·도구·기술이 만들 수 있는 산출의 가능성이고, 생산관계는 누가 토지와 도구를 통제하며 잉여를 가져가는가의 관계다. 두 요소가 충돌한다는 명제는 기술 변화만으로 사회가 바뀌지 않는다는 문제를 제기한다. 그러나 법·정치·종교가 항상 경제적 토대의 일방적 결과라고 가정하면 제도가 생산성을 다시 바꾸는 피드백을 놓친다.\par

\subsection{슬라이드 12 · 사회구성체의 순서}
원시 공산제·노예제·봉건제·자본제는 마르크스 전통의 분석 범주다. 각 범주는 \textbf{잉여가 누구에게 어떤 강제로 이전되는가}를 비교하는 데 도움이 된다. 실제 사회에는 자영농·노예·임금노동이 함께 존재했으며 지역마다 전환 경로도 달랐다. “아시아적 생산양식” 역시 하나의 균일한 아시아 현실로 받아들이기보다 논쟁적인 개념으로 다루어야 한다.\par

\subsection{슬라이드 13 · 마르크스 경제사학의 용도와 한계}
계급 간 분배와 강제의 문제를 전면에 놓는 장점이 있지만, 모든 변화가 정해진 단계로 수렴한다는 설명은 역사적 다양성과 맞지 않는다. 이후 장원 강의에서 부역·지대·토지 권리의 구분에는 생산관계의 언어가 유익하다. 그 설명이 충분한지는 임금, 기술, 정치 권력, 시장 연결의 독립적인 자료로 시험해야 한다.\par

\subsection{슬라이드 14 · 자본주의 맹아론을 읽는 법}
자본주의 “맹아”는 식민 지배가 없었다면 반드시 같은 목적지에 도달했으리라는 예언이 아니라, 외부 충격 이전의 상품 생산과 시장 확대를 찾으려는 문제의식이다. 상인·공방·임금노동의 존재만으로 자본주의 전체가 성립했다고 할 수 없다. 생산수단 소유, 노동계약, 시장을 위한 생산의 규모와 지속성을 함께 살펴야 한다.\par

\subsection{슬라이드 15 · 사회경제사학의 넓은 자료}
가격과 생산량만으로는 가족의 노동, 관습적 권리, 도시의 규칙을 이해하기 어렵다. 블로크와 사회경제사학의 장점은 법률 문서, 생활 관행, 공간 구조를 경제 과정에 결합한 데 있다. 자료가 다양해질수록 서로 충돌하는 증언을 다루는 원칙도 중요해진다. 각 자료가 무엇을 측정하고 무엇을 누락하는지 적어 두자.\par

\subsection{슬라이드 16 · 로스토우의 성장 단계}
“이륙”은 투자가 누적되고 특정 산업의 성장이 다른 산업으로 확산된다는 직관을 준다. 그러나 성장률이 한 번 높아졌다고 모든 지역이 같은 단계에 들어섰다고 볼 수는 없다. 이륙 시점의 판정에는 투자율, 생산성, 산업 구조와 제도 변화의 자료가 필요하며, 후발국의 수입 기술 활용은 선발국과 다른 경로를 만들 수 있다.\par

\subsection{슬라이드 17 · 성장률과 계열 자료}
산출 $Y_t$의 연간 성장률은 \textbf{전년 수준}을 분모로 잡는다. $g_t=(Y_t-Y_{t-1})/Y_{t-1}$이며 백분율은 여기에 100을 곱한다. 서로 다른 시대의 화폐액을 비교할 때에는 물가를 제거한 실질 시계열이 필요하다. 1인당 산출의 로그성장률은 근사적으로 총산출 성장률에서 인구 성장률을 뺀 값이다: $\Delta\ln(Y/N)=\Delta\ln Y-\Delta\ln N$.\par

\subsection{슬라이드 18 · 맬서스의 핵심 제약}
토지가 고정되고 노동이 늘 때 노동의 한계생산이 낮아지는 경제를 생각하자. $Y=A T^{\beta}N^{1-\beta}$에서 $T$는 토지, $N$은 인구, $0<\beta<1$이다. 그러면 1인당 산출은 $y=Y/N=A(T/N)^{\beta}$이므로 기술 $A$와 토지가 그대로일 때 인구가 늘면 $y$는 떨어진다. 산출 총량 증가와 생활 수준 개선이 다른 이유다.\par

\subsection{슬라이드 19 · 맬서스 함정의 동학}
출생률 $b(y)$가 소득과 함께 올라가거나 사망률 $d(y)$가 내려간다고 하자. 인구 증가율은 $\dot N/N=b(y)-d(y)$이고, 장기 균형에서는 두 율이 같다. 기술이 한 번 좋아져 $A$가 증가하면 처음에는 $y$가 상승하지만, 인구가 늘어 토지/사람 비율이 다시 낮아진다. 이 과정은 기술 진보가 \textbf{총인구에는 남고 1인당 소득에는 거의 남지 않을 수 있음}을 설명한다. 예방적 출산 억제나 지속적인 생산성 성장은 이 결론을 바꾼다.\par

\subsection{슬라이드 20 · 근대 성장의 다른 체제}
근대적 경제성장에서는 인구가 늘어도 1인당 산출이 장기간 증가한다. 항등식 $Y=yN$에 로그를 취하면 $g_Y\approx g_y+g_N$이다. 따라서 총산출 성장률이 인구 성장률을 지속적으로 웃돌아야 생활 수준이 오른다. 산업혁명을 한 해의 발명이나 단번의 도약으로 축소하지 말고, 생산성·에너지·인적자본·제도의 누적 변화를 함께 살펴야 한다.\par

\subsection{슬라이드 21 · 맬서스 논리의 적용 범위}
전쟁이나 전염병이 있다고 해서 어떤 지역을 오늘 곧바로 “맬서스 함정”이라고 부를 수는 없다. 토지 제약, 기술 성장률, 출산 선택과 사망률의 소득 반응이 모두 검증되어야 한다. 반대로 산업사회에서도 환경과 자원 제약이 사라진 것은 아니다. 핵심은 기술과 제도가 제약을 얼마나 빠르게 완화하는가다.\par

\subsection{슬라이드 22 · 클리오메트릭스의 질문}
Cliometrics는 경제 이론으로 반박 가능한 역사적 가설을 만들고, 수량 자료로 그 크기를 따진다. 예를 들어 “철도가 미국 성장에 필수였다”는 주장은 철도가 없었다면 같은 화물을 옮기는 비용이 얼마나 더 들었는지를 계산해야 평가할 수 있다. 숫자를 쓰는 것 자체보다 \textbf{비교할 반사실과 필요한 가정}을 명시하는 일이 중요하다.\par

\subsection{슬라이드 23 · 반사실적 실험의 구조}
관찰한 결과 $Y(1)$과 관찰할 수 없는 대안 $Y(0)$의 차이가 인과효과다. 한 사회에 대해 두 결과를 동시에 볼 수 없으므로 모형이나 비교집단이 $Y(0)$을 대신한다. 철도 사회적 절감 추정이라면 대체 수로·도로의 용량과 요금, 운송 경로가 반사실의 핵심 가정이다. 가정이 바뀌면 추정치도 바뀌므로 한 숫자만 외우지 말자.\par

\subsection{슬라이드 24 · 노스와 제도의 역할}
제도는 “게임의 규칙”, 조직은 그 규칙 아래에서 행동하는 기업·길드·국가 기구다. 소유권이 명확해져도 법원이 계약을 집행하지 못하면 거래비용이 높다. 장기 성장에서는 자본 축적의 양뿐 아니라 투자 수익을 회수할 수 있다는 믿음, 분쟁 해결, 정치 권력의 제약이 중요하다. 제도가 성장의 원인인지 성장의 결과인지 구분하는 실증 문제도 남는다.\par

\subsection{슬라이드 25 · 맥클로스키의 질문}
경제사가 경제학에 쓸모 있으려면 “과거에도 흥미로운 일이 있었다”를 넘어, 사실·이론·정책 판단을 바꾸는 구체적인 경로를 보여야 한다. 맥클로스키의 질문은 경제사가 단순한 사례 창고인지, 경제학의 가정을 시험하고 경제학자를 훈련하는 독립적 지식인지 묻는다. 다섯 가지 답은 뒤 슬라이드에서 각각 다른 용도를 설명한다.\par

\subsection{슬라이드 26 · 탈역사화가 낳는 표본 문제}
최근 수십 년의 자료만으로 보편적 경제법칙을 추정하면 표본 밖의 제도와 충격에 약하다. 예컨대 금융위기를 안정기 데이터로만 공부하면 드문 꼬리 사건을 보지 못한다. 과거 자료도 결측과 측정 오류가 있지만, 더 긴 기간과 다양한 제도적 환경을 제공한다. 데이터가 많다는 이유만으로 과거를 무시하거나, 오래되었다는 이유만으로 과거를 낭만화하지 말아야 한다.\par

\subsection{슬라이드 27 · 다섯 가지 효용의 구별}
더 많은 사실은 관찰 범위의 확대, 더 나은 사실은 자료의 깊이와 해석의 개선이다. 더 나은 이론은 기존 가정의 반례를 찾는 것이고, 더 나은 정책은 과거의 유사한 정책이 어떤 제약 아래 작동했는지 배우는 것이다. 마지막 “더 나은 경제학자”는 방정식과 제도적 현실을 왕복할 줄 아는 연구자의 훈련이다. 서로 다른 주장에 서로 다른 증거가 필요하다.\par

\subsection{슬라이드 28 · 과거의 미시 자료}
장원 회계 장부에는 지대·노동일수·작황이, 인구조사 원장에는 직업·가구·이동이 기록될 수 있다. 오늘의 국민계정 이전에도 자료가 없던 것은 아니다. 다만 십일조 기록은 납부 대상과 회피, 센서스는 조사 범위와 직업 분류가 시대마다 달라 그대로 이어 붙일 수 없다. \textbf{오래된 기록을 통계로 바꾸는 과정 자체가 연구}다.\par

\subsection{슬라이드 29 · 기록의 질과 선택 편향}
시간이 지나 공개된 기업 장부나 개인 기록은 현대의 익명 집계보다 세부 정보를 더 줄 수 있다. 그러나 살아남아 보존된 기록은 큰 기업·부유한 가구·관료 조직에 치우칠 수 있다. 누락이 결과와 관련된다면 평균을 그대로 일반화할 수 없다. 표본이 어떻게 만들어졌는지를 숫자의 소수점 자리보다 먼저 확인해야 한다.\par

\subsection{슬라이드 30 · 큰 사건을 비교할 때}
전염병과 대공황은 노동·정책의 효과를 볼 기회를 주지만, 여러 변수가 동시에 변한다. 예컨대 전염병 뒤 임금이 오른다면 노동 감소뿐 아니라 수요 감소와 물가 변화도 실질임금에 영향을 준다. 분석 단위가 도시인지 개인인지에 따라 결론도 달라질 수 있다. 역사적 실험실에서는 \textbf{비교의 설계와 메커니즘의 설명}을 함께 요구한다.\par

\subsection{슬라이드 31 · 정형화된 사실을 시험하기}
성장 모형이 자본/산출 비율이나 저축률을 일정하다고 가정하는 것은 계산에 도움이 되지만, 어느 시대에나 같은 사실이라는 뜻은 아니다. 역사 시계열은 그 비율의 변화와 제도 전환을 보여 준다. 모형이 실패한 곳에서 어떤 마찰—토지 권리, 신용 제약, 가족 노동—을 추가해야 하는지 발견할 수 있다.\par

\subsection{슬라이드 32 · 정책에 필요한 역사적 비교}
정책 논의에서 철도·보호관세·화폐제도를 과거의 “성공 사례” 하나로 정당화하기 쉽다. 역사 비교가 유용하려면 원래의 목표, 당시의 대안, 수혜자와 비용 부담자, 경제의 규모와 국제환경을 함께 비교해야 한다. 반사실적 질문은 \textbf{그 정책이 없었다면 무엇이 일어났을까}이지, 그 정책이 있었던 사회가 성장했다는 관찰만이 아니다.\par

\subsection{슬라이드 33 · 두 다리로 걷는 경제학}
연역은 가정에서 예측을 만들고, 귀납은 관찰에서 패턴과 예외를 찾는다. 둘 중 하나만으로는 “왜”라는 질문을 끝낼 수 없다. 과거의 현상을 오늘의 용어로 표현할 때에는 편의가 있지만, 당시 행위자가 실제로 기대했던 선택지를 복원해야 한다. 이론의 절약성과 역사의 복잡성 사이를 왕복하는 능력이 핵심이다.\par

\subsection{슬라이드 34 · 데이비드의 QWERTY 논문}
표준의 채택은 가격이 즉시 최적기술을 고르는 단순 경쟁과 다르다. 사용자가 이미 배운 배열, 기업의 채용 기준, 기계의 호환성이 서로 영향을 주기 때문이다. 데이비드의 사례는 \textbf{현재의 균형을 설명할 때 과거의 채택 순서가 중요한가}를 묻는다. QWERTY의 실제 효율성 우열은 별도의 실험적 증거로 검증해야 한다.\par

\subsection{슬라이드 35 · 효율성 주장과 경로의존성}
QWERTY가 보편화되었다는 사실은 그것이 가장 빠르다는 증거가 아니다. 반대로 대안 배열이 일부 실험에서 빠르다는 주장만으로 사회 전체의 전환이 이익이라는 결론도 나오지 않는다. 재훈련 비용과 호환성, 실험의 표본·설계를 모두 고려해야 한다. \textbf{경로의존성은 채택 역사에 관한 설명이고, 비효율성은 추가 입증이 필요한 평가}다.\par

\subsection{슬라이드 36 · 초기 우연이 남는 조건}
초기 타자기의 기계 구조와 판매 전략이 표준 형성에 영향을 주었다. 그러나 “활자 엉킴을 피하려고 의도적으로 느린 배열을 만들었다”는 통속적 설명은 논문 자체의 복잡한 역사와 후속 연구를 지나치게 단순화한다. 중요한 경제학적 질문은 첫 우위가 생긴 뒤 \textbf{재학습과 호환성의 누적 효과}가 우위를 얼마나 확대했는가다.\par

\subsection{슬라이드 37 · 세 가지 자기강화 장치}
하드웨어와 숙련의 보완성, 같은 표준 사용자가 많을수록 커지는 네트워크 이익, 과거 훈련을 되돌리기 어려운 비용이 결합한다. 한 요소만으로 고착이 반드시 생기는 것은 아니다. 새 배열의 생산성 이익이 충분히 크거나 전환 비용을 조정할 수 있다면 다른 균형으로 옮겨갈 수도 있다.\par

\subsection{슬라이드 38 · 보완재의 경제학}
키보드와 숙련은 함께 있을 때 생산적이다. 새 기계만 공급해도 숙련된 타자수가 없으면 기업이 채택하지 않고, 교육만 바꿔도 새 기계가 없으면 훈련자가 보상을 받지 못한다. 이는 두 시장이 서로를 기다리는 조정 문제다. 오늘의 소프트웨어 형식과 사용자 훈련, 전기차와 충전망에서도 유사한 구조를 찾을 수 있다.\par

\subsection{슬라이드 39 · 네트워크 외부성과 임계점}
표준 $A$의 사용자가 늘면 추가 사용자가 얻는 호환성 이익이 커질 수 있다. 개인은 자신의 선택이 다른 사용자에게 주는 이익을 모두 계산하지 않으므로 사적 선택과 사회적 최적이 어긋날 수 있다. 초기에 조금 앞선 표준이 임계점을 넘으면 채용·교육·제조가 같은 방향으로 이동한다. 다만 네트워크가 커졌다는 사실만으로 해당 표준의 기술적 열등성이 증명되는 것은 아니다.\par

\subsection{슬라이드 40 · 되돌릴 수 없는 인적자본}
기계의 물리적 교체 비용이 내려가도 숙련된 사람의 재훈련 시간은 남는다. 기존 투자비는 이미 지출한 sunk cost이지만, \textbf{앞으로 전환에 필요한 시간과 생산 손실은 의사결정에 관련된 비용}이다. 사적 전환 비용이 클 때 조정된 집단 전환의 잠재 이익이 있어도 아무도 먼저 움직이려 하지 않을 수 있다.\par

\subsection{슬라이드 41 · 경로의존성을 수식으로 보기}
두 표준 $A,B$의 순편익을 $u_A(n_A)=v_A+\lambda n_A-c_A$, $u_B(n_B)=v_B+\lambda n_B-c_B$로 단순화하자. $n_A$가 이미 크면, $v_B>v_A$라도 $B$의 기술적 이점이 네트워크 차이와 전환 비용을 넘지 못해 $A$를 택한다. 이는 \textbf{초기 채택 순서가 후속 선택의 보수를 바꾼다}는 뜻이다. 효율성 판단에는 모든 사용자의 전환 비용과 장기 편익을 함께 합산해야 한다.\par

\subsection{슬라이드 42 · 골딘이 정리한 혁명}
1993년 포겔과 노스의 수상은 경제사의 정량적·이론적 연구가 경제학 전체의 질문을 바꿨다는 평가와 연결된다. 클리오메트릭스는 단순히 회귀분석을 옛 자료에 적용하는 일이 아니다. \textbf{반사실을 세우고 대안을 수량화하며 자료의 편향을 논쟁 가능한 형태로 드러내는 방법}이다.\par

\subsection{슬라이드 43 · 클리오메트릭스의 정의}
Clio는 역사의 뮤즈이고 metrics는 측정이다. 경제 이론, 수량화, 역사 자료의 세 축이 만나야 한다. 수량 추정이 정교해도 당시 제도를 잘못 정의하면 추정 대상이 흔들리고, 훌륭한 서사가 있어도 반박 가능한 예측이 없으면 경제적 크기를 비교하기 어렵다. 골딘의 논문은 포겔과 노스의 서로 다른 연구가 이 공통 방식을 어떻게 확장했는지 보여 준다.\par

\subsection{슬라이드 44 · 철도의 사회적 절감}
포겔의 질문은 철도가 없으면 모든 운송이 사라지는가가 아니라, 수로·도로로 대체했을 때 비용이 얼마나 늘어나는가였다. 단순화하면 사회적 절감은 동일한 화물량 $Q$에 대한 $(c_{\mbox{대안}}-c_{\mbox{철도}})Q$다. 실제 추정에서는 경로, 계절, 대체망 건설비와 수요의 가격 반응이 필요하다. “필수적이지 않다”는 결론은 철도의 이익이 0이라는 뜻이 아니다.\par

\subsection{슬라이드 45 · 노예제 연구와 평가 기준}
플랜테이션의 산출 효율을 계산하는 것과 노예제가 정당하다는 판단은 완전히 다른 문제다. 강제·폭력·자유의 상실은 회계상 생산성으로 환원되지 않는다. 포겔·엥거먼의 추정은 자료와 방법에 큰 논쟁을 낳았으며, 생산성 비교에는 토지 품질, 작물, 노동시간, 폭력의 비용과 생존자 편향을 살펴야 한다. \textbf{효율성과 인권을 분리하여 평가}하자.\par

\subsection{슬라이드 46 · 노스가 묻는 거래비용}
동일한 기술을 가진 두 사회도 계약을 집행하고 권력을 제한하는 방식이 다르면 투자 유인이 달라진다. 거래비용은 거래자를 찾고, 품질을 확인하고, 약속을 강제하는 데 드는 비용이다. 소유권을 법에 적었다는 사실만으로 실제 집행이 보장되지는 않는다. 제도의 변화는 권력자의 인센티브와 비공식 규범까지 포함한다.\par

\subsection{슬라이드 47 · 신장과 영양의 역사}
키는 어린 시기의 영양과 질병 환경을 반영할 수 있지만, 소득의 단순한 대체 지표가 아니다. 징집 기준이 바뀌거나 자원입대자의 구성이 경기와 함께 변하면 표본 평균도 바뀐다. 포겔의 신체계측 연구는 생활 수준을 화폐 소득 밖으로 넓혔고, 동시에 \textbf{생물학적 지표의 선택 편향}을 엄격히 다루어야 한다는 숙제를 남겼다.\par

\subsection{슬라이드 48 · 경제사의 넓은 진료}
장기 성장에는 노동, 금융, 건강, 가족과 정치가 서로 영향을 준다. 한 분야의 짧은 패널만 보면 이 상호작용을 놓칠 수 있다. 골딘이 강조하는 경제사학자의 역할은 모든 분야의 전문가를 대체하는 것이 아니라, 부분 분석의 결과를 장기 구조와 제도적 맥락 안에서 다시 점검하는 것이다.\par

\subsection{슬라이드 49 · 그라이프의 재평가}
그라이프는 초기 클리오메트릭스의 성과를 인정하면서도, 당시 주류 이론의 가정이 연구 질문을 좁혔다고 지적한다. 동일한 선호와 기술을 주면 한 균형으로 돌아간다는 모형은 역사의 경로와 집단 간 전략적 상호작용을 설명하기 어렵다. 이론의 변화는 경제사에 더 많은 질문을 허용한다.\par

\subsection{슬라이드 50 · 성과와 제약을 동시에 보기}
자료의 수량화는 철도, 노동, 제도에 대한 직관을 검증 가능하게 만들었다. 그러나 컴퓨터 성능의 부족만이 문제가 아니었다. 거래비용이 낮고 완전정보에 가까운 모형을 기본으로 삼으면, 왜 길드나 상인 네트워크가 생기고 유지되는지 질문하기 어렵다. 분석 도구의 선택이 곧 보이는 역사 현상의 선택을 바꾼다.\par

\subsection{슬라이드 51 · 단일 균형 가정의 한계}
생산집합의 볼록성과 완전한 계약을 가정하면 초기 위치가 달라도 같은 가격·배분으로 수렴하기 쉽다. 그러나 구성원의 기대와 상호 신뢰가 서로를 뒷받침하면 여러 균형이 가능하다. 예를 들어 모두가 계약을 지킨다고 믿는 시장과 아무도 상대를 믿지 않는 시장은 같은 기술로도 다른 거래량을 유지할 수 있다.\par

\subsection{슬라이드 52 · 시장 밖 조직도 설명 대상}
가족, 길드, 국가의 규칙은 “시장 실패를 방해하는 잔재”로만 볼 수 없다. 정보가 부족하고 법적 집행이 약할 때에는 평판과 공동체의 제재가 거래를 가능하게 만들 수 있다. 동시에 진입 장벽과 배제도 만든다. 같은 제도의 순효과는 누구의 거래를 열고 누구를 막는지, 대안 제도가 무엇인지에 달려 있다.\par

\subsection{슬라이드 53 · 새 미시이론과 역사}
게임이론, 정보경제학, 계약이론은 과거의 규범과 조직을 경제학의 언어로 분석할 길을 넓혔다. 일회 거래에서는 속이는 편이 유리해도 반복 거래와 평판 손실이 충분하면 협력이 자기강제적 균형이 될 수 있다. 이 조건을 구체적인 상인 집단과 정치 질서에 적용할 때 사료의 세부가 필수다.\par

\subsection{슬라이드 54 · 제도는 왜 스스로 유지되는가}
어떤 규칙이 오래 지속되었다는 사실만으로 효율적이었다고 할 수 없다. 기존 규칙의 수혜자가 변화를 막을 수 있고, 다른 사람은 상대도 함께 바뀔 것이라는 믿음이 없어 움직이지 못할 수 있다. 제도 분석에서는 개인의 보수, 제재를 집행할 조직, 신념의 변화와 외부 충격을 함께 살펴야 한다.\par

\subsection{슬라이드 55 · 성장 모형과 이질성}
총생산 함수 하나로 평균적인 경제를 나타내면 지역·집단 간 권리와 정치력이 사라질 수 있다. 기술 확산은 교육·금융·특허·토지 제도의 차이에 따라 속도가 다르다. 내생적 성장의 관점은 지식 축적을 설명하는 데 유용하지만, 역사에서 어떤 집단이 교육과 투자 기회에서 배제되었는지 따져야 한다.\par

\subsection{슬라이드 56 · 귀납과 연역의 재결합}
구체적인 역사 사례에서 반복되는 문제를 찾고 이를 간결한 모형으로 표현한 뒤, 다른 사례에서 예측이 맞는지 확인한다. 새 이론을 적용해 모든 결과를 사후에 설명할 수 있다면 반증이 불가능해진다. 변수와 시점을 미리 정하고, 같은 모형이 설명하지 못하는 곳까지 공개해야 한다.\par

\subsection{슬라이드 57 · 모형의 다양성과 판별}
한 제도의 형성은 정보 부족, 권력 유지, 네트워크 효과 등 여러 모형으로 설명될 수 있다. 어느 설명이 맞는지 판별하려면 각 모형만의 예측을 찾아야 한다. 예컨대 정보 문제 설명은 거래 상대의 평판 자료와 제재 방식을, 정치경제 설명은 제도 변경의 수혜자와 거부권을 더 직접적으로 예측한다. \textbf{설명 가능성보다 구별 가능성}이 중요하다.\par

\subsection{슬라이드 58 · 유용한 경제사의 기준}
방법이 정교해도 다른 경제학자와 역사학자, 정책 담당자가 이해할 질문을 만들지 못하면 연구의 영향은 좁다. 경제사의 유용성은 현대 문제와 같은 답을 낸다는 뜻이 아니다. 과거의 다른 조건을 이용해 오늘 당연하게 받아들이는 제도와 가정을 낯설게 만드는 데 있다.\par

\subsection{슬라이드 59 · 학문 간 단절의 비용}
계량 방법을 공유하는 것만으로 과거의 의미가 전달되지는 않는다. 경제학자는 추정치를 원하고 역사학자는 행위자의 언어와 맥락을 중시할 수 있다. 좋은 연구는 둘을 연결해 추정 대상이 과거 사람들이 실제로 경험한 어떤 변화인지 설명한다. “유용하다”는 말도 청중에 따라 학술 이론, 정책, 공적 역사 이해라는 다른 기준을 갖는다.\par

\subsection{슬라이드 60 · 신장 자료와 제도 설명의 자기 점검}
신장 자료는 영양과 질병을 포착하지만 군대·감옥·학교 표본은 인구 전체를 대표하지 않을 수 있다. 제도 설명도 “좋은 제도가 성장했다”는 상관만 제시하면 순환 논리에 빠진다. 자료와 이론이 새 시각을 주는 만큼 \textbf{표본 선택과 인과 방향}을 함께 공개해야 한다.\par

\subsection{슬라이드 61 · 세 청중에 대한 세 질문}
첫째, 경제학에는 새 가설이나 더 나은 사실을 주는가. 둘째, 정책에는 비슷한 선택의 장기 비용을 가늠할 근거를 주는가. 셋째, 대중에게는 단순한 영웅·몰락 서사 대신 복잡하지만 이해 가능한 과거의 그림을 주는가. 한 연구가 세 질문 모두에 같은 방식으로 답할 필요는 없지만, 누구를 설득하는지 명확해야 한다.\par

\subsection{슬라이드 62 · 공간·시간·인간을 되살리기}
공간은 단순 운송비 변수가 아니라 강, 항구, 도시의 밀도와 통치 범위다. 시간은 할인율의 지수만이 아니라 숙련·관습·인프라가 누적되는 과정이다. 인간은 평균 노동자 하나로 합쳐지지 않는 성별·지위·가족의 차이를 가진다. 이런 차원을 되살리면 추상 모형의 변수에 역사적 내용이 생긴다.\par

\subsection{슬라이드 63 · 노동 조건의 질}
실질임금이 같아도 사고 위험, 강제, 노동시간, 직장 내 존중은 다를 수 있다. 가사 노동과 임금 노동의 선택을 볼 때에도 시장 임금만이 아니라 집 안팎의 협상력과 안전을 고려해야 한다. 생활 수준을 소득 하나로 환원하면 제도 변화의 중요한 이익과 비용을 놓친다.\par

\subsection{슬라이드 64 · 앙시앵 레짐 재정의 유사성}
재정 적자, 전쟁비용과 조세 특권이라는 조합은 시대를 넘나드는 질문을 만든다. 하지만 18세기 프랑스의 조세 징수권·대표제·국채 시장을 오늘의 정부와 동일하게 취급해서는 안 된다. 유사성은 \textbf{어떤 정치적 이해관계가 세입 개혁을 막는가}라는 가설을 세우는 데 쓰고, 차이는 그 가설의 적용 범위를 정한다.\par

\subsection{슬라이드 65 · 공공투자의 시간 지평}
오래 쓰는 도로·상하수도·학교는 현재 비용과 여러 세대의 편익을 나눈다. 정치인은 짧은 선거 주기에서 유지보수보다 눈에 띄는 새 사업이나 즉시 지출을 선호할 수 있다. 역사적 도시 인프라 사례는 장기 투자의 필요를 떠올리게 하지만, 당시 건설이 누가 낸 세금과 노동으로 가능했는지도 함께 물어야 한다.\par

\subsection{슬라이드 66 · 만들어 낸 과거를 경계하기}
현대 논쟁에 편리한 “옛날에는 모두…” 식의 이야기는 대개 지역·계층의 차이를 지운다. 자료로 확인할 수 있는 명제와 도덕적 평가, 오늘에 대한 비유를 분리하자. 경제사는 현재의 정책을 정당화하는 전설을 만드는 대신, 과거의 실제 선택지와 예상치 못한 결과를 밝혀야 한다.\par

\subsection{슬라이드 67 · 네 가지 큰 논쟁}
Dobb–Sweezy는 봉건제 내부의 계급관계와 외부 상업 중 무엇이 이행을 이끌었는지, Brenner는 토지 소유와 계급관계의 지역 차이를 물었다. 산업혁명 논쟁은 기술·임금·에너지·식민지의 상대적 역할을, Great Divergence 논쟁은 유럽과 아시아의 장기 분기를 다룬다. 답을 외우기보다 \textbf{각 주장에 필요한 관찰}을 비교해 보자.\par

\subsection{슬라이드 68 · 역사비교제도분석}
중앙 권력이 계약을 모두 강제하지 못할 때 상인들은 반복 거래, 평판, 공동체의 처벌로 신뢰를 만들 수 있다. 어떤 제도가 자기강제적이라는 것은 규칙을 지키는 것이 각 참여자에게 현재와 미래를 합쳐 유리하다는 뜻이다. Nash 균형이라는 이름만으로 제도의 정의로움을 보장하지는 않는다. 배제된 사람의 비용과 외부의 대안도 살펴야 한다.\par

\subsection{슬라이드 69 · 근대 시장 모형의 조건}
익명의 거래가 작동하려면 물건의 품질을 확인하고 계약 불이행을 제재할 장치가 필요하다. 사적 평판망이 그 일을 할 수도 있고 법원·표준화·금융기관이 넓은 시장에서 맡을 수도 있다. 그림의 시장 참여자와 정보 흐름을 따라가며, 어떤 관계가 깨졌을 때 거래량이 줄어드는지 확인하자. 가격만으로 시장의 존재를 설명할 수 없다.\par

\subsection{슬라이드 70 · 아날학파와 일상생활}
브로델의 장기지속은 정치적 사건보다 느리게 변하는 지리, 인구, 물질생활을 전면에 둔다. 식사·주거·가족의 관습은 소비의 제약이자 노동의 조직 방식이다. 국민소득 통계가 없는 시대에도 음식·주거·물품 목록을 통해 생활 수준을 복원할 수 있지만, 소수의 기록을 전체 사회로 일반화하는 데에는 주의가 필요하다.\par

\subsection{슬라이드 71 · 확장되는 연구 의제}
금융위기, 환경, 신체계측, 역사적 국민계정은 서로 다른 자료와 추정 방법을 요구한다. 공통 질문은 \textbf{변화의 크기와 분배, 제도적 매개}다. 과거의 화폐액을 오늘과 비교할 때는 물가뿐 아니라 상품의 질과 소비 바구니의 변화도 고려해야 한다. 새 데이터베이스가 생겨도 해석의 문제가 사라지지는 않는다.\par

\subsection{슬라이드 72 · 부의 격차를 설명하는 질문}
나라 간 소득 격차는 기술 차이 한 줄로 끝나지 않는다. 인구 압력, 교육·건강, 토지 권리, 무역, 국가 역량, 식민 지배와 에너지 접근이 서로 영향을 준다. 산업혁명이 왜 영국에서, 왜 그 시점에 일어났는지 묻는 것은 각 설명이 \textbf{다른 나라와 더 이른 시대를 어떻게 예측하는가}를 함께 묻는 일이다.\par

\subsection{슬라이드 73 · 소득과 행복의 거리}
1인당 GDP는 평균적인 시장 생산의 척도이지 건강·자유·환경·관계의 완전한 지표가 아니다. 소득이 늘어도 기대와 욕망이 함께 변하면 주관적 행복의 상승 폭이 작을 수 있다는 것이 Easterlin paradox의 문제의식이다. 단면 국가 비교와 한 나라의 장기 추세를 구별하고, 소득·수명·교육·분배를 함께 보자.\par

\subsection{슬라이드 74 · 긴 시간의 그림을 읽는 법}
1800년 전후의 소득 도약을 보여 주는 그림은 강한 요약이지만, 모든 지역이 한 날에 맬서스 체제에서 벗어난 것은 아니다. 로그축인지 선형축인지, 실질 구매력을 어떻게 환산했는지, 국가 평균 뒤의 분배 차이가 무엇인지 확인해야 한다. 다음 인구 강의에서는 이 그림의 분모와 맬서스 기제를 더 자세히 푼다.\par

\subsection{읽기 · Does the Past Have Useful Economics?}
\subsubsection{논문의 질문}

맥클로스키의 질문은 과거에 흥미로운 이야기가 많은가가 아니다. \textbf{경제학자가 과거를 공부했을 때 현재의 사실, 이론, 정책 판단이 어떻게 달라지는가}다. 20세기 중반 경제학의 수학화와 거시·계량 방법의 발전은 분석을 정교하게 만들었지만, 역사적 사례와 자료를 주변으로 밀어낼 위험도 있었다. 저자는 경제사가 경제학에 주는 효용을 다섯 갈래로 정리한다. 이 갈래는 단순한 찬사가 아니라 각각 다른 연구 설계를 요구한다. \href{https://ideas.repec.org/a/aea/jeclit/v14y1976i2p434-61.html}{서지 정보}.\par

\subsubsection{1. 더 많은 경제적 사실}

현대 통계가 시작되기 전에도 장원 회계, 십일조, 가격 장부, 항해 기록, 임금 계약과 인구조사가 남아 있다. 자료가 처음부터 GDP 형태로 쓰인 것은 아니므로, 기록의 단위와 누락을 해석해야 한다. 예컨대 중세 장원의 장부에서 노동 의무가 며칠로 적혔다면 이는 노동 공급의 기록이면서 동시에 영주의 청구권 기록이다. 두 가구의 “노동일수”를 합쳐 실제 작업시간으로 쓰려면 의무일과 수행일을 구별해야 한다.\par

경제학적으로 사실을 늘린다는 것은 표본의 길이와 제도적 범위를 넓힌다는 뜻이다. 최근 30년만 본 성장 모형은 전쟁과 인구 붕괴, 산업 전환처럼 드문 사건에 약하다. 과거를 추가하면 이론이 관찰하지 못하던 경우를 시험할 수 있다. 그러나 옛 기록의 숫자가 많다고 신뢰도가 자동으로 높아지지는 않는다.\par

\subsubsection{2. 더 나은 경제적 사실}

오랜 시간이 지나 공개된 원장에는 개인 이름, 가족 관계와 사업의 실제 거래가 남아 오늘의 익명 집계보다 풍부할 수 있다. 하지만 \textbf{기록자가 왜 기록했는가}를 먼저 물어야 한다. 세금을 피하려는 축소 신고, 군대의 신장 기준, 큰 기업만 남은 장부는 표본 편향을 만든다. 경제사 연구자는 현대의 잘 정제된 데이터셋을 소비하는 데서 멈추지 않고 데이터의 생성 과정까지 재구성한다.\par

그 결과 “영국 노동자의 임금이 올랐다”는 문장도 세부 질문으로 나뉜다. 명목임금인가 실질임금인가, 성인 남성인가 전체 노동자인가, 정해진 임금인가 실제 지급액인가, 식품 가격을 어떤 바구니로 나눴는가. 이런 질문은 현대 통계에도 그대로 유효하다.\par

\subsubsection{3. 더 나은 이론}

이론은 단순한 가정을 통해 변수 사이의 관계를 분명히 하지만, 관찰되지 않은 시대까지 자동으로 적용되지는 않는다. 예를 들어 성장 모형이 일정한 자본/산출 비율을 전제로 한다면 장기 역사 자료는 그 비율이 제도와 기술 변화 중에도 안정적인지 검사한다. 예상과 다르면 이론을 버리거나 새로운 마찰을 넣어야 한다. 역사 속의 길드와 장원은 \textbf{거래비용, 집행, 위험 분담}을 설명하는 이론의 자극이 될 수 있다.\par

이론의 실패를 찾는 데에도 방법이 필요하다. 특정 시대가 모형과 다르다는 사실만으로 “경제학은 쓸모없다”고 말할 수 없다. 모형의 어떤 가정—자유로운 이동, 완전한 계약, 동일한 기술—이 달랐는지 식별해야 이론이 개선된다.\par

\subsubsection{4. 더 나은 정책}

정책 논의에서 과거는 경고와 가능성을 동시에 준다. 철도가 없었다면 미국의 성장이 불가능했을 것이라는 통념을 검증하려면 수로와 도로라는 실제 대안을 계산해야 한다. 과거의 특정 제도가 성공한 사회에서 오늘도 성공하리라고 가정해서는 안 된다. 인구 구조, 행정능력, 국제시장과 그 제도에 대한 기대가 달라졌을 수 있다.\par

정책 비교의 기본 순서는 다음과 같다.\par

\noindent\textbullet 당시 정책의 목표와 실제로 가능한 대안은 무엇이었는가?\par
\noindent\textbullet 누가 비용을 내고 누가 이익을 얻었는가?\par
\noindent\textbullet 다른 변화가 동시에 일어났는가?\par
\noindent\textbullet 오늘의 상황과 공통인 기제와 다른 제약은 무엇인가?\par

이 순서를 거치면 단순한 역사적 비유가 인과 질문으로 바뀐다.\par

\subsubsection{5. 더 나은 경제학자}

자료가 왜 그런 형태로 남았는지, 오늘 당연한 제도가 과거에는 왜 없었는지, 평균 뒤에 어떤 사람들이 있었는지 묻는 습관이 경제학자의 판단을 넓힌다. 연역적 이론과 귀납적 역사 사이를 왕복해야 한다는 주장은 수식을 피하라는 뜻이 아니다. 오히려 수식이 \textbf{무엇을 생략했는가}를 볼 수 있어야 제대로 사용한다는 뜻이다.\par

\subsubsection{적용: 흑사병 뒤 실질임금}

노동 $L$, 고정 토지 $T$인 생산함수 $Y=A T^\alpha L^{1-\alpha}$를 두면 경쟁적 실질임금은 다음과 같다.\par

\[
\frac{w}{P}=\frac{\partial Y}{\partial L}=(1-\alpha)A\left(\frac{T}{L}\right)^\alpha.
\]

인구 충격으로 $L$이 줄면 이론은 노동의 한계생산과 실질임금의 상승을 예측한다. 이제 역사 자료가 필요하다. 실제 노동을 떠날 수 있었는가, 임금 상한이 시행되었는가, 식량가격 $P$가 어떻게 움직였는가, 영주가 현물로 지급했는가에 따라 관찰된 결과가 다르다. \textbf{이론이 질문을 선명하게 하고 사료가 이론의 조건을 검사}하는 왕복이 논문의 실천적 요지다.\par

\subsubsection{읽은 뒤 확인할 질문}

1. “더 많은 사실”과 “더 나은 사실”은 왜 다른가?\par
2. 반례가 하나 나오면 어떤 가정을 먼저 검토해야 하는가?\par
3. 현대 정책에 과거 사례를 쓸 때 단순한 유사성보다 중요한 것은 무엇인가?\par


\subsection{읽기 · Clio and the Economics of QWERTY}
\subsubsection{왜 키보드 배열을 경제사로 연구하는가}

데이비드의 관심은 타자기 발명 연대 자체보다 \textbf{현재의 선택이 과거의 선택에 의존하는 경제}다. 기술적 성능만으로 한 표준의 시장 점유율을 예측할 수 있다면 역사는 초기 조건에 불과하다. 그러나 하드웨어, 숙련과 사용자 수가 서로 보완하면 과거의 작은 우위가 후속 선택의 보수를 바꾼다. QWERTY는 이론적 메커니즘을 설명하는 사례다. \href{https://www.jstor.org/stable/1805621}{논문 서지}.\par

\subsubsection{이야기의 순서: 기계에서 숙련으로}

초기 타자기는 기계 구조와 특허, 판매망에 따라 배열이 달랐다. 어느 배열이 먼저 보급되느냐는 처음부터 완성된 최적화 문제로 해결되지 않았다. QWERTY 기계를 쓰는 사무실이 늘면 그 배열로 훈련된 사람을 채용하기 쉽다. QWERTY를 배운 사람은 같은 배열의 기계가 많은 곳에서 더 높은 생산성을 발휘하므로 그 배열을 더 배우려 한다. \textbf{기계의 설치 기반과 사람의 인적자본}이 서로를 강화한다.\par

이 과정은 첫 기계가 반드시 가장 좋았다는 증거도, 반드시 최악이었다는 증거도 아니다. 채택 과정의 설명과 효율성 평가는 구별해야 한다. 데이비드가 강조하는 것은 분권화된 현재의 결정이 모두 합리적이어도 사회 전체의 표준 전환이 어려울 수 있다는 가능성이다.\par

\subsubsection{세 가지 메커니즘}

\paragraph{1. 기술적 보완성}

기계와 사용자의 숙련은 함께 생산성을 낸다. 새 배열의 기계만 공급되어도 이를 다룰 직원이 없으면 기업은 구매하지 않는다. 직원만 새 배열을 배워도 기계가 없으면 훈련 수익을 얻지 못한다. 이는 두 시장이 동시에 움직여야 하는 조정 문제다.\par

\paragraph{2. 네트워크 외부성}

같은 배열의 사용자가 많아지면 교육 기관, 채용 시장, 수리·제조업체가 그 배열에 맞춰진다. 신규 사용자는 자신의 선택이 다른 사람에게 주는 호환성 편익을 모두 보상받지 못한다. 네트워크 효과가 크면 초기에 조금 앞선 배열이 더 빨리 커진다.\par

\paragraph{3. 전환 비용과 되돌리기 어려운 투자}

기계를 교체하는 비용만 생각하면 현대 컴퓨터에서 배열을 바꾸는 일은 쉬워 보인다. 하지만 이미 훈련받은 사람의 재학습 시간과 전환 중 생산 손실은 여전히 비용이다. 이전 훈련비 자체는 sunk cost이므로 현재 결정에 다시 더하면 안 된다. \textbf{앞으로 발생할 재훈련 비용}만 현재의 전환 판단에 관련된다.\par

\subsubsection{간단한 임계점 모형}

새 사용자가 표준 $A$ 또는 $B$를 고른다고 하자. $v_j$는 배열 자체의 생산성, $n_j$는 이미 그 표준을 쓰는 사람 수, $\lambda$는 네트워크 효과, $c_j$는 새 사용자의 전환·훈련 비용이다.\par

\[
u_A=v_A+\lambda n_A-c_A,\qquad
u_B=v_B+\lambda n_B-c_B.
\]

$B$가 기술적으로 더 좋아 $v_B>v_A$여도, 신규 사용자는 $u_A\geq u_B$이면 $A$를 선택한다. 이를 정리하면 $B$의 기술 우위가 $A$의 기존 사용자 우위와 추가 훈련 비용을 넘어야 한다.\par

\[
v_B-v_A>\lambda(n_A-n_B)+(c_B-c_A).
\]

이 식은 \textbf{개인의 현재 선택}을 나타낸다. 사회 전체가 동시에 바꿀 때의 장기 편익과 비용을 합한 사회적 판단은 별도다. 따라서 현재 $A$를 계속 선택한다고 해서 $A$가 반드시 사회적으로 최적이라는 결론은 나오지 않는다.\par

\subsubsection{비에르고드성과 역사적 우연}

같은 현재의 기술 후보가 있어도 초기 사건의 순서가 다르면 다른 표준에 도달할 수 있다. 이를 경로의존성이라고 한다. 더 엄밀한 논의에서 비에르고드성은 초기의 역사적 사건이 장기 결과의 분포에 사라지지 않을 수 있음을 뜻한다. 그러나 \textbf{경로의존성 일반과 Pareto 비효율적 lock-in은 다른 주장}이다. 역사에 의존해도 결과가 효율적일 수 있다. 데이비드의 \href{https://siepr.stanford.edu/publications/working-paper/path-dependence-foundational-concept-historical-social-science-revised}{후속 개념 정리}는 이 구분을 강조한다.\par

\subsubsection{실증 논쟁: QWERTY가 실제로 열등한가}

QWERTY와 Dvorak의 속도 비교는 훈련시간, 피험자의 기존 숙련, 과업, 실험을 후원한 기관에 민감하다. \href{https://www.journals.uchicago.edu/doi/abs/10.1086/467198}{Liebowitz와 Margolis의 반론}은 QWERTY의 열등성이 통설만큼 강하게 입증되지 않았다고 주장한다. 이 반론이 네트워크 효과라는 이론 자체를 없애는 것은 아니다. \textbf{이론적 가능성, 키보드의 역사 서술, 효율성에 관한 실증 판단}을 세 층으로 구분해야 논쟁이 명료해진다.\par

\subsubsection{응용: 다른 표준은 언제 바뀌나}

파일 형식, 결제망, 철도 궤간에서도 설치 기반과 숙련은 보완적이다. 하지만 전환을 돕는 변환기, 표준 간 호환성과 공동의 전환 일정이 있으면 비용이 줄어든다. 정책적으로는 기존 표준을 금지할지보다 \textbf{상호운용성과 정보 제공으로 조정 문제를 줄일 수 있는지}가 더 구체적인 질문이다.\par

\subsubsection{점검 질문}

1. QWERTY가 널리 쓰인다는 사실과 QWERTY가 최적이라는 주장은 왜 다른가?\par
2. sunk cost와 앞으로 지불할 전환 비용을 어떻게 구분하는가?\par
3. 네트워크가 강할 때 개인의 합리성과 사회적 효율성이 왜 어긋날 수 있는가?\par


\subsection{읽기 · Cliometrics and the Nobel}
\subsubsection{논문의 위치}

1993년 Robert Fogel과 Douglass North가 노벨경제학상을 받았을 때, Claudia Goldin은 두 사람이 무엇을 바꾸었는지 경제사 내부의 역사로 정리했다. Cliometrics는 역사 자료에 \textbf{경제 이론과 정량 방법을 결합}하는 연구다. 논문의 초점은 두 사람의 모든 결론이 최종적으로 옳았다는 찬사가 아니라, 과거에 대한 큰 주장도 계산 가능한 반사실과 자료의 비판을 거쳐야 한다는 연구 방식이다. \href{https://www.aeaweb.org/articles?id=10.1257/jep.9.2.191}{논문 페이지}.\par

\subsubsection{1960년대의 새 경제사}

기존 서술은 철도·노예제·기업가정신처럼 큰 주제를 다뤘지만, “철도가 성장에 필수였다” 같은 말을 정확히 어느 크기의 기여로 이해해야 하는지는 불분명했다. 젊은 연구자들은 신고전파의 가격·대체·기회비용 개념과 데이터를 사용해 통념을 검증했다. 정량화의 의의는 역사 이야기를 없애는 것이 아니라 \textbf{그 이야기가 참이라면 어떤 숫자가 관찰되어야 하는지}를 묻는 데 있었다.\par

\subsubsection{포겔: 철도가 없었다면}

철도가 없었을 때의 대안은 운송 자체가 0인 세계가 아니다. 수로와 도로를 이용하되 경로와 비용이 바뀐 세계다. 철도가 있는 상태의 단위 운송비를 $p_R$, 대안 운송비를 $p_A$로 두고, 운송 수요를 $q(p)$라고 하자. 대안 비용이 더 높다면 소비자에게 생긴 사회적 절감은 단순화해 다음과 같다.\par

\[
SS=\int_{p_R}^{p_A} q(p)\,dp.
\]

수요량이 가격에 따라 변하지 않는다고 가정하면 $SS=(p_A-p_R)Q$가 된다. 그러나 운송비가 오르면 멀리 보내는 물량과 생산지가 바뀔 수 있으므로 고정 물량 계산은 보통 거친 근사다. 강·운하가 어느 지역에 있었는지, 얼어붙는 계절은 언제인지, 철도가 없을 때 추가 수로를 건설할 수 있었는지가 결과에 들어간다.\par

포겔의 큰 교훈은 철도의 이익이 작았다는 단순 결론보다 \textbf{“필수불가결”이라는 말을 대체 가능성과 기회비용의 언어로 바꾸었다}는 점이다. 철도가 유용했어도 다른 교통이 성장의 일부를 담당할 수 있었다. 사회적 절감의 추정 범위와 전제는 후속 연구에서 논쟁될 수 있다.\par

\subsubsection{노스: 교역과 제도의 비용}

North는 19세기 해운과 미국 성장, 더 길게는 제도의 변화에 관심을 두었다. 운송비가 내려가는 이유를 선박의 물리적 속도만으로 찾지 않는다. 해적·전쟁 위험의 관리, 보험, 항구 제도, 계약 집행과 교역 조직이 비용을 바꾼다. 이를 거래비용이라고 부를 때에는 \textbf{정보를 얻고, 약속을 지키게 하고, 분쟁을 해결하는 실제 비용}을 뜻한다.\par

제도는 법률 문구만이 아니다. 관행과 신뢰, 권력자의 약속을 믿을 수 있는 정도도 투자 결정을 바꾼다. 다만 재산권이 좋아진 나라가 성장했다는 상관만으로 제도가 원인이라고 결론 내리기 어렵다. 성장한 경제가 더 나은 법원을 만들었을 수도 있으므로 변화의 시점과 기제를 찾아야 한다.\par

\subsubsection{포겔과 엥거먼의 노예제 연구}

남북전쟁 전 미국 남부의 플랜테이션이 어떤 수익성과 생산성을 가졌는지는 노예제가 곧 비효율 때문에 사라졌다는 도식에 도전했다. 그러나 플랜테이션의 산출 효율을 측정하는 일과 제도의 도덕적 정당성을 평가하는 일은 다른 질문이다. \textbf{강제·가족 분리·폭력과 자유 상실은 생산량의 높은 숫자로 지워지지 않는다.} 산출 비교에도 토지·작물·노동시간과 폭력의 비용이 들어가야 한다. 이 연구와 해석은 강한 비판을 받아 왔음을 기억해야 한다.\par

\subsubsection{몸의 역사: 신장과 생존}

포겔의 후기 연구는 소득 통계 밖의 생활 수준으로 확장된다. 키는 어린 시절의 영양과 질병 환경을 반영할 수 있고 연금·징집·인구조사 기록은 여러 시점의 개인을 연결할 가능성을 준다. 그러나 징집된 사람만의 신장, 측정 기준의 변화, 경기별 입대자의 차이는 편향을 만든다. 평균 키 한 수치가 모든 사람의 복지를 뜻하지 않는다. \href{https://www.nobelprize.org/uploads/2018/06/fogel-lecture.pdf}{포겔의 노벨 강연}은 영양·생리와 장기 성장의 관계를 더 풀어 설명한다.\par

\subsubsection{노벨상 이후의 역설}

경제사 연구가 경제학의 수학·계량 방법을 능숙하게 쓰면서 경제학의 여러 하위 분야로 흩어졌다. 성공적으로 주류 방법을 공유한 만큼 “경제사만이 제공하는 무엇인가”라는 정체성 질문도 남았다. 골딘은 장기적 사건에서 노동·건강·금융·제도를 동시에 보아야 한다는 넓은 시야를 강조한다.\par

\subsubsection{세 논문을 연결하면}

맥클로스키는 과거가 경제학에 \textbf{왜 유용한가}를, 골딘은 정량적 반사실이 그 유용성을 \textbf{어떻게 실천했는가}를, 그라이프는 초기에 선택한 이론이 \textbf{어떤 질문을 가렸는가}를 묻는다. 하나의 방법이 완성된 답이 아니라 점검 가능한 주장과 새로운 질문을 만드는 과정이다.\par

\subsubsection{점검 질문}

1. 철도가 없다는 반사실에서 수로와 도로를 고려해야 하는 이유는 무엇인가?\par
2. 노예제의 수익성·생산성과 제도의 정당성은 왜 독립적인가?\par
3. 신장 자료가 소득 자료를 보완하면서도 새로운 선택 편향을 만드는 이유는 무엇인가?\par


\subsection{읽기 · Cliometrics after 40 years}
\subsubsection{40년 뒤에 다시 묻는 방법}

그라이프는 클리오메트릭스가 과거에 이론과 정량 분석을 적용한 성과를 인정한다. 그러나 초기 연구에 쓰인 경제 이론은 당시 이용 가능한 컴퓨터와 통계만큼이나 강한 제약을 주었다. \textbf{어떤 이론을 도구로 삼느냐가 어떤 역사 현상을 설명 대상으로 볼 수 있는지를 결정}한다는 것이 논문의 핵심 문제다. \href{https://web.stanford.edu/~avner/Greif_Papers/1997%20Clio%20AER.pdf}{저자 제공 원문}.\par

\subsubsection{초기 신고전파 모형이 잘하는 일}

개인에게 일관된 선호가 있고, 가격에 반응하며, 계약과 정보의 비용이 낮다고 가정하면 자원 배분의 방향과 교역 이익을 명확하게 계산할 수 있다. 철도 비용 비교처럼 \textbf{이미 존재하는 시장의 상대가격}이 중심인 문제에는 강력하다. 그러나 이 모형을 모든 역사 질문에 곧바로 적용하면 시장이 왜 특정 방식으로 형성되었는가를 설명하기 어렵다.\par

\subsubsection{가정이 가리는 현상}

완전정보와 낮은 거래비용을 기본값으로 놓으면 낯선 상인끼리 왜 길드와 평판망을 만들었는지가 이상한 잔재처럼 보인다. 계약이 자동 집행된다고 가정하면 법원·국가·종교 공동체의 역할은 생산함수 바깥의 잡음이 된다. 동일한 선호와 단일 균형을 가정하면 같은 기술을 가진 도시가 왜 다른 제도 아래에 오래 머무는지도 설명하기 어렵다.\par

그라이프가 시장 밖 조직을 옹호하기만 하는 것은 아니다. 길드가 정보와 보험을 제공하면서 외부인의 진입을 막을 수 있듯, \textbf{제도의 편익과 분배 효과는 동시에 분석}해야 한다.\par

\subsubsection{게임이론이 여는 질문}

계약 위반을 국가 법원이 즉시 처벌하지 못하는 교역을 생각하자. 상인이 오늘 속여서 얻는 이익이 $G$, 앞으로 정직한 회원으로 남아 매기 얻을 수 있는 이익이 $R$, 미래의 가치를 오늘로 할인하는 인자가 $0<\delta<1$이라고 하자. 한 번 속이면 공동체에서 배제된다는 신뢰할 만한 규칙이 있다면 정직할 조건은 다음과 같다.\par

\[
G\leq \delta R+\delta^2R+\cdots=\frac{\delta R}{1-\delta}.
\]

이 식은 제재가 강력해서가 아니라 \textbf{미래 거래의 가치가 현재 일회성 이익보다 클 때} 협력이 유지된다는 직관을 보여 준다. 실제 역사에서는 속임수를 누가 알아내는지, 회원이 배제 결정을 함께 지키는지, 다른 시장으로 옮길 수 있는지가 조건에 들어간다. 따라서 제도는 단순 법 조항보다 조직과 신념을 포함한다.\par

\subsubsection{복수 균형과 경로의존성}

같은 기술·자원 아래에서도 모두가 계약을 지킬 것으로 기대하는 시장과 모두가 상대를 불신하는 시장이 각각 자기강제적일 수 있다. 기대는 다른 사람의 행동을 바꾸고, 그 행동이 다시 기대를 확인한다. 초기 사건이나 정치적 선택이 한 경로를 유리하게 만들어 오래 남을 수 있다. 데이비드의 기술 표준 사례와 닮았지만, 제도에서는 \textbf{권력과 분배, 정당성에 대한 믿음}이 더 직접적으로 작동한다.\par

복수 균형이라는 말을 “어떤 결과든 사후에 설명할 수 있다”는 면허로 쓰면 안 된다. 각 균형이 지속되려면 어떤 제재·정보·기대가 있어야 하는지 미리 적고 자료에서 그 증거를 찾아야 한다.\par

\subsubsection{역사 비교를 통한 판별}

어떤 상인 집단이 친족과 종교 공동체를 통해 거래했다면, 정보 문제가 핵심인 모형은 평판 공유와 위반자 제재의 기록을 예측한다. 정치권력 모형은 허가권과 조세 이익을 가진 집단의 저항을 예측한다. 같은 “길드가 있었다”는 사실만으로 두 이론을 고를 수 없다. \textbf{각 이론이 다르게 예측하는 자료}를 찾는 것이 역사비교제도분석의 힘이다.\par

\subsubsection{초기 연구의 성과를 버리지 않는 방법}

그라이프의 비판은 수량 분석을 폐기하자는 요구가 아니다. 더 풍부한 미시이론과 자료 비판을 결합하자는 요구다. 철도의 운송비나 임금의 반사실 계산과, 시장 형성 이전의 조직·문화·권력 분석은 서로 보완한다. 경제사가 유용해지려면 역사적 사례가 이론의 \textbf{적용 대상}일 뿐 아니라 이론의 \textbf{가정 자체를 수정할 근거}가 되어야 한다.\par

\subsubsection{비교하며 읽기}

\noindent\textbullet 맥클로스키: 과거가 경제학에 주는 다섯 가지 효용은 무엇인가?\par
\noindent\textbullet 골딘: 클리오메트릭스는 어떤 통념을 수량적으로 시험했는가?\par
\noindent\textbullet 그라이프: 그 시험에 사용한 이론이 어떤 역사적 질문을 놓쳤는가?\par

세 논문을 이어 보면 “더 많은 자료 → 더 정교한 반사실 → 더 넓은 이론”이라는 방법의 발전이 보인다.\par


\section{인구: 생존, 이동, 성장}
인구는 산출을 나누는 분모이면서 노동, 수요, 전염과 이주의 주체다. 출생·사망·이동의 회계식에서 시작해 장기 성장과 세계적 인구 이동을 분석한다.\par

\subsection{슬라이드 01 · 인구를 경제사의 출발점으로}
인구는 생산을 나누는 분모이면서 생산을 담당하는 노동력이다. 따라서 인구 증가가 곧 풍요도 빈곤도 아니다. 토지·기술·자본이 그대로일 때와 함께 변할 때 결과가 달라진다. 이 강의는 출생·사망·이동을 각각 분해하고, 1인당 생활 수준과 생산 구조에 어떤 효과를 주는지 추적한다.\par

\subsection{슬라이드 02 · 인구경제학의 범위}
출산은 가구의 선택이고, 사망은 영양·질병·공공재의 결과이며, 이주는 임금 차이와 이동 비용에 반응한다. 경제학이 인구를 다룬다는 것은 사람을 숫자로만 환원한다는 뜻이 아니다. 생존과 가족 결정이 당시의 소득·제도·권력 구조 속에서 어떻게 가능했는지 묻는 것이다. 현재의 저출산을 과거의 다산과 단순 비교할 때 생존율과 교육비 차이를 놓치기 쉽다.\par

\subsection{슬라이드 03 · 맹자의 인구 문제}
흉년 지역에서 다른 지역으로 사람을 옮겨도 한 나라 전체의 식량과 조세 제도가 그대로라면 구조적 문제는 남는다. 이 일화는 인구의 \textbf{공간적 재배치}와 \textbf{전체 생활 수준}을 구별하게 한다. 현대에도 지방에서 수도권으로 이동하면 개인의 기회는 개선될 수 있지만 주거·교통과 지역 간 재정 문제는 새로 생긴다.\par

\subsection{슬라이드 04 · 분모로서의 인구}
1인당 GDP는 $y=Y/N$이다. 총산출 $Y$가 커져도 인구 $N$이 더 빨리 늘면 평균 산출은 줄어든다. 미분하면 $\dot y/y=\dot Y/Y-\dot N/N$가 된다. 총량은 국가의 세수·군사력과 관련되고 1인당 수치는 평균 생활 수준과 더 가깝다. 다만 소득분배와 비시장 가사노동까지 알려 주지는 않는다.\par

\paragraph{성장률을 차근차근 계산하기}

$y=Y/N$의 양변에 로그를 취하면 $\ln y=\ln Y-\ln N$이다. 시간을 따라 미분할 때 $d\ln x/dt=\dot x/x$이므로 $g_y=g_Y-g_N$을 얻는다. 연도별 자료의 \textbf{정확한} 이산식은 $1+g_y=(1+g_Y)/(1+g_N)$이다. 예를 들어 총산출이 3\%, 인구가 2\% 늘면 1인당 성장률은 $1.03/1.02-1\approx0.98\%$다. 1\%라는 뺄셈은 작은 성장률에서 유용한 근사다.\par

\subsection{슬라이드 05 · per capita의 여러 의미}
1인당 산출, 소득, 소비는 분자가 다르다. 생산이 국내에서 이루어져도 이윤이 해외로 지급되면 국내총생산과 국민소득이 달라질 수 있다. 가계 소비가 늘지 않고 투자와 군사비가 늘면 1인당 GDP가 상승해도 체감 생활 수준은 다를 수 있다. 평균이므로 부유층의 소득 집중도 함께 확인해야 한다.\par

\subsection{슬라이드 06 · 맬서스와 고전파의 갈림길}
맬서스는 식량 생산의 제약과 인구 반응을 통해 빈곤의 지속을 설명했다. 리카도도 수확체감을 다뤘지만 토지 지대와 분배에 더 관심을 두었다. 케인스와 맬서스를 곧바로 한 계보로 묶는 것은 총수요 부족에 대한 문제의식이라는 제한된 의미에서 이해해야 한다. 맬서스의 인구론과 케인스의 유효수요 이론은 같은 모형이 아니다.\par

\subsection{슬라이드 07 · 학파 계보를 사용할 때}
학파의 계보는 주요 질문을 정리하는 지도이지 단일한 진화의 사다리가 아니다. 중상주의의 국가·무역, 중농주의의 토지, 고전파의 분배, 신고전파의 한계 선택이 서로 다른 시대의 문제를 반영한다. 어떤 연구자의 주장이 이후 학파의 모든 주장과 일치한다고 생각하지 말고, \textbf{공유한 질문과 달리 제시한 답}을 따로 읽자.\par

\subsection{슬라이드 08 · 맬서스 함정의 직관}
맬서스가 제시한 기하급수적 인구와 산술급수적 식량의 대비는 기제의 그림이지 모든 시대의 실측 성장 법칙은 아니다. 핵심은 고정 토지에서 노동이 증가하면 1인당 식량과 임금이 감소하고, 생존·출산 반응이 다시 인구를 바꾼다는 점이다. 기술이 지속적으로 발전하거나 출산 행동이 바뀌면 장기 결과가 달라진다.\par

\subsection{슬라이드 09 · 생산함수로 보는 함정}
토지 $T$가 고정되고 $Y=A T^{\alpha}N^{1-\alpha}$라고 하자. 1인당 산출은 $y=A(T/N)^\alpha$다. 기술 $A$가 한 번 상승하면 $y$가 먼저 오른다. 출생이 사망을 웃돌아 $N$이 증가하면 $(T/N)^\alpha$가 작아져 소득이 되돌아간다. 균형 생계수준 $\bar y$가 고정이라면 $N^*=T(A/\bar y)^{1/\alpha}$이므로 기술 개선은 주로 \textbf{부양 가능한 사람 수}에 남는다.\par

\paragraph{단기 충격에서 장기 균형까지}

생산함수를 사람 수로 나누면 $Y/N=A T^\alpha N^{-\alpha}$가 된다. 따라서 토지와 기술이 고정일 때 인구가 1\% 늘면 1인당 산출은 대략 $\alpha$\% 감소한다. 안정 소득을 $\bar y$라고 가정하고 $\bar y=A(T/N^*)^\alpha$를 풀면 먼저 $(\bar y/A)^{1/\alpha}=T/N^*$, 이어서 $N^*=T(A/\bar y)^{1/\alpha}$가 된다. $A$가 영구적으로 상승해도 출산·생존 반응이 충분히 강하면 새 교점은 오른쪽으로 이동하고 장기 소득은 다시 $\bar y$가 된다. 이 결론에는 \textbf{고정 토지, 일정한 인구 반응, 지속 성장하지 않는 기술}이라는 조건이 필요하다.\par

\subsection{슬라이드 10 · 인구 회계식}
기간 중 인구 변화는 $\Delta N=B-D+I-E$다. $B,D$는 출생과 사망, $I,E$는 유입과 유출이다. 두 항의 순차이가 같아도 구성은 다를 수 있다. 출생과 사망이 모두 높은 사회는 둘 다 낮은 사회보다 가족의 생존 경험과 연령 구조가 크게 다르다. 강의 전반에서 자연증가와 순이동을 혼동하지 않도록 먼저 분해하자.\par

\subsection{슬라이드 11 · 조출생률과 조사망률}
조출생률 $b=1000B/\bar N$, 조사망률 $d=1000D/\bar N$에서 $\bar N$은 보통 연평균 인구다. 단위는 퍼밀(‰), 즉 인구 천 명당 건수다. 조율은 계산이 쉬우나 연령 구성에 영향을 받는다. 고령자가 많은 나라의 조사망률이 높아도 같은 연령에서의 사망 위험은 오히려 낮을 수 있다. 비교에는 연령별 사망률과 표준화율이 필요하다.\par

\subsection{슬라이드 12 · 출생·사망률의 범위}
높은 출산력 집단의 관찰값은 생물학적 가능성과 사회 전체의 지속 가능한 출생률을 구분하게 한다. 혼인 연령, 모유 수유, 산모 건강, 피임, 유아 생존이 모두 실제 출생률을 바꾼다. “최대치”를 모든 사람에게 적용 가능한 일정한 상수로 보지 말고, \textbf{특정 집단과 기간에서 관찰된 상한의 근거}를 확인해야 한다.\par

\subsection{슬라이드 13 · 자연증가율의 단위}
조출생률 50‰, 조사망률 30‰이면 자연증가율은 $(50-30)/1000=0.02$, 즉 연 2\%다. 인구 이동이 없다면 $N_{t+1}\approx1.02N_t$이고, 복리로 $t$년 뒤 $N_t\approx N_0(1.02)^t$다. 천 명당 20명을 곧 20\%라고 읽는 단위 오류가 흔하다. 순이동률이 있으면 이를 별도로 더해야 한다.\par

천 명당 출생 50명에서 사망 30명을 빼면 순증가 20명이다. 이를 다시 천 명으로 나누어 $20/1000=0.02$를 얻는다. 2\%가 매년 유지되면 1년 뒤 $1.02N_0$, 2년 뒤 $1.02^2N_0$가 된다. 두 배가 되는 기간은 $(1.02)^t=2$의 양변에 로그를 취해 $t=\ln 2/\ln 1.02\approx35$년이다. 실제 역사에서 율은 해마다 바뀌므로 이 수치는 \textbf{일정한 율이라는 가정 아래의 비교 기준}이다.\par

\subsection{슬라이드 14 · 적극적 억제와 예방적 억제}
적극적 억제는 높은 사망률이 인구를 제한하는 경로, 예방적 억제는 만혼·독신·출산 조절 등으로 출생이 줄어드는 경로다. 둘 다 $b(y)=d(y)$라는 안정 조건을 만들 수 있어도 안정 소득은 다를 수 있다. 이 도식은 역사적 제도를 설명하는 출발점이며, 실제 사회에서 출산 선택을 소득만의 함수로 환원해서는 안 된다.\par

\subsection{슬라이드 15 · 사망률이 균형을 정할 때}
소득이 증가할수록 사망률 $d(y)$가 내려가고 출생률 $b$가 높은 수준에서 일정하면 두 곡선의 교점 $y^*$가 장기 소득이다. $y>y^*$에서는 $b>d$로 인구가 늘어 토지/사람 비율이 낮아지고, $y<y^*$에서는 반대다. 안정성은 이 방향의 피드백에서 나온다. 그래프의 교점은 높은 사망을 “바람직하다”고 평가하는 점이 아니다.\par

\subsection{슬라이드 16 · 일시적 소득 증가의 반응}
새 기술이나 좋은 작황으로 현재 소득이 $y^*$보다 높아지면 출생이 사망보다 많아질 수 있다. 이때 인구는 즉시 점프하지 않고 여러 세대에 걸쳐 증가한다. 단기적인 풍요와 장기 균형을 같은 시점으로 읽으면 동학을 잃는다. 초기 충격이 지속적 생산성 성장으로 이어지면 인구 증가가 이를 완전히 상쇄하지 못할 수도 있다.\par

\subsection{슬라이드 17 · 출생 행동이 바뀌면}
예방적 억제로 고소득에서 출생률이 떨어지는 함수 $b(y)$를 생각하자. 사망률 곡선과의 교점이 더 높은 $y$에 놓일 수 있다. 여기서 중요한 것은 단순히 “아이를 적게 낳으면 부유해진다”가 아니다. 생존율, 여성 교육, 노동시장, 양육비와 자녀에 대한 기대가 함께 변하며 출산과 소득 사이에는 양방향 인과가 있다.\par

\subsection{슬라이드 18 · 근대적 경제성장의 탈출}
기술·인적자본·조직의 생산성이 인구 증가보다 충분히 빨리 상승하면 $y=Y/N$이 지속적으로 높아진다. 한 차례의 생산성 상승과 매년 반복되는 상승을 구분해야 한다. 전자는 맬서스적 반응에 의해 흡수될 수 있고, 후자는 인구 조정 중에도 1인당 소득을 밀어 올릴 수 있다. 인구변천으로 출생률까지 하락하면 1인당 자본과 교육 투자에도 여유가 생긴다.\par

\subsection{슬라이드 19 · 세계경제사 그래프의 축}
장기간 거의 평평한 1인당 소득선과 근대의 상승은 “언제 탈출했는가”라는 질문을 만든다. 그러나 로그축, 구매력 평가, 고대의 희박한 관찰치를 어떻게 보간했는지에 따라 기울기의 인상이 달라진다. 영국의 선행과 다른 지역의 후발을 한 방향의 자동 추격으로 가정하지 말고 기술 이전과 제도적 장벽을 따져야 한다.\par

\subsection{슬라이드 20 · 경기순환과 추세}
한 해의 흉작이나 전염병은 소득을 추세 아래로 떨어뜨리지만, 장기 성장률의 지속적 변화와 같지 않다. 추세를 추정할 때에는 시작과 끝의 특이한 해, 전쟁과 물가 환산을 조심해야 한다. 반대로 긴 평균을 쓰면 사람들의 실제 위기 경험이 지워질 수 있다. 장기 경향과 단기 충격을 서로 다른 시간 척도로 읽자.\par

\subsection{슬라이드 21 · 거대한 분기의 뜻}
Great Divergence는 유럽 일부와 아시아 주요 지역의 생활 수준이 언제, 왜 크게 벌어졌는지를 둘러싼 논쟁이다. 처음부터 서유럽 전체가 아시아 전체보다 부유했다고 단정할 수 없다. 도시 임금, 식품 소비, 노동시간과 1인당 GDP 복원은 서로 다른 지표다. 원인도 석탄·신대륙 자원·제도·기술의 상호작용으로 다뤄야 한다.\par

\subsection{슬라이드 22 · 한국의 장기 1인당 소득}
한국의 시계열에서 수준의 상승과 성장률의 상승은 다르다. 그래프가 일정한 가격 연도를 쓰는지, 식민지기의 국민계정과 전후 자료의 연결 방법이 무엇인지 확인하자. 한국의 take-off를 한 연도로 정하더라도 이전의 교육·토지·산업 기반을 무시하면 인과를 과도하게 단순화한다. 실질 GDP의 증가와 그 분배도 구분해야 한다.\par

\subsection{슬라이드 23 · 사망 충격의 여러 유형}
전염병, 기근, 전쟁, 기후 충격은 사망자 수뿐 아니라 나이별 사망, 이동, 생산, 가족 형성을 함께 바꾼다. 맬서스 모형의 “적극적 억제”는 인구가 줄면 토지/사람 비율이 오를 수 있음을 말하지만, 실제 재난은 자본과 제도도 파괴한다. 따라서 재난 뒤 임금 상승을 기계적인 축복으로 해석해서는 안 된다.\par

\subsection{슬라이드 24 · 현대 전염병과 역사 비교}
1918년 인플루엔자와 최근의 유행병은 병원체, 감염 연령, 의료기술, 이동망이 다르다. 단순 사망률 비교보다 초과사망, 노동 공급, 학교 폐쇄와 소득 손실을 구분해야 한다. 세계적 연결은 전파 속도를 높일 수 있지만 과학 정보와 대응 기술도 함께 확산시킨다. 한 사건의 결과를 다른 시대로 그대로 옮기지 말자.\par

\subsection{슬라이드 25 · 흑사병의 경제사적 위치}
흑사병은 인구 급감이 노동과 토지의 상대적 희소성을 어떻게 바꾸는지 관찰할 중요한 사례다. 그러나 중세 말의 변화가 병 하나로 시작되고 끝난 것은 아니다. 이전의 인구 압력, 장원의 의무, 도시 시장과 권력 관계가 충격의 결과를 조정했다. 다음 자료에서는 확산 경로와 사회적 반응을 따로 본다.\par

\subsection{슬라이드 26 · 페스트의 시기와 이름}
14세기 유행은 여러 차례의 파동이었고 사망도 지역마다 달랐다. “Black Death”라는 이름과 피부의 검은 반점만으로 전파 방식이나 치명률을 추정할 수 없다. 역사 기록, 고고학적 유골, 병원체 연구는 서로 다른 종류의 증거다. 사회경제적 효과를 따질 때에는 충격 전후의 인구와 임금 자료의 측정 불확실성을 함께 적어야 한다.\par

\subsection{슬라이드 27 · 기원에 관한 사료의 한계}
발병의 “최초”를 묻는 것은 현재 남은 기록 중 어디가 가장 이른지와 병원체가 실제로 어디서 생겼는지를 구분해야 한다. 유럽 기록에 먼저 보인 지점을 세계적 기원으로 착각하면 유럽 중심적 서술이 된다. 이동 경로의 복원에는 교역로, 선박, 동물 숙주, 유전적 증거가 필요하며 단일한 직선 화살표만으로는 부족하다.\par

\subsection{슬라이드 28 · 카파에서 유럽으로}
흑해 교역항 카파는 제노바 상인과 유라시아의 이동망이 만나는 곳이었다. 상인·선박·화물이 이동한 네트워크는 상품뿐 아니라 질병도 옮길 수 있었다. 카파에서 유럽으로 전파되었다는 서술은 유력한 역사적 경로 중 하나지만, 각 지역의 도착 시점과 전파 기제를 독립 자료로 확인해야 한다. 교역의 이익과 감염 외부효과가 공존한다.\par

\subsection{슬라이드 29 · 전파 지도 읽기}
지도상의 화살표는 확인된 모든 감염자 이동을 뜻하지 않는다. 항구와 내륙 도시의 연도 표기를 읽어 이동 속도와 연결망을 추정하고, 보고가 늦어진 지역도 생각해야 한다. 지리적으로 가까운 곳보다 교역으로 강하게 연결된 곳이 먼저 노출될 수 있다. 전파 지도는 경제적 네트워크의 간접적인 지도이기도 하다.\par

\subsection{슬라이드 30 · 그림 자료의 빈칸}
이미지나 지도가 단독으로 제시될 때에는 제작 시기, 표현 대상, 후대의 재구성 여부를 먼저 확인해야 한다. 전염병을 묘사한 그림은 사람들이 어떻게 두려움과 죽음을 이해했는지 보여 주지만, 발병률의 정확한 통계가 되지는 않는다. 시각 자료의 \textbf{증언하는 바와 증언하지 못하는 바}를 나누는 습관이 중요하다.\par

\subsection{슬라이드 31 · 생물학전 이야기의 검증}
포위군이 시신을 던졌다는 전승은 전쟁과 질병을 연결하는 강렬한 이야기다. 그러나 전승의 신뢰성과 실제 전파에서의 비중은 별도로 검증해야 한다. 한 사건이 알려진 전파 경로 전체를 설명하지 못할 수 있다. 특히 교역망·설치류·벼룩·사람 사이의 전파 가능성을 구분하고 단일한 기원담으로 고정하지 말자.\par

\subsection{슬라이드 32 · 증상 묘사의 용도}
그림의 반점은 당대의 질병 인식을 보여 주지만 현대 의학적 진단서가 아니다. 같은 시기의 다른 질병도 비슷하게 묘사될 수 있다. 역사 인구 연구에서는 증상 그림보다 매장 기록, 유언장, 장부의 공백과 가구 단위의 사망 흔적이 사망 규모를 추정하는 데 더 직접적일 수 있다.\par

\subsection{슬라이드 33 · 공동묘지와 사망 규모}
대량 매장은 충격의 심각성을 시각화한다. 하지만 한 도시의 묘지 그림으로 유럽 전체의 사망률을 계산할 수는 없다. 사망자가 집중된 계절과 지역, 매장 관행, 기록이 남지 않은 빈민을 고려해야 한다. 노동시장 효과도 같은 비율로 나타나지 않는다. 생존 노동자의 기술과 토지 이용 방식이 달라지기 때문이다.\par

\subsection{슬라이드 34 · 유대인 박해와 위기 대응}
질병 원인을 모를 때 소수집단에 책임을 떠넘기는 폭력은 경제 충격의 사회적 경로다. 거래·신용 네트워크의 붕괴와 재산 약탈도 뒤따를 수 있다. 이런 박해는 방역의 성과가 아니라 정보 부족과 권력, 편견의 결과로 분석해야 한다. 전염병의 경제사를 단순 노동 공급 곡선만으로 설명할 수 없는 이유다.\par

\subsection{슬라이드 35 · 채찍파의 행진}
고행 운동은 사람들이 재난을 종교적·도덕적 언어로 해석한 한 방식이다. 믿음은 이동, 집회, 기부, 폭력과 같은 실제 행동을 바꾸므로 경제사에서 배제할 수 없다. 현대의 질병 지식을 과거 행위자에게 미리 부여하면 왜 특정 대응이 설득력을 가졌는지 이해하기 어렵다. 동시에 그림 한 장을 모든 신앙인의 전형으로 일반화하지 말자.\par

\subsection{슬라이드 36 · 교역과 질병의 공동 확산}
교역량이 늘면 상품의 다양성과 도시 소득이 늘 수 있지만, 병원체가 이동할 경로도 늘어난다. 감염 위험은 거래 당사자 밖의 사람에게 비용을 주는 외부효과다. 무역을 금지하면 이익도 잃으므로 역사적으로 검역, 정보 공유, 항구 통제와 같은 절충이 등장했다. 흑사병을 세계화의 한 부산물로 볼 수 있지만 전파의 세부는 증거로 따져야 한다.\par

\subsection{슬라이드 37 · 시각 자료의 비교 기준}
시각 자료가 앞뒤의 전파·인구 도표 사이에 놓였다면 무엇을 연결하는지 질문하자. 이동 경로에서 사망 규모로 넘어갈 때 필요한 중간 변수는 접촉률, 감염 확률과 생존율이다. 그림의 상징적 의미와 정량 자료를 혼동하지 않으면 사건의 서사와 측정 가능한 결과를 함께 읽을 수 있다.\par

\subsection{슬라이드 38 · 대륙별 인구의 장기 변화}
600\~{}1600년의 인구 추정치는 세기별로 자료의 질이 크게 다르다. 흑사병 전후 유럽의 감소를 읽을 때 세계 전체의 인구 감소와 동일시하지 말자. 지역별 농업 생산력, 전염병, 이주가 서로 다른 경로를 만들었다. 그래프의 절대 인구와 비율을 구분하고, 관측값인지 역사적 추정값인지 확인하는 것이 중요하다.\par

\subsection{슬라이드 39 · 천연두와 신대륙}
신대륙의 천연두는 병원체에 대한 이전 노출과 면역 경험의 차이가 인구에 큰 충격을 줄 수 있음을 보여 준다. 그러나 원주민 인구 감소를 한 질병만의 결과로 쓰면 전쟁·강제노동·이주·기근을 지운다. 질병은 식민지 경제의 노동 수요와 강제적 노동 조직에도 영향을 주었고, 그 결과가 다시 생존에 영향을 미쳤다.\par

\subsection{슬라이드 40 · 총과 균을 함께 보기}
소수 정복자의 승리를 무기 하나나 질병 하나로 설명할 수 없다. 동맹, 정치적 분열, 군사기술, 감염과 보급이 결합했다. “균”은 의도 없이 퍼졌더라도 결과적으로 힘의 불균형을 키웠다. \textbf{인구 감소 → 노동·조세 기반 약화 → 지배력 변화}라는 경로를 검토하되, 지역마다 시점과 크기가 달랐음을 기억하자.\par

\subsection{슬라이드 41 · 멕시코 중부 인구 그래프}
1532\~{}1608년의 급감은 절대 숫자의 불확실성이 큰 사례다. 기준 연도의 인구 추정이 바뀌면 감소율도 달라진다. 그래프가 보여 주는 큰 방향과 정확한 수치를 구분하고, 질병·전쟁·이주·기록 범위의 변화가 각각 어떤 영향을 주는지 나누어 보자. 인구의 붕괴는 식민 경제의 노동 조달 방식과도 연결된다.\par

\subsection{슬라이드 42 · 원주민의 시각 자료}
원주민이 그린 질병 장면은 정복자 기록 외의 시각을 제공한다. 증상만이 아니라 가족·공동체의 돌봄과 충격을 읽을 수 있다. 다만 그림의 제작 목적과 전승 과정을 확인해야 하고, 개별 장면을 인구 전체의 임상 자료로 간주해서는 안 된다. 다양한 출처를 교차할 때 역사의 설명이 견고해진다.\par

\subsection{슬라이드 43 · 인구 이동을 분석하는 틀}
이주는 자연증가와 달리 출발지와 도착지의 인구를 반대 방향으로 바꾼다. 개인의 예상 소득 차이에서 이동 비용과 위험을 뺀 값이 선택에 영향을 주지만, 모든 이동이 자발적인 것은 아니다. 이후 사례에서는 \textbf{강제의 정도, 누가 이동 비용을 지불하는가, 이동 뒤 권리가 무엇인가}를 나누어 읽는다.\par

\subsection{슬라이드 44 · 인류 초기의 확산}
인류의 아프리카 기원과 확산은 매우 긴 기간의 생태 적응과 이동을 포함한다. 현대 국경의 “이민” 개념을 선사시대에 그대로 대입할 수 없다. 환경, 식량, 인구 압력과 기술이 이동 가능성을 바꾸었을 것이다. 고고학·유전학적 증거는 이동 경로를 추정하지만 한 지도의 화살표가 단일 집단의 직선 여행을 뜻하지는 않는다.\par

\subsection{슬라이드 45 · 세계 확산 지도}
지도의 연대는 최초 도달의 추정 시점이지 모든 지역에 균일한 정착이 이루어진 시점이 아니다. 이동에는 세대를 거친 확산, 재이동, 기존 집단과의 혼합이 있다. 넓은 시간·공간 축을 읽을 때에는 지도의 해상도보다 더 세밀한 경제적 이유를 단정하지 않도록 주의해야 한다.\par

\subsection{슬라이드 46 · 고대 말의 이동과 제도}
게르만족 이동은 정치적 경계와 토지·군사 조직을 바꾸었다. 이를 단순한 “민족의 대이동”으로만 보면 여러 집단의 계약·동맹·정착 과정을 놓친다. 이동은 노동 공급과 토지 점유, 조세 권리의 재배치이기도 하다. 다음 두 장에서 압력과 제국 내부의 수요를 함께 본다.\par

\subsection{슬라이드 47 · 훈족의 압력과 연쇄 이동}
훈족의 서진은 한 집단의 이동이 다른 집단의 선택지를 바꾸는 연쇄 효과를 보여 준다. 로마 제국의 쇠퇴도 외부 압력만이 아니라 군사·재정의 내부 문제와 결합했다. “밀려났다”라는 표현은 출발 요인을 담지만 각 집단이 어디에 정착하고 어떤 권리를 얻었는지까지 설명하지 못한다.\par

\subsection{슬라이드 48 · 이동 지도와 중세 유럽}
지도에서 집단명과 도착지의 화살표를 읽되, 하나의 집단이 본래부터 고정된 정체성을 유지했다고 가정하지 말자. 이동 뒤에는 현지 인구와 섞이고 기존 정치 제도를 사용하거나 변형했다. 경제사적으로 중요한 것은 토지 소유와 농민 의무, 도시 연결망이 어떻게 재편되었는가다.\par

\subsection{슬라이드 49 · 강제 이동이라는 별도의 범주}
대서양 노예무역은 높은 임금에 반응해 이동한 이민과 다른 범주다. 사람은 계약을 거절하거나 돌아갈 권리가 없었고, 이동 과정과 플랜테이션의 폭력이 비용의 핵심이었다. 경제적 분석을 할 때에도 거래된 인원 수만이 아니라 사망, 가족 분리, 자유 상실과 후대의 효과를 함께 다뤄야 한다.\par

\subsection{슬라이드 50 · 플랜테이션의 노동 수요}
사탕수수·담배·면화는 원격지 시장을 위한 대규모 상품 작물이었다. 토지 확보와 유럽 수요는 노동력의 조달 문제와 결합했고, 식민 권력은 강제노동을 제도화했다. 생산량의 증가가 현지인의 후생 증가를 뜻하지 않는다. 누가 토지와 상품의 소유권을 가졌고 누가 위험과 폭력을 감당했는지 따로 적어야 한다.\par

\subsection{슬라이드 51 · 아시엔토의 성격}
Asiento는 스페인 식민지에 노예를 공급하는 독점 계약권이었다. 국가는 독점권에서 재정 수입을 얻고 상인은 정치적 허가로 거래 이익을 얻었다. 이는 시장가격만으로 교역량이 정해지는 경쟁시장과 다르다. 계약권의 이전은 유럽 국가 간 전쟁·협상의 결과였으며, 피해를 입은 사람의 권리는 협상 대상에서 배제되었다.\par

\subsection{슬라이드 52 · 위트레흐트 조약과 영국}
1713년의 계약권 이전은 국제 조약이 특정 상업 독점을 어떻게 배분하는지 보여 준다. 독점의 법적 권리와 실제 운송·금융·항구 통제 능력이 결합해야 수익이 실현된다. 영국의 참여를 단순히 “무역 발전”이라 표현하면 강제 이송의 성격을 지우므로, 거래 주체와 피지배자의 위치를 분명히 구분한다.\par

\subsection{슬라이드 53 · 전쟁과 계약권의 변화}
스페인과 영국의 충돌은 독점 계약이 영구적인 재산이 아니라 정치적 권력에 의존한다는 점을 드러낸다. 전쟁은 선박 손실, 보험료, 항구 접근과 밀무역의 유인을 함께 바꾼다. 법적 계약의 명목상 수량과 실제 이송된 인원은 다를 수 있으므로 자료 출처와 기간을 확인해야 한다.\par

\subsection{슬라이드 54 · 규모 추정의 기준}
노예무역의 규모를 말할 때에는 아프리카에서 승선한 인원, 대서양을 살아 건넌 인원, 각 식민지에 도착한 인원을 구분해야 한다. 국내 강제 이동과 사하라 횡단·인도양 교역까지 포함하는지도 명시해야 한다. 숫자의 차이는 단순 오차가 아니라 \textbf{무엇을 셌는가}의 차이일 수 있다. \href{https://www.slavevoyages.org/}{SlaveVoyages 데이터베이스}가 항해 단위의 추정과 출처를 제공한다.\par

\subsection{슬라이드 55 · 출발지별 도표}
막대의 높은 지역이 그 지역만의 “자발적 공급”을 의미하지 않는다. 전쟁, 노예 사냥, 중개 무역과 해안 항구의 접근성이 기록된 출발지를 결정했다. 인구 손실의 경제 효과는 단순히 이송 인원만이 아니라 젊은 노동자의 선택적 이탈, 공동체의 불안정, 정치적 폭력의 유인을 포함한다.\par

\subsection{슬라이드 56 · 서점운동의 의미}
미국의 서부 이동은 토지 기회와 국가의 영토 확장이 결합한 사례다. 이주민에게 제공된 토지는 원래 비어 있던 땅이 아니었다. 원주민의 축출과 조약 위반, 노예제의 서진을 함께 보아야 한다. “인구 이동”의 통계 뒤에는 서로 다른 집단의 권리 충돌이 있다.\par

\subsection{슬라이드 57 · Westward Movement 지도}
지도에서 이동 방향과 새로 편입된 지역을 읽을 때, 국가가 인정한 소유권의 변화와 실제 거주자의 삶을 분리하자. 저렴한 토지 가격은 이주민의 기회였지만 그 가격이 원주민의 손실을 포함하지 않았다. 서부 이주는 농업 생산, 운송 수요, 도시 형성과 정치적 대표권도 함께 바꿨다.\par

\subsection{슬라이드 58 · 명백한 운명과 정책}
Manifest Destiny는 확장을 정당화한 정치적 언어다. 토지 조례와 국유지 매각은 누가 토지를 얻는지 결정하는 제도였고, 인구 기준의 주 승격은 정치적 힘을 바꿨다. 이념만으로 이주를 설명할 수 없지만, 이념이 강제와 재산권 재편을 합리화하는 역할도 분석해야 한다.\par

\subsection{슬라이드 59 · 콜럼비아의 이미지 읽기}
그림 속 빛과 통신·철도는 확장을 “진보”로 묘사하지만, 뒤로 밀려나는 사람들은 다른 경험을 보여 준다. 시각 자료를 경제성과의 중립적 기록으로 보지 말고 제작자의 설득 목적을 읽자. 같은 기술의 도입이 집단마다 이익과 손실을 달리 배분할 수 있다.\par

\subsection{슬라이드 60 · 골드러시의 유인}
금 채굴의 높은 예상 수익은 먼 이동과 탐색의 위험을 감수하게 한다. 하지만 대박 사례가 알려질수록 더 많은 진입자가 몰려 평균 수익이 낮아지고 운송·식량 가격이 오른다. 1849년의 이주를 이해하려면 발견된 금의 가치뿐 아니라 \textbf{이동 비용과 성공 확률의 불확실성}을 생각해야 한다.\par

\subsection{슬라이드 61 · 토지 분배와 가구 형성}
국유지 매각과 자작농 창설 정책은 농장 규모, 신용 접근과 이주 속도를 바꾼다. 토지를 받는 법적 조건과 실제로 필요한 이주 비용·농기구 비용이 다르므로 모든 빈민이 같은 기회를 가진 것은 아니다. 인구 증가가 주 승격의 기준이 되면서 토지정책은 정치적 대표권과도 연결되었다.\par

\subsection{슬라이드 62 · 면화왕국의 구조}
면화 수출의 성장은 영국 섬유공업의 수요와 미국 남부의 토지·강제노동을 연결했다. 수출 비중이 높다는 사실만으로 남부 주민 모두가 부유했다는 뜻은 아니다. 플랜테이션 소유자의 수익, 노예 노동자의 삶, 토양 고갈과 금융 위험은 분리해서 평가해야 한다.\par

\subsection{슬라이드 63 · 조면기와 보완 요소}
조면기는 섬유와 씨앗을 분리하는 비용을 낮춰 면화 공급을 늘렸다. 그러나 기계만으로 생산량이 결정되지는 않는다. 재배 토지, 노동력, 운송로와 영국의 수요가 함께 확대되었다. 생산성 기술이 등장해도 특정 집단의 강제노동 수요가 오히려 증가할 수 있다는 점은 기술진보가 자동으로 자유를 넓히지 않음을 보여 준다.\par

\subsection{슬라이드 64 · 노예 인구와 면화 생산}
두 시계열이 함께 증가하면 상호 보완 관계를 의심할 수 있지만 어느 쪽이 원인인지 그림만으로 결정할 수 없다. 세계 면화 가격과 토지 개간, 법적 노예제 확대가 두 변수 모두를 밀었을 수 있다. 생산성의 변화를 보려면 면화 산출을 노동자 수만이 아니라 토지와 자본도 고려해 나누어야 한다.\par

\subsection{슬라이드 65 · 계약노동의 강제성}
Indentured labor는 일정 기간의 노동으로 도항비나 채무를 갚는 계약으로 설명되지만, 정보 부족과 이동의 제약, 폭력 때문에 형식적 동의가 실질적 자유를 뜻하지 않을 수 있다. 노예제와 동일한 법적 지위로 단순화하지도, 완전한 자유계약으로 미화하지도 말자. \textbf{계약 종료권과 이동·재산·가족의 권리}를 비교해야 한다.\par

\subsection{슬라이드 66 · 채찍질 통계의 해석}
연평균 폭력 횟수는 강제의 한 단면을 보여 줄 뿐이다. 낮은 평균이라고 폭력이 드물었다거나 노동자가 동의했다는 결론은 나오지 않는다. 폭력의 위협만으로도 행동이 통제될 수 있고, 기록되지 않은 사건과 집단별 편차가 있다. 노동 인센티브를 논할 때 물리적 강제와 자유로운 임금 계약의 차이를 지워서는 안 된다.\par

\subsection{슬라이드 67 · 유럽의 해외 이주}
19세기의 대서양 이주는 노동이 많은 지역에서 토지와 노동 수요가 큰 지역으로 움직인 사례다. 그러나 임금 격차만으로 이동량을 설명할 수 없다. 선박 운임, 정보망, 언어, 이민법과 가족의 지원이 모두 선택을 바꾼다. 이주자가 젊고 건강한 사람에 치우치면 출발지와 도착지의 인구 구조도 달라진다.\par

\subsection{슬라이드 68 · 1820\~{}1913년의 두 물결}
초기와 후기 이주의 출발국이 달라졌다는 점은 유럽 내부의 개발 격차와 운송비 하락을 보여 준다. 어느 시기든 전체 유럽인이 같은 확률로 떠난 것이 아니다. 필요한 운임을 감당할 수 없는 극빈층은 오히려 이동하기 어려울 수 있다. 숫자를 읽을 때에는 순이민과 총입국, 귀환 이민의 차이도 확인해야 한다.\par

\subsection{슬라이드 69 · 출발국과 도착국의 비대칭}
지도나 표의 흐름은 인구가 많은 국가가 반드시 가장 많이 보내는 나라라는 뜻은 아니다. 임금 차이, 항구 접근, 이민 제한, 친척·동향인 네트워크가 경로를 집중시킨다. 출발국의 노동 공급이 줄면 임금이 오를 수 있지만, 기술과 자본도 함께 빠져나가면 지역별 효과가 달라진다.\par

\subsection{슬라이드 70 · Push와 pull}
아일랜드의 감자 기근은 생존을 압박하는 push 요인의 강한 사례다. 미국의 토지·임금 기회는 pull 요인이다. 두 요인은 동시에 작동하고, 출발지의 충격만으로 도착지를 선택할 수는 없다. 기근의 사망과 해외 이주를 합쳐 단순 인구 감소로만 읽으면 강제된 선택의 차이를 지운다.\par

\subsection{슬라이드 71 · 선박 운임과 회항 화물}
대서양을 건넌 선박은 양방향 운송 수요를 맞춰야 했다. 한쪽 방향의 화물이 부족하면 빈 배로 돌아오는 대신 승객 운임을 낮출 유인이 생긴다. 운임 하락은 이동의 고정 비용을 낮춰 더 많은 계층이 이주할 수 있게 한다. 다만 실제 비용에는 항구까지 이동, 정착 자금, 여행 중 위험이 포함된다.\par

\subsection{슬라이드 72 · 이주 증가 그래프}
미국 이주 증가와 유럽 인구 증가는 함께 나타나도 이주율과 절대 인원을 구별해야 한다. 유럽 인구가 늘면 이주자가 많아져도 인구 대비 비율은 다르게 움직일 수 있다. 연도별 큰 변동은 경기, 전쟁, 법률, 선박 기술의 영향을 받을 수 있으므로 장기 추세와 충격을 나누어 읽자.\par

\subsection{슬라이드 73 · 미국의 조상 집단}
오늘의 ancestry 통계는 19세기 입국 기록과 같은 개념이 아니다. 한 사람은 여러 조상을 선택할 수 있고, 세대가 지나면서 정체성의 표현도 바뀐다. 현재의 집단 규모를 과거의 순이민량으로 직접 역산해서는 안 된다. 누적 이주, 출산, 사망, 재이주와 조사 방법을 모두 고려해야 한다.\par

\subsection{슬라이드 74 · 직업 구성의 변화}
이주자의 직업 표는 출발지의 교육·산업 구조와 도착지의 수요를 반영한다. 직업 미기재자와 농민의 계절 노동을 어떻게 분류했는지도 중요하다. 의사가 이주하면 출발지는 인적자본을 잃고 도착지는 얻을 수 있지만, 송금·귀환·지식 이전이 이를 부분적으로 바꾼다. 단순한 “두뇌 유출” 하나로 평가하지 말자.\par

\subsection{슬라이드 75 · 엘리스섬 기록의 용도}
이민 기록에는 이름, 출발지, 가족 관계, 때로는 직업이 남아 개인과 가구의 이동을 복원할 수 있다. 이름의 철자 변화와 누락, 다른 항구를 통한 입국을 조심해야 한다. 행정자료가 허용한 관찰 범위와 전체 이주자의 범위를 구분하면 한 가문의 미시사와 거시적 이민 흐름을 연결할 수 있다.\par

\subsection{슬라이드 76 · 스타인웨이의 사례}
이민 기업가의 성공은 기술·숙련이 국경을 넘어 이동한다는 사례다. 그러나 성공한 기업 하나를 평균 이민자의 결과로 일반화하면 생존자 편향이 생긴다. 도착지의 시장 규모, 자본 접근, 특허와 유통망이 숙련이 사업으로 바뀌는 조건이었다. 개인 역량과 제도의 보완 관계를 보자.\par

\subsection{슬라이드 77 · 출발지와 도착지의 효과}
노동자가 떠나면 출발지의 노동 공급이 줄어 잔류자의 임금이 오를 수 있다. 그러나 가족 분리와 숙련 손실도 있다. 도착지는 노동과 기술을 얻지만 토지·주택·학교의 공급이 즉시 늘지 않으면 갈등이 생길 수 있다. 이주의 순효과는 평균 임금 한 개가 아니라 집단별 임금, 지대, 재정과 후대 성과로 나누어야 한다.\par

\subsection{슬라이드 78 · 아시아와 열대 지역의 계약 이주}
노예제 폐지 뒤 플랜테이션의 노동 수요가 사라진 것은 아니었다. 인도·중국 등지에서 계약노동자를 모집했지만 모집의 강제, 부채, 이동권 제한이 적지 않았다. 같은 “계약”이라는 용어 아래 법적 지위와 실제 자유가 크게 달랐다. 출발지·도착지의 국가 권력과 상인의 중개 수익을 함께 살펴야 한다.\par

\subsection{슬라이드 79 · 계약노동자 수의 비교}
출신 국가별 이주 인원을 비교하려면 기간, 목적지, 귀환자와 재이주자를 어떻게 세었는지 먼저 보자. 표의 국가명이 당시의 제국·식민 경계와 오늘의 국경이 다를 수 있다. 집계된 사람은 노동시장의 추상적 공급량이 아니라 계약 조건과 위험이 달랐던 개인들이었다.\par

\subsection{슬라이드 80 · 중국의 인구와 내부 이동}
중국처럼 넓은 지역에서는 해외 이민보다 국내의 개척·도시 이동이 인구 분포에 더 큰 영향을 줄 수 있다. 한 왕조의 총인구 증가만으로 동부와 서부의 토지 압력이나 생활 수준을 알 수 없다. 이후 지도들은 인구가 어느 지역으로 이동했고 농업의 노동 집약도가 어떻게 달라졌는지를 보여 준다.\par

\subsection{슬라이드 81 · 매디슨의 장기 추정}
역사적 GDP와 인구 추정은 서로 다른 출처를 연결한 재구성이다. 매디슨의 수치로 장기 질서를 비교할 수 있지만 고대·중세의 한 자리 숫자를 현대 국민계정처럼 읽으면 안 된다. 명목 환율이 아닌 구매력 기준의 1인당 소득과 총량을 구별하고, 불확실한 시대일수록 큰 방향에 집중하자. \href{https://www.rug.nl/ggdc/historicaldevelopment/maddison/}{Maddison Project Database}는 추정의 갱신판을 제공한다.\par

\subsection{슬라이드 82 · 교와 쇄국의 대비}
해외 교역과 이민의 개방 정도는 시대마다 달랐다. “쇄국”도 모든 사람·상품의 이동이 0이었다는 뜻으로 받아들이면 안 된다. 법적 제한, 실제 밀무역, 내부 이동과 국가의 해안 통제를 따로 확인해야 한다. 교역 정책 하나만으로 동서양의 소득 격차를 설명하기보다 농업·에너지·국가 재정과 연결해야 한다.\par

\subsection{슬라이드 83 · 14\~{}16세기의 연결망}
이 시기의 지도는 항로와 항구가 만들어 낸 상업 연결을 보여 준다. 국가가 원정과 교역을 지원한 시기에도 민간 상인의 움직임과 해외 공동체가 같은 방식으로 움직인 것은 아니다. 항로의 존재와 지속적인 대규모 이민, 생산성 향상은 각각 다른 명제다. 지도에서 목적지와 출발지를 구별해 보자.\par

\subsection{슬라이드 84 · 정화 원정의 경제적 질문}
대규모 함대를 보낼 수 있는 능력은 국가의 조선·재정·조직력을 보여 준다. 그러나 원정의 성공이 곧 민간 교역의 지속적 이익이나 산업혁명으로 이어지지는 않는다. 원정의 목적, 비용을 낸 집단, 후속 항해의 인센티브와 정책 변화가 무엇이었는지 묻자. 유럽의 이후 항해와 단순 선후 비교만으로 “기회를 놓쳤다”고 결론 낼 수 없다.\par

\subsection{슬라이드 85 · 해안 통제의 시기}
정책 변화를 표시한 지도에서 법률의 도입 시기와 효과가 나타난 시기를 구분하자. 해안 금지 정책이 강화되어도 내부 농업 개척과 인구 이동은 계속될 수 있다. 국가 정책은 해상 교역의 수익률을 바꾸지만, 수요·밀무역·국경의 집행 능력이 실제 거래량을 결정한다.\par

\subsection{슬라이드 86 · 거대한 분기의 설명 경쟁}
중국의 정책 전환을 Great Divergence의 단일 원인으로 놓으면 유럽의 석탄·식민지 자원, 중국의 지역별 시장과 농업 생산성을 놓친다. 대분기는 \textbf{언제부터 어느 지역 사이에서 어떤 지표가 벌어졌는가}가 먼저 정해져야 설명할 수 있다. “동양”과 “서양”을 각 한 점으로 표시하는 비교는 지나치게 거칠다.\par

\subsection{슬라이드 87 · 수전농업과 인구 부양력}
논은 면적당 많은 식량을 생산할 수 있고, 높은 인구 밀도를 유지할 수 있다. 그러나 토지당 산출이 높은 것과 노동자 한 사람당 산출이 높은 것은 다르다. 더 많은 노동을 투입해 생산량을 늘리면 토지 생산성은 올라가도 노동 생산성은 정체할 수 있다. 인구 부양력과 1인당 소득을 구별하는 핵심 사례다.\par

\subsection{슬라이드 88 · 쌀과 밀의 면적당 수확}
헥타르당 곡물 무게만 비교하면 열량, 이모작 가능성, 종자·관개·노동 투입을 빠뜨린다. 같은 면적에서 더 많은 사람을 먹일 수 있다는 주장은 토지 제약 아래 중요하지만, 그 사회가 더 높은 시간당 임금을 얻었다는 뜻은 아니다. 비교할 때에는 \textbf{열량/면적, 산출/노동시간, 가구당 소비}를 각각 계산해야 한다.\par

\subsection{슬라이드 89 · 근면혁명과 산업혁명}
Industrious revolution은 노동시간과 노동 집약도를 높여 가계 산출을 늘리는 경로, industrial revolution은 기계·에너지와 자본을 써 노동자당 산출을 높이는 경로로 대비된다. 두 경로는 배타적이지 않고 지역별 자원 가격과 시장 수요에 따라 섞인다. 노동이 풍부하고 자본이 비싸면 노동 집약적 기술이 사적으로 합리적일 수 있다.\par

\subsection{슬라이드 90 · 등량곡선으로 읽는 기술 선택}
세로축 $K$는 자본, 가로축 $L$은 노동이다. 같은 등량곡선 위에서는 산출이 같지만 $P_2$는 $P_3$보다 자본을 더, 노동을 덜 쓴다. 비용 최소화 조건은 $MPL/MPK=w/r$이다. 임금 $w$가 상대적으로 높으면 노동 절약적 $P_2$가, 자본 비용 $r$가 높고 노동이 풍부하면 $P_3$가 유리하다. 그림의 위쪽 등량곡선은 더 높은 산출을 나타내므로 \textbf{투입 구성의 이동과 기술·산출의 이동}을 나누어 읽어야 한다.\par

\paragraph{등량곡선의 접점을 도출하기}

정해진 산출 $\bar Y$를 가장 싸게 만드는 문제는 $\min_{L,K} wL+rK$이며 제약은 $F(L,K)=\bar Y$다. 생산량을 그대로 두고 노동을 한 단위 늘리면 자본은 $MPL/MPK$ 단위 줄일 수 있다. 이 교환의 비용 변화가 0이어야 하므로 $w-r(MPL/MPK)=0$, 곧 $MPL/MPK=w/r$다. 임금이 오르면 왼쪽의 노동 한계생산이 더 높은, 즉 노동을 아끼는 투입 조합으로 이동한다. 다만 실제 선택에는 도구의 도입 비용과 기술 학습도 들어간다.\par

\subsection{슬라이드 91 · 내권화의 뜻과 논쟁}
Involution은 더 많은 노동을 같은 땅에 넣어 총산출을 올리되 노동자당 산출이나 한계수익의 개선이 작을 수 있는 현상을 가리킨다. 예를 들어 $Y=A T^\alpha L^{1-\alpha}$에서 $T$가 고정이면 $L$ 증가로 $Y$는 늘지만 $Y/L$은 줄어든다. 실제 중국 농업이 어느 정도 이 경로를 따랐는지는 지역과 시대의 자료로 검증해야 한다. 모든 노동 집약화를 비효율이라고 부를 수는 없다.\par

\subsection{슬라이드 92 · 낫과 호미의 비교}
도구의 모양은 작물, 토양, 밭의 크기와 노동·자본의 상대가격에 적응한다. 더 큰 도구가 언제나 “선진적”인 것은 아니다. 작은 경지와 집약적 재배에서는 섬세한 노동 도구가 토지당 산출을 높일 수 있다. 경제적 비교에는 도구의 구매비용과 노동시간, 수확량, 다른 작물에 대한 적합성을 함께 넣어야 한다.\par

\subsection{슬라이드 93 · 18세기 중국의 내부 이주}
지도는 산시·후베이 등 여러 출발지에서 쓰촨과 남서 지역으로 향하는 복수의 화살표를 그린다. 전쟁 뒤의 재정착, 인구 압력, 개간 가능한 토지와 국가의 정책이 경로를 바꾸었다. 하나의 화살표를 실제 이동 인원의 정확한 크기로 읽지 말고, 출발·도착지의 상대적 토지 희소성이 왜 달랐는지 물어야 한다.\par

\subsection{슬라이드 94 · 청대의 개발 지역 구분}
청대 이동 지도는 이미 개발된 동부, 개발 중인 내륙, 상대적으로 개발되지 않은 지역을 나누고 이동 방향을 표시한다. 범례의 “개발”은 해당 지도의 농업 정착과 행정 기준이지 그 지역에 사람이 없었다는 뜻이 아니다. 새로운 경지 개간은 단기적으로 토지/노동 비율을 높이지만, 시간이 지나 인구가 늘면 이 이익이 줄 수 있다.\par

\subsection{슬라이드 95 · 총인구와 지역 비중}
위 그래프는 중국 전체 인구의 장기 추정, 아래 그래프는 중원과 강남의 비중 변화를 나타낸다. 총인구가 늘어도 중심 지역의 비중은 이동과 개발로 달라질 수 있다. 두 그래프의 세로축 단위가 각각 백만 명과 퍼센트라는 점을 확인하자. 비중 하락이 그 지역의 절대 인구 감소를 뜻하지는 않는다.\par

\subsection{슬라이드 96 · 초기 재배 지역의 확산}
지도는 오래된 벼 재배 흔적의 시기를 지역별 음영으로 표시한다. 가장 이른 흔적의 위치와 이후 경작지가 확장된 방향을 구분해야 한다. 농업 기술의 전파는 사람의 이동, 작물의 이동, 현지 집단의 학습 중 어느 경로로도 일어날 수 있다. 현재 국경선은 선사시대 확산의 원인이 아니다.\par

\subsection{슬라이드 97 · 현대의 코끼리 이미지}
이 장면은 이동이 인간 사회만의 현상이 아니며 토지 이용·도시화·환경 변화가 서식지와 충돌할 수 있음을 환기한다. 사진 속 특정 사건 하나에서 장기 기후 원인을 단정할 수는 없다. 인구경제사에서 중요한 연결은 \textbf{사람의 생산과 정착이 다른 생물의 이동 공간을 어떻게 바꾸는가}라는 환경사적 질문이다.\par

\subsection{읽기 · 인구와 경제: 기본 개념 검토}
\subsubsection{인구는 경제에 어떻게 들어오는가}

인구는 노동·소비·저축을 하는 사람의 수이면서 1인당 산출을 계산하는 분모다. 총생산 $Y$가 증가해도 인구 $N$이 더 빨리 증가하면 $Y/N$은 낮아진다. 반대로 총인구가 정체해도 노동 가능한 연령대의 비중이 달라지면 같은 $N$에서 일하는 사람 수가 바뀐다. 따라서 \textbf{인구 규모, 성장률, 연령 구조}를 서로 다른 변수로 다뤄야 한다. \href{https://youtu.be/luHh0S9SlH4}{영상 보기}.\par

\subsubsection{생산가능인구와 부양비}

관례적으로 15\~{}64세를 생산가능인구 $W$, 0\~{}14세를 유소년 $C$, 65세 이상을 고령인구 $O$로 분류한다. 총인구는 $N=C+W+O$다. 유소년부양비는 $C/W$, 노년부양비는 $O/W$, 총부양비는 $(C+O)/W$다. 분모가 \textbf{전체 인구가 아닌 생산가능인구}라는 점이 핵심이다.\par

예를 들어 $C=20$, $W=60$, $O=20$명인 사회의 총부양비는 $40/60=0.667$이다. 20년 뒤 $C=10$, $W=65$, $O=25$라면 총부양비는 $35/65\approx0.538$로 낮아진다. 고령인구는 늘었어도 유소년 감소와 생산가능인구 증가가 더 크기 때문이다. “고령화”와 “총부양 부담 증가”가 같은 시점에 나타난다는 보장은 없다.\par

\subsubsection{인구배당을 항등식으로 보기}

노동자 한 사람당 산출을 $Y/E$, 고용자 수를 $E$, 생산가능인구를 $W$, 전체 인구를 $N$으로 두면 다음 항등식이 성립한다.\par

\[
\frac{Y}{N}=\frac{Y}{E}\cdot\frac{E}{W}\cdot\frac{W}{N}.
\]

1인당 산출은 노동자당 생산성, 생산가능인구 중 고용 비율, 생산가능인구의 전체 인구 비중의 곱이다. $W/N$이 상승하면 다른 두 요소가 같을 때 1인당 산출이 올라갈 여지가 있다. 이것이 인구배당의 회계적 직관이다. 실제 성장에는 일자리, 여성의 노동 참여, 교육과 자본 축적이 필요하다. 생산가능인구 비중 상승만으로 생산성이 자동 상승하는 것은 아니다.\par

\subsubsection{한국의 장기 질문}

한국의 빠른 성장과 생산가능인구 비중의 상승은 함께 나타난 시기가 있다. 두 현상을 같은 그래프에 놓는 것은 출발점이고, 인과관계를 결론 내리려면 교육·수출·투자·기술과 정책을 함께 봐야 한다. 출산율 하락은 처음에는 유소년부양비를 낮출 수 있고, 더 시간이 지나면 일할 연령대의 크기를 줄이고 노년부양비를 높인다. 같은 출산 변화의 경제 효과가 \textbf{세대에 따라 부호를 바꿀 수 있음}을 기억하자.\par

\subsubsection{인구 자료의 범위}

조선의 호적, 일제하 조사, 전후 인구조사는 조사 목적과 누락 방식이 다르다. 호적의 등록 인구가 실제 거주 인구와 같다고 가정하면 안 된다. 장기 비교에서 정의의 변화가 진짜 인구 변화처럼 보이지 않도록 자료의 작성·등록 규칙을 확인해야 한다. \href{https://kosis.kr/}{우리나라 인구통계 설명}과 \href{https://population.un.org/wpp/}{UN World Population Prospects}는 현대 지표의 정의와 비교에 도움이 된다.\par

\subsubsection{확인할 질문}

1. 총인구가 증가해도 1인당 GDP가 떨어질 수 있는 조건은 무엇인가?\par
2. 생산가능인구의 비중 증가가 실제 소득 증가로 이어지려면 무엇이 필요한가?\par
3. 유소년부양비와 노년부양비가 서로 다른 방향으로 움직일 수 있는가?\par


\subsection{읽기 · 인구변천}
\subsubsection{다산다사에서 소산소사로}

인구변천은 출생률과 사망률이 모두 높은 체제에서 둘 다 낮은 체제로 이동하는 전형적인 패턴이다. 중요한 것은 두 율이 \textbf{같은 시점에 하락하지 않는다}는 점이다. 많은 사례에서 사망률, 특히 아동 사망률이 먼저 낮아지고 출생률은 뒤따른다. 그 시차 동안 인구의 자연증가율이 높아진다. \href{https://youtu.be/J4jeB1QHjOY}{영상 보기}, \href{https://ourworldindata.org/demographic-transition}{여러 나라의 장기 패턴}.\par

\subsubsection{네 단계의 회계}

1단계에서는 높은 출생률 $b$와 높은 사망률 $d$가 비슷하여 $b-d$가 작다. 2단계에서는 위생·영양·의료 개선으로 $d$가 내려가지만 $b$는 높게 남아 $b-d$가 커진다. 3단계에서는 출산 선택이 바뀌어 $b$가 내려가므로 인구는 아직 증가해도 \textbf{증가 속도는 둔화}한다. 4단계에서는 둘 다 낮아 자연증가율이 작다. 출생률이 사망률보다 낮으면 자연감소도 가능하다.\par

\[
g_N^{\mbox{자연}}=\frac{B-D}{\bar N}=\frac{b-d}{1000}
\]

여기서 $b,d$를 천 명당 인원(‰)으로 썼다. $b=30$, $d=10$이면 자연증가율은 연 2\%다. 이주가 있다면 $g_N=(B-D+I-E)/\bar N$을 써야 한다. 출생과 사망의 차이를 총인구 성장률로 곧바로 읽으면 국제·국내 이동의 영향을 놓친다.\par

\subsubsection{그래프의 면적과 기울기}

출생률 곡선이 사망률 곡선보다 얼마나 위에 있는지가 그 해의 자연증가율이다. 두 곡선 사이의 *면적*은 여러 해에 걸친 증가율의 누적 직관을 주지만, 정확한 인구 배수를 얻으려면 해마다 성장률을 복리로 곱해야 한다. $g_t$가 작을 때 $\ln(N_T/N_0)\approx\sum_t g_t$다. 3단계에서 곡선 간격이 줄면 인구 수준이 즉시 감소하는 것이 아니라 증가율이 낮아진다. 이것이 흔한 그래프 해석 오류다.\par

\subsubsection{연령 구조의 파도}

2단계에서 아동이 더 많이 살아남으면 유소년 비중과 유소년부양비가 높아질 수 있다. 시간이 지나 그 집단이 노동 연령에 들어가면 생산가능인구 비중이 높아진다. 출생률이 계속 낮으면 더 뒤에는 고령 인구 비중이 상승한다. 사망률과 출생률의 변동이 \textbf{시간 지연을 두고 연령 구조를 바꾸는 것}이다. 학교 수요, 노동 공급, 연금과 의료 수요도 순서대로 움직일 수 있다.\par

\subsubsection{왜 국가별 속도가 다른가}

스웨덴의 사례처럼 긴 세기에 걸쳐 사망과 출생이 천천히 하락한 곳과, 보건기술이 이미 존재하던 20세기에 사망률이 빠르게 떨어진 곳은 2단계의 폭이 다르다. 기술의 국제 확산, 교육·도시화의 속도와 가족계획에 대한 접근이 시차를 바꾼다. “선진국”과 “개발도상국”이라는 큰 분류 안에도 지역·계층의 차이가 있으므로 각 나라의 실제 시계열을 확인해야 한다.\par

\subsubsection{4단계 그림의 한계}

이 도식은 평균적 순서를 보여 주지만 모든 국가에 동일한 연도와 정책을 지정하지 않는다. 전쟁·전염병과 이주가 곡선을 흔들 수 있고, 4단계 이후 낮은 출생률이 오래 지속될 수도 있다. 5단계 모형을 덧붙이는 경우도 있으나 그것이 인구 감소의 필연성을 뜻하지는 않는다. \href{https://population.un.org/wpp/}{UN 인구 전망}에서 출생·사망·이동 가정을 따로 확인해야 한다.\par

\subsubsection{확인할 질문}

1. 사망률과 출생률이 동시에 같은 폭으로 떨어지면 왜 인구 폭증이 일어나지 않는가?\par
2. 3단계의 인구가 늘면서 증가율은 낮아질 수 있는가?\par
3. 낮은 출산의 경제 효과가 처음과 수십 년 뒤에 다른 이유는 무엇인가?\par


\subsection{읽기 · 사망률 변천}
\subsubsection{사망률을 먼저 보는 이유}

많은 인구변천 사례에서 사망률의 하락이 출생률 하락보다 앞선다. 살아남는 아이의 수가 먼저 늘면 가족의 출산 목표와 사회의 연령 구조가 시간이 지나 바뀐다. 이 흐름을 이해하려면 “사망률” 하나 안에 어떤 분모와 연령이 들어 있는지 구분해야 한다. \href{https://youtu.be/953x7Le4PnY}{영상 보기}.\par

\subsubsection{조사망률과 연령별 사망률}

조사망률은 기간 중 사망자 수 $D$를 연앙인구 $N_m$으로 나눈 뒤 천 명당 값으로 표시한다.\par

\[
CDR=1000\frac{D}{N_m}.
\]

연앙인구는 보통 7월 1일의 인구처럼 연중 가운데의 규모를 사용한다. 조율은 계산하기 쉽지만 고령자가 많은 사회에서 높아질 수 있다. 연령집단 $a$의 사망률 $m_a=D_a/N_a$와 그 집단 비중 $s_a=N_a/N$를 두면 전체 사망률은 $\sum_a s_a m_a$다. 따라서 $m_a$가 모두 낮아져도 고령층의 $s_a$가 충분히 커지면 전체 조사망률은 상승할 수 있다.\par

예를 들어 젊은층 사망률이 1‰, 고령층이 50‰이고 인구의 10\%가 고령이면 전체는 $0.9(1)+0.1(50)=5.9$‰다. 다른 사회의 연령별 사망률이 각각 0.5‰, 30‰로 더 낮아도 고령층이 30\%라면 $0.7(0.5)+0.3(30)=9.35$‰다. 두 번째 사회의 높은 조율을 의료가 뒤떨어진다는 증거로 읽으면 오류다. \textbf{연령 표준화}는 두 사회에 같은 연령 비중을 적용해 이 효과를 제거한다.\par

\subsubsection{기대여명과 생명표}

출생 시 기대수명은 “그해 태어난 아이가 반드시 그 나이에 사망한다”는 예언이 아니다. 특정 기간의 연령별 사망률이 평생 유지된다고 가정한 가상 집단의 평균 생존 연수다. 생명표에서 $q_x$를 나이 $x$에서 다음 나이 전에 사망할 확률, $l_x$를 그 나이까지 살아남은 사람 수라 두면 $l_{x+1}=l_x(1-q_x)$다. 각 연령 구간에서 산 사람들의 총 생존 연수 $L_x$를 더하면 출생 시 기대수명은 $e_0=\sum_x L_x/l_0$가 된다.\par

어린 나이에 많은 사람이 사망하면 $e_0$가 크게 낮아진다. 그러나 \textbf{출생 시 기대수명 37세가 대부분의 성인이 37세에 죽었다는 뜻은 아니다.} 30세까지 살아남은 사람의 조건부 기대여명은 별도로 계산해야 한다. 성인 생존도 오늘보다 어려웠는지는 조건부 생존율 자료로 확인할 수 있다.\par

\subsubsection{소득, 위생, 교육의 연결}

소득 상승은 식품과 주거를 개선하고 위생·교육에 투자할 자원을 만든다. 깨끗한 식수와 하수도, 예방접종과 치료는 특히 감염성 질환을 줄였다. 하지만 국가 사이 소득과 기대수명의 양의 상관관계만으로 소득의 인과효과를 바로 추정할 수 없다. 공공의료가 소득과 생존을 함께 바꿀 수도 있다. 시대·지역별로 어떤 기술과 정책이 먼저 도입되었는지 비교해야 한다.\par

여성 교육은 건강 정보의 접근, 의료 이용과 가구 내 의사결정에 영향을 줄 수 있다. 그러나 교육을 많이 받은 가구가 원래 부유하고 위생시설에 접근하기 쉬웠을 가능성도 있다. 국가나 종교권의 평균 차이를 문화적 본질로 설명하기보다 교육·소득·공공서비스와 법적 권리를 분리해 검사해야 한다.\par

\subsubsection{감염성·비감염성 질환의 그래프}

20세기 미국에서 감염성 질환 사망이 크게 줄고 비감염성 질환의 \textbf{조사망 수치}가 덜 변한 그래프를 읽을 때, 고령자가 많아진 인구 구성을 고려해야 한다. 사람들이 더 오래 살면 암·심혈관 질환이 나타날 나이까지 생존하는 사람도 늘어난다. 비감염성 질환의 연령별 사망률이 낮아져도 총인구 천 명당 사망자는 다른 모습을 보일 수 있다. 발병률의 변화와 생존자의 연령 구성, 진단 기준의 변화는 별도 설명이다.\par

\subsubsection{역사 자료의 비교}

조선 과거 급제자의 수명은 성인 남성 중 선택된 엘리트 표본이고 현대 전체 남성의 생명표와 자료 구조가 다르다. 비교할 때에는 30세까지 생존한 집단의 60세 도달 확률처럼 조건부 확률을 맞추어야 한다. \href{https://www.mortality.org/}{Human Mortality Database}의 생명표와 \href{https://ourworldindata.org/child-mortality-global-overview}{Our World in Data의 아동 사망 자료}는 지표의 정의를 확인하는 데 도움이 된다.\par

\subsubsection{확인할 질문}

1. 어느 나라의 조사망률이 더 높아도 모든 연령의 사망 위험은 더 낮을 수 있는가?\par
2. 출생 시 기대수명이 낮았다는 사실로 성인의 평균 사망 나이를 바로 정할 수 없는 이유는 무엇인가?\par
3. 소득·교육과 기대수명의 상관관계에서 인과를 확인하려면 무엇을 비교해야 하는가?\par


\subsection{읽기 · 출생률 변천}
\subsubsection{출생을 어떤 지표로 볼 것인가}

조출생률 $CBR=1000B/N_m$은 총인구 천 명당 출생자 수다. 가임연령의 여성 비중이 작으면 여성 한 사람의 출산 행동이 바뀌지 않아도 조출생률이 낮을 수 있다. 합계출산율(TFR)은 특정 연도의 연령별 출산율을 한 여성이 가임기간 내내 경험한다고 가정할 때의 출생아 수다. 1세 연령별 자료에서는 $TFR=\sum_{a=15}^{49} f_a$, 5세 구간을 쓰면 각 구간의 연율을 합쳐 5를 곱한다. \href{https://youtu.be/d1JINF3af34}{영상 보기}.\par

TFR은 \textbf{그해 태어난 여성의 미래 출산을 예언한 수치가 아니다.} 실제 한 출생 코호트의 완결 출산율과 다르며, 출산 시기를 늦추는 변화가 빠를 때에는 한동안 기간 TFR이 낮아질 수 있다. 과거 기혼 여성의 평생 출산 자녀 수와 오늘 전체 여성의 기간 TFR을 그대로 비교하면 혼인율과 지표의 정의 차이가 섞인다.\par

\subsubsection{사망률 하락과 생존 자녀 목표}

부모가 성인이 될 자녀의 기대 수를 $k$명 원하고 각 아이의 생존 확률이 $s$라면, 아주 단순한 평균 모형에서 필요한 출생 수 $n$은 $sn=k$, 즉 $n=k/s$다. $k=3$, $s=0.5$이면 평균적으로 6명을 출산한다는 계산이 나온다. $s=0.9$이면 약 3.33명이다. 실제 가구는 자녀의 생존을 확인한 뒤 다시 결정하고, 생존 결과의 불확실성·양육비와 피임 접근을 고려한다. 이 식은 \textbf{생존율이 높아질 때 같은 목표의 출생 수가 줄어드는 방향}을 보여 주는 직관이다.\par

사망률은 먼저 낮아지고 출생 행동은 늦게 바뀔 수 있다. 부모가 새로운 생존 환경을 인지하고 기대를 조정하는 데 시간이 걸리기 때문이다. 그동안 살아남는 아동의 수가 늘어 인구가 빠르게 증가한다.\par

\subsubsection{자녀의 수와 인적자본의 상충관계}

각 자녀에게 기본 양육비 $c_0$와 교육 수준 $q$에 비례한 비용 $c_1q$가 든다고 하자. 부모 소득 $I$와 다른 소비 $C$의 예산제약은 다음과 같다.\par

\[
C+n(c_0+c_1q)\leq I.
\]

자녀가 $n$명일 때 교육 수준을 한 단위 높이는 총비용은 $nc_1$이다. 자녀 수가 많을수록 같은 교육 개선의 비용이 커진다. 반대로 교육을 더 많이 시키면 자녀 한 명의 비용 $c_0+c_1q$가 올라가 추가 출산의 비용도 커진다. 이 교차 비용이 quantity–quality trade-off의 기초다. 교육의 노동시장 보상이 커지면 $q$의 편익이 높아질 수 있다. 하지만 소득이 늘면 자녀 수에 대한 소득효과도 가능하므로 \textbf{소득 상승만으로 출산 감소를 논리적으로 필연화할 수는 없다.} 가격·기회비용·규범과 정책이 함께 바뀐다.\par

\subsubsection{여성 교육과 시간의 기회비용}

출산·양육에 필요한 시간을 $h$시간, 시장에서 벌 수 있는 시간당 임금을 $w_f$라고 하면 단순한 시간 기회비용은 $w_fh$다. 교육이 $w_f$를 높이면 시장 일을 포기하는 비용이 커진다. 그러나 교육은 건강 지식과 가구 협상력도 바꾸고, 보육·부모휴가 정책은 양육의 시간 비용을 낮출 수 있다. 여성의 노동 참여가 늘었다는 사실만으로 모든 사회의 출산율이 같은 방향으로 움직인다고 단정하면 안 된다.\par

\subsubsection{혼인과 출산 시기의 변화}

첫 출산 연령이 늦어지는 기간에는 각 연령의 출생 건수가 일시적으로 줄어 기간 TFR이 크게 낮아질 수 있다. 나중에 출산이 회복되는 정도에 따라 완결 출산율은 다르게 움직인다. 역사 비교에서는 혼인율, 혼인 연령, 비혼 출산의 규칙, 피임 기술을 함께 살펴야 한다. “혼인 지연”을 단순한 선호 변화로만 볼 수 없고 주거·고용의 제약도 중요하다.\par

\subsubsection{인구변천으로 돌아가기}

사망률 하락은 당장의 아동 생존을 높인다. 출생률 하락이 뒤따르면 자연증가율은 줄지만 인구가 즉시 감소하는 것은 아니다. 이미 태어난 큰 코호트가 노동 연령과 출산 연령에 들어가면 \textbf{인구 모멘텀}이 남는다. 이후 그 코호트가 고령층으로 이동하면 부양비의 구조가 바뀐다. \href{https://ourworldindata.org/demographic-transition}{인구변천 장기 비교}와 \href{https://population.un.org/wpp/}{UN World Population Prospects}는 출산·사망·이주를 나눠 볼 수 있는 보조 자료다.\par

\subsubsection{확인할 질문}

1. 조출생률은 낮아졌는데 합계출산율은 그대로일 수 있는가?\par
2. 자녀 생존율이 높아진 직후와 수십 년 뒤의 출산 행동은 왜 다를 수 있는가?\par
3. 교육 투자 확대가 출산 수를 줄이는 예산 기제와 그 결론의 한계는 무엇인가?\par


\section{농업: 토지와 노동의 제도}
농업 생산력과 토지의 권리 구조를 함께 보면 장원, 흑사병, 농민 해방, 인클로저와 공업화가 한 인과 사슬로 연결된다.\par

\subsection{슬라이드 01 · 농업을 읽는 세 축}
농업사는 작물과 농법의 역사만이 아니다. 같은 토지와 도구로 누가 일하고 누가 수확을 가져가는지, 땅을 어떻게 상속·거래하는지, 식량이 도시와 공업을 얼마나 부양하는지가 함께 변한다. 토지 생산성(면적당 산출)과 노동 생산성(노동시간당 산출)을 분리해 두면 뒤의 장원·흑사병·인클로저가 연결된다.\par

\subsection{슬라이드 02 · 빙하기의 환경 제약}
기후와 식생은 초기 인간이 이용할 수 있는 야생 동식물과 이동 경로를 제한했다. 그러나 환경이 같은 곳에서 모두 같은 농경 제도가 생긴 것은 아니다. 기술·인구 밀도·지식의 전파와 사회적 협력이 중요하다. 지도에서 빙하와 해안선의 당시 모습을 현재 국경과 혼동하지 말고, 어떤 자원 접근성이 달라졌는지 보자.\par

\subsection{슬라이드 03 · 채취에서 재배로}
수렵·채집에서 농경과 목축으로의 전환은 식량을 계획적으로 생산·저장하게 했지만 노동시간과 질병 노출도 바꾸었다. “혁명”은 단일한 연도의 급변보다 긴 적응 과정을 가리킨다. 생산의 잉여가 생기면 전문직과 국가가 가능해지지만, 그 잉여가 균등하게 분배된다는 보장은 없다.\par

\subsection{슬라이드 04 · 수렵·채집의 쇠퇴}
인구 증가와 경지 확대는 기존 사냥터와 채집 공간을 줄였을 수 있다. 농경이 언제나 개인에게 더 나은 선택이어서만 확산된 것은 아니다. 한 지역이 정착 농업으로 바뀌면 주변 집단의 이동과 자원 이용도 제약된다. 기술 채택의 편익을 개인의 영양·노동과 사회 전체의 인구 부양력으로 나누어 비교하자.\par

\subsection{슬라이드 05 · 구석기 식생활의 비교}
현대의 “구석기식”을 과거의 실제 평균 식단으로 생각하면 안 된다. 수렵·채집의 식품은 계절·지역에 따라 달랐고, 생존한 유골이나 유적도 표본 선택을 겪었다. 농경 이전과 이후의 생활 수준을 비교하려면 열량뿐 아니라 영양 다양성, 감염, 작업량과 사망 위험을 함께 보아야 한다.\par

\subsection{슬라이드 06 · 비옥한 초승달 지대}
서남아시아의 야생 곡물과 가축화 가능한 동물은 초기 식량 생산의 조건이었다. 하지만 작물이 존재한다는 것과 사람이 그것을 재배하기 시작한 이유는 다르다. 저장 가능성, 계절 위험, 인구 압력, 지식의 축적이 선택의 편익을 바꾸었다. 농업의 한 발원지를 모든 지역의 유일한 기원으로 읽지 말자.\par

\subsection{슬라이드 07 · 여러 농경 중심지}
작물의 가축화와 재배는 여러 지역에서 서로 다른 시기에 나타났다. 지도에는 작물의 원산지, 최초 재배 흔적, 후대 확산 경로가 섞일 수 있으니 범례를 확인하자. 비슷한 농경이 여러 곳에서 생겼다는 사실은 특정 문명의 선천적 우월성보다 지역별 생태와 학습의 가능성을 보게 한다.\par

\subsection{슬라이드 08 · 농경의 생활 수준 역설}
농경은 면적당 식량을 늘려 더 많은 사람을 부양할 수 있지만, 인구가 함께 증가하면 1인당 식량은 정체할 수 있다. 정착과 가축 사육은 감염병 노출도 높였다. 평균 신장이 줄었다는 자료는 영양·질병의 변화와 일치할 수 있지만, 지역과 유골 표본의 차이를 통제해야 한다. 총산출의 증가와 개인의 후생을 구별하는 맬서스적 사례다.\par

\subsection{슬라이드 09 · 유골 신장 자료}
뼈에서 추정한 평균 키는 성장기의 영양과 질병 환경을 반영한다. 그러나 유골이 남을 확률, 발굴된 묘지의 사회계층, 남녀·연령 구성의 변화가 시계열을 왜곡할 수 있다. 키가 낮아진 시기를 곧바로 소득 하락이라고 결론 내리지 말고, 치아·병변·매장 방식과 다른 자료를 교차해야 한다.\par

\subsection{슬라이드 10 · 골격 건강의 여러 차원}
신장 외에 치아 손상, 감염 흔적, 빈혈 징후 등은 삶의 질에 관한 다른 정보를 준다. 한 지표가 개선되고 다른 지표가 악화할 수 있다. 농경의 안정적 식량과 단조로운 식단, 인구 밀집의 감염을 함께 고려해야 한다. 그래프의 표본 수와 신뢰구간을 확인하면 작은 차이를 과도하게 해석하지 않을 수 있다.\par

\subsection{슬라이드 11 · 토지와 생산의 세 요소}
토지 $T$, 노동 $L$, 자본 $K$는 서로 보완하거나 대체할 수 있다. 초기 사회의 자본은 기계뿐 아니라 저장 시설·종자·가축과 도구를 포함한다. 공동 이용지는 “누구의 것도 아닌 땅”과 다르다. 이용·배제·상속의 규칙이 있는 공유 제도일 수 있다. 소유권의 명칭보다 실제 권리 묶음을 살펴야 한다.\par

\subsection{슬라이드 12 · 라티푼디움에서 콜로나투스로}
로마의 대토지 경영은 노동을 조직하는 방식에 따라 달라졌다. 노예 노동에 기대던 라티푼디움과 소작·종속 경작자를 쓰는 콜로나투스는 모두 대지주의 토지 통제라는 공통점이 있지만 노동자의 법적 지위와 경영 유인이 다르다. 제도의 이행을 한 번의 교체로 보지 말고 지역별 혼재를 살펴야 한다.\par

\subsection{슬라이드 13 · 라티푼디움의 비용 구조}
노예제 대농장은 넓은 토지와 감독 조직을 결합한다. 노예의 취득·유지·감시 비용과 도주 위험이 있으므로 임금이 0인 노동이 아니다. 생산량을 비교할 때에는 토지 품질과 작물, 숙련과 폭력의 비용을 포함해야 한다. 고대의 대농장 이윤이 자유 농민의 삶을 개선했음을 의미하지 않는다.\par

\subsection{슬라이드 14 · 콜로누스의 종속}
콜로누스는 토지를 경작하면서 지대나 세금을 부담했고 이동의 자유가 제한될 수 있었다. 노예와 동일하지 않지만 완전한 자유 계약의 소작농도 아니다. 노동자에게 일정한 경영 재량을 주면 감독비용은 줄 수 있으나, 토지와 법적 권리를 지주가 통제하면 잉여의 분배는 여전히 불평등하다.\par

\subsection{슬라이드 15 · 봉건제라는 말의 범위}
“봉건제”는 서유럽의 주군·봉신 관계를 가리키기도, 영주와 농민의 생산관계를 가리키기도, 더 넓은 전근대적 지배를 뜻하기도 한다. 같은 단어를 쓸 때 어느 층위를 말하는지 먼저 정해야 한다. 정치적 충성의 사슬과 장원의 노동 의무가 언제나 동일한 공간에 일치하지는 않는다.\par

\subsection{슬라이드 16 · 상부 구조의 교환}
봉주는 군사·정치적 봉사를 받고 봉신에게 토지 수익이나 보호를 제공했다. 인적 주종 관계와 물권 관계가 결합하되 각각의 권리·의무는 다르다. 현금 국가가 모든 군인을 직접 고용하는 방식과 달리, 권력과 수입원을 분산해 동맹을 유지하는 체계였다. 충성의 계약이 곧 농민에게 자유로운 선택을 주는 것은 아니다.\par

\subsection{슬라이드 17 · 중층적 수봉}
한 봉신이 다른 사람의 봉주가 될 수 있으므로 권리와 의무가 여러 층으로 이어진다. 그림의 피라미드는 이해를 돕지만 실제 충성 관계는 겹치고 충돌할 수 있었다. 위계가 길어질수록 중앙의 명령 집행과 지방의 자율성이 달라진다. 토지에서 나오는 수입이 군사·사법 권력과 연결되는 구조에 주목하자.\par

\subsection{슬라이드 18 · 장원이라는 하부 구조}
장원제는 생산 활동에 직접 참여하는 영주와 농민의 관계를 뜻한다. 고전적 장원에서는 영주 직영지를 농민의 부역으로 경작하고, 농민은 자기 보유지도 경작한다. 순수 또는 지대 장원에서는 영주가 직접 생산을 조직하기보다 지대를 수취한다. 이 차이는 노동의 상대가격과 시장 발달에 따라 중요해진다.\par

\subsection{슬라이드 19 · 용어의 포함 관계}
중세 유럽의 봉건제, 장원제, 농노제는 서로 겹치지만 같은 말이 아니다. 정치적 봉건제는 봉주·봉신, 장원제는 영주·농민, 농노제는 농민의 법적·인신적 제약에 초점을 둔다. 중국의 “봉건”이나 현대의 “봉건적”이라는 비유를 유럽 장원의 세부 제도와 바로 등치하지 않도록 주의하자.\par

\subsection{슬라이드 20 · 농노는 어떤 사람인가}
농노는 노예와 달리 보유지에서 가구 단위로 경영할 수 있었지만 영주에 대한 부역·지대와 이동·혼인의 제약을 졌다. 자유 농민과의 경계도 시대·지역에 따라 달랐다. 경영의 재량과 인신의 자유를 하나의 축으로 합치면 혼동된다. \textbf{누가 땅을 경작하고, 누가 잉여를 요구하며, 떠날 수 있는가}를 나누어 묻자.\par

\subsection{슬라이드 21 · 농노 형성의 여러 경로}
노예가 토착 소작인으로 바뀐 경로, 자유 농민이 부채·전쟁 때문에 보호를 구하다 종속된 경로가 모두 가능하다. 한 가지 기원설로 모든 지역을 설명하기 어렵다. 자유를 일부 양도하고 보호를 받는 거래처럼 보이는 경우에도 당시의 강제와 선택 가능한 대안을 따져야 한다.\par

\subsection{슬라이드 22 · 영주권의 두 성격}
영주는 토지에서 지대를 받을 권리와 농민에게 명령·재판을 행사할 권력을 함께 가질 수 있었다. 전자는 경제적 소유와 가깝고 후자는 국가의 공권력이 분산된 모습이다. 현대의 토지 임대차에서는 지주가 임차인을 형사 처벌할 권리까지 갖지 않는다. 이 차이가 장원제를 단순한 소작제와 구분한다.\par

\subsection{슬라이드 23 · 경제외적 강제}
시장에서 지대를 내지 않으면 계약을 끝내는 관계와, 영주가 사법·신분 권한으로 부역을 강제하는 관계는 다르다. 경제외적 강제는 가격 교섭만으로 노동 의무가 정해지지 않음을 뜻한다. 그러나 강제가 있다고 경제적 유인이 사라지는 것은 아니다. 농민은 부역의 감시, 도주 기회, 자기 땅의 수확을 비교하며 대응했다.\par

\subsection{슬라이드 24 · 이동·혼인·상속의 비용}
장원 밖 결혼에 대한 납부금과 토지 상속 비용은 농민 가구의 장기 결정을 바꾸었다. 인신적 지배를 추상적인 “부자유” 한 단어로 끝내지 말고, 구체적으로 어떤 이동과 가족 결정의 비용이 올랐는지 보자. 영주는 노동력의 유출을 막고 수입을 확보하려 했고 농민은 더 나은 기회를 얻으려 했다.\par

\subsection{슬라이드 25 · 중층적 토지 소유}
같은 땅에 영주의 상급 권리와 농민의 세습 경작권이 겹칠 수 있다. 오늘의 배타적 소유권 개념으로는 이를 한 사람의 절대 소유로 정리하기 어렵다. 직영지와 보유지를 구분하고, 농민이 작물·상속·양도에 대해 실제로 어떤 권리를 가졌는지 확인해야 토지 시장의 형성과정을 이해할 수 있다.\par

\subsection{슬라이드 26 · 부역의 현장}
그림 속 감독자는 노동을 징발했을 때 생기는 유인 문제를 보여 준다. 농민이 자기 보유지에서는 수확을 더 가져가지만 영주 직영지에서는 한계노력의 보상이 작으면 감독이 필요하다. 감독 비용과 계절별 노동 수요가 커질수록 영주가 부역 대신 지대를 받고 임금 노동자를 쓰는 선택도 매력적일 수 있다.\par

\subsection{슬라이드 27 · 공조와 현물 납부}
농민이 생산물 일부를 납부하는 제도는 수확의 위험을 어떻게 나누는지에 따라 다르다. 고정량 현물 지대는 흉작 때 농민에게 큰 위험을 남기고, 수확 비율 지대는 영주가 위험 일부를 공유한다. 다만 비율 지대에서는 농민이 추가 노력의 성과를 모두 받지 못한다. 지대의 이름보다 계산 방식과 집행을 확인해야 한다.\par

\subsection{슬라이드 28 · 장원 도식의 공간}
그림에서 영주의 직영지, 농민 보유지, 공동 목초와 마을을 구별해 보자. 장원은 하나의 담장 안에 모인 농장이 아니라 흩어진 지조와 여러 사용권의 조합이었다. 생산은 가구 단위로도 이루어졌지만 농법과 수확 시기는 공동의 규칙에 묶였다. 공간 구조가 거래와 감독의 비용을 결정한다.\par

\subsection{슬라이드 29 · 생산양식 도식의 한계}
노예제·봉건제·자본제를 나란히 놓으면 잉여 이전의 주된 관계를 비교할 수 있다. 그러나 각 시대의 모든 사람이 한 관계에 속한 것은 아니다. 중세에도 자유 농민과 임금 노동자가 있었고 현대에도 강제노동이 존재한다. 그림을 \textbf{분류 도구}로 쓰고, 역사 전체가 한 줄로 진행한다는 증거로 읽지 말자.\par

\subsection{슬라이드 30 · 노예·농노·노동자의 세 축}
표는 인신의 자유, 토지의 소유, 경영의 독립을 별개의 축으로 나눈다. 농노는 경영 독립이 있어도 이동의 자유가 제한될 수 있고, 임금 노동자는 자유롭지만 생산수단을 소유하지 않을 수 있다. 한국의 노비 유형과 유럽 농노를 비교할 때도 법적 지위와 실제 노동 형태가 일치하지 않는 사례를 인정해야 한다.\par

\subsection{슬라이드 31 · 택지·채원과 공동 경지}
농가는 집과 텃밭을 중심으로 생활하고, 멀리 흩어진 경지를 이용했다. 집 주변의 채원은 집약적으로 가꾸기 쉬웠지만 공동 경지는 경작 순서와 가축 방목의 규칙을 따라야 했다. 같은 가구의 토지가 여러 조각인 이유를 비효율 하나로만 판단하지 말자. 토질과 물 접근의 위험을 분산하는 기능도 있었다.\par

\subsection{슬라이드 32 · 촌락 항공사진 읽기}
항공사진에서 긴 경지의 흔적과 도로·주거의 위치를 찾으면 문헌의 장원 구조가 공간으로 보인다. 다만 오늘 남은 형태를 곧바로 중세의 원형으로 간주해서는 안 된다. 수백 년의 인클로저와 도로 건설이 경계를 바꿨을 수 있다. 지형 자료는 토지대장과 함께 읽을 때 더 강한 증거가 된다.\par

\subsection{슬라이드 33 · 총유와 보유지}
가구가 이용하는 택지, 경지, 공동 이용지는 권리의 성격이 다르다. 총유는 개인이 제멋대로 매각할 수 없는 공동 사용과 규칙을 포함한다. 여러 조각을 묶어 한 가구의 생계를 이룬다는 뜻에서 후페·만스 같은 단위가 쓰였다. 면적 단위 하나를 현대의 헥타르로 단순 환산하면 지역별 수확력 차이를 놓친다.\par

\subsection{슬라이드 34 · 모르겐이라는 시간 단위}
Morgen은 원래 일정한 시간에 갈 수 있는 땅의 면적이라는 뜻과 연결된다. 같은 노동시간이라도 토양·쟁기·소의 수가 다르면 실제 면적이 달라진다. 전근대의 도량형은 물리적 면적뿐 아니라 노동과 경작 관행의 정보를 담는다. 지역별 모르겐을 그대로 더하거나 평균내기 전에 정의를 확인해야 한다.\par

\subsection{슬라이드 35 · 개방 경지제의 제약과 이익}
긴 지조는 공동의 쟁기와 소를 이용하기 편하고 토질이 다른 구획을 나누어 가질 수 있다. 반면 한 사람만 파종 시기와 울타리를 바꾸기 어렵다. 제도가 유지된 것은 단순 무지 때문이 아니라 공동 작업과 위험 분산의 이익도 있었기 때문이다. 생산성 비교에서는 경지 통합 이익과 기존 협력의 손실을 함께 본다.\par

\subsection{슬라이드 36 · 지조의 모양}
길고 좁은 지조는 큰 쟁기의 회전 횟수를 줄이고 여러 가구가 좋은 땅과 나쁜 땅을 나누는 데 유리할 수 있다. 그림의 경계가 실제 울타리라는 뜻은 아니다. 작물이 자라고 수확이 끝난 뒤의 방목 규칙이 달랐다. 형태만 보고 개인 소유인지 공동 사용인지 단정하지 말고 권리 문서를 확인해야 한다.\par

\subsection{슬라이드 37 · 쟁기와 토양의 관계}
무거운 토양을 깊게 뒤집는 북유럽의 쟁기와 건조한 지중해 토양의 가벼운 쟁기는 다른 환경에 적응했다. 어느 쟁기가 무조건 더 효율적인지 비교하려면 토양, 강우, 축력 비용과 수확량을 같게 맞춰야 한다. 기술의 “선진성”은 목적 함수와 생산 요소의 상대가격에 달려 있다.\par

\subsection{슬라이드 38 · 공동 축력의 필요}
여러 마리의 소를 묶어 쟁기질하려면 소유자들의 협력이 필요하고, 이웃의 경지 경계와 일정도 맞춰야 한다. 한 가구가 소를 모두 보유하지 못해도 공동 사용으로 깊은 경작이 가능해진다. 협력이 무너지면 개별 농가의 생산성이 떨어질 수 있어 경작 강제와 공동 규칙이 생겼다.\par

\subsection{슬라이드 39 · 정전법과 비교의 한계}
동아시아의 정전법 도식과 유럽의 개방 경지제는 토지의 구획과 공동 의무를 비교하는 재료다. 그러나 같은 사각형 배치가 같은 소유권이나 실제 운영을 뜻하지 않는다. 정전법이 어느 시대에 얼마나 현실적으로 시행되었는지와 중세 유럽 지조의 문서 증거는 자료의 성격부터 다르다. 형태보다 지대·부역의 규칙을 비교하자.\par

\subsection{슬라이드 40 · 왜 개방 경지였는가}
깊은 쟁기질, 소의 공동 이용, 지력과 작물의 순환이 개방 경지의 조정 이익을 만든다. 반면 혼재된 지조는 개별 농가가 새로운 작물만 단독으로 실험하는 비용을 높인다. 제도의 순효과는 시대에 따라 바뀐다. 축력과 공동 방목이 중요할 때는 유리하고, 개별적인 농법 개선의 가치가 커지면 제약이 커질 수 있다.\par

\subsection{슬라이드 41 · 혼재 경지제}
영주의 직영지와 농민의 보유지가 여러 경포에 흩어져 있으면 각 가구가 토질과 수확 위험을 분산한다. 그러나 멀리 떨어진 조각 사이의 이동 시간이 들고 공동 일정이 필요하다. 한 지주가 경지를 한곳에 모을 때 얻는 효율과 여러 농민이 잃는 보험 기능을 동시에 평가해야 한다.\par

\subsection{슬라이드 42 · 북부 프랑스의 토지 지도}
좁고 긴 필지가 반복되는 지도는 혼재 경지의 공간적 흔적이다. 검은 영역과 사선 영역의 소유자·사용권을 범례로 확인하고, 한 가구가 서로 다른 위치의 땅을 얼마나 보유했는지 찾아보자. 지도만으로 수확량을 읽을 수는 없지만, 공동 경작에 왜 조정 규칙이 필요했는지는 보인다.\par

\subsection{슬라이드 43 · 경작 강제의 경제학}
한 사람의 지조가 다른 사람의 지조 사이에 끼어 있으면 독자적으로 울타리, 파종, 방목을 결정하기 어렵다. Flurzwang은 공동의 작물 선택과 일정에 따르도록 만든다. 이 규칙은 무임승차를 막을 수 있지만 혁신의 실험을 늦출 수 있다. 공동 조정의 편익과 개인 선택의 제약이 함께 존재한다.\par

\subsection{슬라이드 44 · 인클로저 이전의 경지}
지도에 울타리가 없다고 토지가 “주인 없는 공유지”였던 것은 아니다. 경작권, 수확 뒤 가축 방목권, 길 통행권이 사람마다 달랐다. 인클로저는 단순히 경계선을 긋는 일이 아니라 이 권리들을 재배분하는 과정이었다. 이후 생산성 변화와 함께 기존 이용자의 손실도 살펴야 한다.\par

\subsection{슬라이드 45 · 삼포제의 회전}
겨울곡물, 봄곡물, 휴한지를 매년 돌리면 두 포제보다 한 해 경작 면적의 비율이 $1/2$에서 $2/3$으로 늘 수 있다. 모든 조건이 같다면 경작 가능한 몫은 $(2/3)/(1/2)=4/3$배다. 실제 식량 증가는 작물별 수확률·노동·비료와 지력 변화에 달려 있어 33\% 증가라고 단정할 수 없다. 휴한은 낭비가 아니라 당시의 지력 회복 방식이었다.\par

\subsection{슬라이드 46 · 전형적 장원의 첫 그림}
집, 공동 경지, 초지와 영주의 영역을 한 그림에서 구분하자. 한 가구는 자기 보유지를 경작하면서도 직영지의 노동과 공동지의 규칙에 묶인다. 농업의 생산조직과 정치·사법 권한이 같은 공간에서 겹친다는 점이 현대의 독립 농장과 다르다.\par

\subsection{슬라이드 47 · 지조와 권리의 혼합}
같은 경포 안에 영주와 농민의 지조가 번갈아 있으면 누가 언제 쟁기를 쓰고 가축을 들일지 공동 결정해야 한다. 물리적 혼합이 사회적 평등을 의미하지는 않는다. 영주는 지대와 재판권을 가졌고 농민은 토지 이용을 위해 협력했다. 그림을 보며 각 화살표가 노동 의무인지 토지 사용인지 나누어 보자.\par

\subsection{슬라이드 48 · 삼포제의 실제 경작}
겨울곡물과 봄곡물은 파종·수확 시기가 달라 노동 수요도 분산한다. 휴한지는 가축 방목과 토양 회복에 쓰일 수 있다. 어떤 작물을 어떤 경포에 심을지는 공동 방목과 연결되므로 개인이 임의로 바꾸기 어렵다. 기술적 회전 방식과 공동 규칙이 결합되어야 장원 농업을 이해할 수 있다.\par

\subsection{슬라이드 49 · 직영지에서 지대 장원으로}
영주가 직영지를 부역으로 경작하는 비용이 커지면 토지를 임대하고 지대를 받는 방식이 유리해질 수 있다. 이 변화는 흑사병 이전에도 시장과 화폐 사용이 늘면서 진행되었다. 직영지의 해체가 농민의 완전한 자유를 즉시 뜻하지는 않는다. 부역이 현물·화폐 지대로 바뀌어도 사법권과 신분 의무는 남을 수 있다.\par

\subsection{슬라이드 50 · 상대가격이 계약을 바꿀 때}
노동이 희소해지면 임금 $w$가 지대 $r_T$에 비해 높아질 수 있다. 영주의 직접 경영 이익을 $pY-wL-C_{\mbox{감독}}$으로 쓰면, 노동비와 감독비가 올라갈수록 고정 지대 $R$을 받는 임대가 매력적이다. 그러나 영주와 농민의 교섭력, 법적 강제와 토지 시장이 실제 계약을 결정한다. 상대가격은 방향을 설명하고 제도는 속도와 분배를 설명한다.\par

\subsection{슬라이드 51 · 자유를 매각한다는 뜻}
영주는 농민의 의무를 일시금이나 정기 지대로 바꾸어 수입을 확보할 수 있고, 농민은 이동과 경영의 여지를 얻는다. “해방”은 모든 세금·지대가 사라진다는 뜻이 아니다. 어떤 권리를 얼마에 사들였는지, 부자유를 입증하는 재판권이 남았는지를 확인해야 한다. 지대 수취자로의 변화는 장원 경영 구조의 전환이다.\par

\subsection{슬라이드 52 · 등본 보유농의 권리}
Copyhold는 장원 재판 기록에 보유권을 적고 관습에 따라 상속·이용을 인정하는 형태다. 법률상의 완전한 자유보유권과 같지 않지만 농민에게 예측 가능한 권리를 줄 수 있었다. 기록이 남는다는 것은 거래와 분쟁 해결에 도움을 주지만, 갱신료나 관습 의무의 부담도 확인해야 한다.\par

\subsection{슬라이드 53 · 요먼의 성장}
Yeoman은 자영 농민 또는 상당한 경작권을 가진 농민을 가리키며, 모든 이가 똑같은 크기의 땅을 가진 것은 아니다. 흑사병 뒤 토지/노동 비율이 높아져 살아남은 농민의 교섭 여지가 커졌다. 그러나 토지 접근과 상속이 불평등하면 일부는 땅을 늘리고 다른 일부는 토지를 잃는다. 평균 임금 상승과 계층 분화를 함께 봐야 한다.\par

\subsection{슬라이드 54 · 농민층의 분화}
토지 시장과 인클로저는 큰 경영의 확장과 토지 없는 노동자의 증가를 동시에 만들 수 있다. 생산성이 올라도 그 잉여의 분배는 토지 소유와 협상력에 따라 달라진다. 자본가적 농업가와 임금 노동자가 생기는 과정은 순수하게 기술 변화만이 아니라 공유권의 재배분과 법적 강제의 결과이기도 하다.\par

\subsection{슬라이드 55 · 흑사병의 노동 충격}
대규모 사망으로 노동이 토지에 비해 희소해지면 노동의 한계생산과 협상력이 상승할 수 있다. 단순 Cobb–Douglas 모형 $Y=A T^\alpha L^{1-\alpha}$에서 경쟁적 실질임금은 $w/p=\partial Y/\partial L=(1-\alpha)A(T/L)^\alpha$다. $L$이 감소하면 이 값은 오른다. 그러나 실제 명목임금·식량가격·법적 임금 상한이 다르게 움직여 실질임금의 시점과 크기는 지역별로 달랐다.\par

\paragraph{노동 감소의 크기를 읽기}

토지·기술이 그대로라면 $w/p$는 $L^{-\alpha}$에 비례한다. 양변에 로그를 취하면 $d\ln(w/p)/d\ln L=-\alpha$이므로 노동이 1\% 줄 때 실질임금의 모형상 변화는 약 $\alpha$\% 상승이다. 예를 들어 $\alpha=0.3$이고 노동이 30\% 줄었다면 임금 배율은 $(0.7)^{-0.3}\approx1.11$이다. 이는 \textbf{경쟁적 고용과 변하지 않는 산출가격} 아래의 수치 예시다. 도시의 전염 위험, 강제 노동과 임금 규제는 관찰값을 이 예측에서 벗어나게 할 수 있다.\par

\subsection{슬라이드 56 · 실질임금 그래프}
명목임금이 올라도 식품 가격이 더 오르면 노동자가 살 수 있는 양은 줄 수 있다. 실질임금은 보통 $w/P$이므로 그래프의 가격지수와 기준 연도를 확인하자. 흑사병 직후의 단기 혼란과 여러 세대에 걸친 상승은 같은 현상이 아니다. 도시와 농촌, 남성과 여성, 숙련 여부에 따라 평균도 달랐을 수 있다.\par

\subsection{슬라이드 57 · 임금/지대 비율}
임금/지대 비율은 노동과 토지의 상대적 가격을 보여 준다. 인구가 줄어 $T/L$이 늘면 노동의 희소성은 커지고 토지의 희소성은 줄어드는 방향이 예상된다. 그러나 그림의 지대가 고정 현금 지대인지 산출 비율인지에 따라 인플레이션의 효과가 달라진다. 비율 상승을 바로 모든 농민의 소득 상승으로 읽지 말자.\par

\subsection{슬라이드 58 · 노동자 조례의 정치경제}
영주와 국가는 높은 임금과 이동을 제한하려 했다. 1351년의 노동자 조례는 시장에서 노동이 희소해졌다는 신호와 제도권의 대응이 충돌한 사례다. 법에 적힌 임금 상한이 실제 임금을 완전히 결정하는지는 집행력과 회피 가능성에 달렸다. 가격 통제가 왜 등장했는지와 효과가 무엇이었는지를 분리해서 물어야 한다.\par

\subsection{슬라이드 59 · 조례 문장을 읽는 법}
건전한 노동자의 취업 의무와 관습적 임금 상한은 노동의 자유로운 이동을 막으려는 조항이다. “공정한 임금”은 현대의 노동권 기준이 아니라 당시의 관습과 위계에 맞는 가격을 의미할 수 있다. 흑사병 이전 임금으로 돌아가려는 법의 의도와 계약 현장의 실제 지불액은 다를 수 있으므로 재판 기록과 장부를 대조해야 한다.\par

\subsection{슬라이드 60 · 폐촌의 여러 원인}
버려진 마을은 전염병 이후의 인구 감소를 보여 주지만 모두 흑사병 한 번으로 생긴 것은 아니다. 토지의 용도 변경, 경작지 통합, 기후·전쟁·시장 변화도 영향을 줄 수 있다. 폐촌의 위치와 시기를 비교하면 어느 설명이 가능한지 좁힐 수 있다. 한 사진의 흔적에서 전체 지역의 인과를 단정하지 말자.\par

\subsection{슬라이드 61 · 봉건 반동의 범위}
노동이 부족할 때 영주가 농민의 이동을 법과 강제로 막으려는 반응을 봉건 반동이라 부른다. 이를 모든 유럽에서 동일하게 성공한 정책으로 생각하면 안 된다. 영주의 집행력, 도시 임금 기회, 국가와 농민의 동맹이 다르면 결과가 달라졌다. 상대가격만큼 \textbf{누가 규칙을 만들고 집행하는가}가 중요하다.\par

\subsection{슬라이드 62 · 동부 독일의 농장 영주제}
동유럽의 Gutsherrschaft에서는 영주가 대농장을 운영하며 노동 의무와 재판권을 강화했다. 서유럽의 농민 해방과 대비되지만 “동유럽 전체”를 한 제도로 묶으면 지역 차이가 사라진다. 장거리 곡물 수출의 이익과 도시·국가 권력의 균형이 영주에게 유리한 조건을 만들었는지 살펴보자.\par

\subsection{슬라이드 63 · 엘베강의 경계}
엘베강은 서유럽과 동유럽의 다른 농업 제도를 비교할 때 쓰는 대략적 경계다. 지도상의 선이 실제 제도의 날카로운 단절을 뜻하지는 않는다. 강 동서의 도시 밀도, 토지 소유, 수출 시장과 정치 권력을 함께 비교해야 “같은 흑사병 뒤 왜 다른 대응이 나왔는가”에 답할 수 있다.\par

\subsection{슬라이드 64 · 농민 봉기의 배경}
더 나은 임금·이동의 기회가 생긴 농민이 옛 의무의 회복에 저항하면 분쟁이 발생한다. 봉기를 단순한 소득 상승의 기계적 결과로 보지 말고 세금, 재판권, 토지 권리와 종교·정치적 요구를 함께 읽자. 집단행동에는 조직과 위험 분담이 필요하므로 불만이 모두 봉기로 이어지는 것은 아니다.\par

\subsection{슬라이드 65 · 장원의 장기 변화 도식}
부역 중심 장원 → 현물·화폐 지대 중심 장원 → 농민 권리 확대라는 흐름은 유용한 지도다. 그러나 반동과 봉기가 끼어 있고 지역마다 순서와 속도가 다르다. 그림의 화살표가 “반드시 그렇게 된다”는 법칙인지 “자주 관찰된 경로”인지 구분하자. 토지의 사적 소유권도 단번에 생기지 않았다.\par

\subsection{슬라이드 66 · 서유럽과 동유럽이 갈라진 이유}
시장 발달은 농민에게 대체 고용과 현금 수입을, 영주에게 임대의 이익을 줄 수 있다. 반대로 영주가 수출 경로와 국가 권력을 통제하면 노동 강제가 유지될 수 있다. 같은 노동 희소성 충격도 \textbf{농민의 탈출 가능성과 영주의 집행력}이 다르면 다른 균형으로 간다. 이것이 단순 요소가격 설명을 제도 분석으로 확장하는 지점이다.\par

\subsection{슬라이드 67 · 시민혁명의 경제적 과제}
혁명은 왕과 의회를 바꾸는 정치 사건이면서 지대·독점·조세 특권을 재정의하는 경제 사건이다. 봉건적 특권을 없애도 토지의 불평등이 자동으로 사라지는 것은 아니다. 장벽이 없어지면 거래와 투자 선택은 달라지지만 누가 기존 토지를 소유했는지가 새 시장의 출발점을 결정한다.\par

\subsection{슬라이드 68 · 소유권의 재구성}
영주의 사법·조세 특권이 폐지되면 토지를 임대하는 지주의 경제적 권리와 사람을 강제하는 공권력이 분리된다. 오늘의 지주·소작 관계에서도 불평등은 가능하지만, 계약과 재판의 주체가 달라진다. “사유화”를 권리의 명확화로만 보지 말고 공동 이용권의 소멸과 분배 효과를 함께 검토해야 한다.\par

\subsection{슬라이드 69 · 분익 소작의 유인}
수확의 절반을 지대로 내는 계약이라면 농민은 한 단위 추가 산출에서 절반만 얻는다. 단순화해 생산 $Y(e)$와 노력비용 $c(e)$를 두면 농민의 선택은 $\max_e (1-s)Y(e)-c(e)$이고, 최적 조건은 $(1-s)Y'(e)=c'(e)$다. 자가 경영의 $Y'(e)=c'(e)$보다 노력이 적을 수 있지만, 수확 위험을 영주와 나누는 이익도 있다. 정률 지대의 순효과는 위험과 신용 제약에 달린다.\par

\paragraph{유인과 보험을 한 모형에서 보기}

농민의 몫을 $1-s$, 지주의 몫을 $s$라고 두자. $Y(e)=ae-be^2/2$, $c(e)=ke^2/2$이면 농민이 보유하는 추가 생산의 가치는 $(1-s)(a-be)$이고 추가 노력비용은 $ke$다. 같게 놓고 풀면 $e_{\mbox{분익}}=(1-s)a/[k+(1-s)b]$다. 지대가 없는 경우는 $e_{\mbox{자가}}=a/(k+b)$이므로 $s>0$에서는 전자가 작다. 그러나 수확 $Y$가 날씨 때문에 불확실하면 농민 수입 $(1-s)Y$의 분산은 $(1-s)^2\operatorname{Var}(Y)$로 줄어든다. 신용이 약하고 흉작의 생계 피해가 큰 사회에서는 낮은 노력 유인과 위험 분담의 편익을 함께 비교해야 한다.\par

\subsection{슬라이드 70 · 혁명 전 프랑스의 신분}
신분표는 법적 특권과 조세 부담이 소득 분배와 동일하지 않음을 보여 준다. 귀족·성직자 안에도 부의 차이가 있고 제3신분도 농민부터 상인까지 다양했다. 혁명 전의 재정 위기를 이해하려면 누가 얼마를 벌었는가뿐 아니라 \textbf{누가 세금을 면제받고 누가 의사결정을 막을 수 있었는가}를 봐야 한다.\par

\subsection{슬라이드 71 · 농업혁명이라는 명칭}
“농업혁명”은 작물·윤작·도구·토지 제도가 오랜 기간 바뀐 과정을 묶는 말이다. 곡물이 국제 교역품이 되려면 생산의 잉여뿐 아니라 저장·운송과 상업 신용이 필요하다. 영국이 어느 시기 순수출국이었는지를 말할 때에는 산업화가 진행되는 동안 수입 구조가 바뀌었음도 함께 보아야 한다.\par

\subsection{슬라이드 72 · 생산성과 도시 노동}
농업 노동생산성이 오르면 같은 식량을 더 적은 노동으로 만들 수 있어 비농업 노동을 부양할 여지가 생긴다. 그러나 농촌 노동자가 자동으로 공장에 이동하는 것은 아니다. 공업의 노동 수요, 도시 주거, 이동 비용과 가족의 네트워크가 필요하다. 농업 발전이 공업화의 필요조건인지 충분조건인지 구별해야 한다.\par

\subsection{슬라이드 73 · 제1차와 제2차 인클로저}
초기 인클로저는 양모·목양과 관련된 경우가 많았고, 뒤의 의회 인클로저는 경지의 재편과 개별 농법의 적용을 목표로 한 사례가 많다. 둘을 단지 “불법/합법”으로 나누면 실제 강제와 보상 구조를 놓칠 수 있다. 개방 경지가 얼마나 남았는지와 지역별 시기를 확인해야 인클로저의 생산성 효과를 추정할 수 있다.\par

\subsection{슬라이드 74 · 초기 인클로저의 그림}
울타리가 생기면 가축을 통제하고 작물을 선택하기 쉬워지지만 공동 방목권을 가진 가구는 사용권을 잃을 수 있다. 그림의 면적 변화가 곧 총생산 증가를 뜻하지는 않는다. 양모 수익과 식량 공급, 노동 수요, 토지 없는 사람의 이동을 함께 비교해야 한다.\par

\subsection{슬라이드 75 · 지도의 시기 비교}
세계사 지도에 표시된 영국의 도시·경지 변화는 산업화 전야의 공간적 조건을 보여 준다. 지도에서 농지의 통합과 도시 성장의 시간 순서를 확인하되, 둘이 같이 보인다는 사실만으로 어느 쪽이 원인인지 결론 내리지 말자. 식량 가격과 공업 임금, 운송망이 중간 경로다.\par

\subsection{슬라이드 76 · 장기 19세기라는 구분}
1789\~{}1914년을 하나의 “장기 19세기”로 묶는 것은 달력보다 혁명·산업·제국의 연속성을 강조한다. 시대 구분은 설명을 위한 분석 도구이므로 시작과 끝의 사건을 외우는 것만으로 충분하지 않다. 농업·공업·정치권력이 이 기간 어떻게 서로를 바꾸었는지 따라가자.\par

\subsection{슬라이드 77 · 의회 인클로저의 유인}
인구 증가, 도시 수요, 전쟁기의 곡물 가격 상승은 경지의 집약적 이용을 유리하게 만들었다. 통합된 땅에서는 개별 농가가 울타리, 배수, 작물 순환을 조정하기 쉽다. 그러나 법적으로 권리를 증명해야 보상받을 수 있었다면 문서 없는 공동 이용자는 손실을 볼 수 있다. 효율과 분배를 분리해 평가하자.\par

\subsection{슬라이드 78 · 『유토피아』의 양}
“양이 사람을 잡아먹는다”는 비유는 목양 인클로저가 사람의 거주와 경작을 밀어냈다는 사회적 비판이다. 문학적 비유를 통계 수치로 읽지 말고 당시의 갈등과 정치적 언어를 보여 주는 사료로 읽자. 실제 퇴거의 규모와 임금·고용 효과는 지역의 토지 문서와 인구 자료로 따로 검증해야 한다.\par

\subsection{슬라이드 79 · 인구와 도시의 변화}
도시 인구 비중이 올라갈 때 농촌 인구가 반드시 절대적으로 감소하는 것은 아니다. 총인구가 더 빨리 증가하면 농촌 인구도 늘면서 비중만 줄 수 있다. 도시화는 공업 노동 수요와 농업 식량 공급의 결합을 요구한다. 지도·그래프의 분모를 확인하면 인클로저와 이주의 관계를 정확히 읽을 수 있다.\par

\subsection{슬라이드 80 · 공장 노동의 출처 논쟁}
농민을 토지에서 쫓아낸 인클로저가 공장 노동자를 만들었다는 설명은 직관적이지만, 실제 지역 이동은 비용과 거리의 제약을 받았다. 인구 자체의 증가, 도시 출산과 다른 산업에서의 이동도 노동 공급원이었다. 원인별 기여를 따지려면 노동자의 출신지와 시기, 공장 위치의 기록이 필요하다.\par

\subsection{슬라이드 81 · 인클로저의 제도적 결과}
토지 통합은 경작 결정을 빠르게 할 수 있지만 중소농의 공동권 상실과 토지 집중을 동반할 수 있다. 농업 생산성이 올랐다면 누가 이익을 얻었는지—지주, 임차 농업가, 임금노동자, 식품 소비자—나누어 보자. 자본주의적 농업의 형성은 사유권 강화만이 아니라 임대차·고용·시장 판매의 확장까지 포함한다.\par

\subsection{슬라이드 82 · 휴한을 줄이는 기술}
휴한은 지력을 회복하는 한 방식이지만 생산하는 작물을 심지 못하는 면적의 기회비용이 있다. 사료 작물과 콩과 작물을 돌려 심으면 토양과 가축 사육을 함께 개선할 수 있다. 어떤 농법이 유리한지는 종자·노동·가축·비료와 시장 수요에 달려 있으며, 이름만 새로운 농법이 자동으로 수확을 늘리는 것은 아니다.\par

\subsection{슬라이드 83 · 노포크 윤작의 구조}
밀→순무→보리→클로버의 순환은 곡물, 사료, 토양의 질소 보충을 결합한다. 순무는 겨울철 가축 사료가 되고 클로버는 콩과식물로 토양에 질소를 더할 수 있다. 이 순환은 가축의 분뇨를 경지로 돌려 식량과 사료의 순환을 만든다. 수확량 개선은 지역별 토양과 실행 비용을 고려해 평가해야 한다.\par

\subsection{슬라이드 84 · 노포크의 위치와 확산}
지도나 사진에서 노포크를 찾으면 농법의 이름이 한 지역의 실험과 확산에서 왔음을 알 수 있다. 어떤 기술이 다른 지역으로 옮겨갈 때 토양, 강우, 가축 수요와 토지 경계가 같지 않다. 채택 지역을 그대로 복사하기보다 기술의 핵심 원리와 현지 제약을 나누어 이해하자.\par

\subsection{슬라이드 85 · 순무와 클로버의 보완성}
휴한지에 사료작물을 심으면 경작지의 활용률이 높아지고 겨울 가축의 생존율이 오를 수 있다. 더 많은 가축은 분뇨를 공급해 곡물 수확을 다시 높일 수 있다. 이 선순환은 농가가 작물과 가축을 함께 경영할 때 강해진다. 다만 순무·클로버의 노동 투입과 초기 종자 비용도 계산해야 순생산성 증가를 평가할 수 있다.\par

\subsection{슬라이드 86 · 겨울 사료와 고기 보존}
건초가 부족하면 겨울에 가축을 많이 도살해야 하고, 저장·보존 기술에 의존하게 된다. 관개 목초지와 사료작물은 가축의 계절적 공급 제약을 완화한다. 후추의 사용을 부패를 막는 단일 원인으로 설명하기보다 향신료의 맛·지위·거래와 보존 기술을 구분하자. 농법 변화는 식품 소비의 계절성도 바꿀 수 있다.\par

\subsection{슬라이드 87 · 농업에서 공업으로 가는 여러 경로}
농업 생산성은 도시의 식량을 공급하고 수출이나 수입 대체로 외화 제약을 완화할 수 있다. 농민의 저축이 산업 투자로 이동할 가능성도 있다. 하지만 농업의 잉여가 자동으로 공장 자본으로 바뀌지 않는다. 은행·조세·지주 소비와 투자 선택의 경로가 필요하다. 한 경로의 증거를 다른 경로의 증거로 대신하지 말자.\par

\subsection{슬라이드 88 · 식량 가격과 공업 임금}
노동자의 명목임금 $w$와 식품 가격 $P_f$를 구분하면 실질 구매력은 대략 $w/P_f$다. 식품 공급이 늘어 $P_f$가 낮아지면 같은 실질임금에서 기업의 명목임금 부담이 낮아질 수 있다. 그러나 노동시장 경쟁, 주거비와 비식품 가격도 영향을 준다. “싼 식량 → 싼 노동 → 공업화”는 가능한 경로이지 기계적인 법칙이 아니다.\par

\subsection{슬라이드 89 · 20세기 농업의 기계화}
트랙터와 콤바인은 화석연료와 고정자본을 써 노동시간당 산출을 높인다. 이는 18세기의 윤작·토지제도 변화와 다른 시기의 “농업 공업화”다. 대형 기계의 이익은 넓고 평탄한 경지에서 크지만 소규모 경지에서는 비용을 회수하기 어렵다. 에너지 가격과 환경 비용까지 넣어야 장기 효과를 평가할 수 있다.\par

\section{공업: 길드에서 공장 이전까지}
생산이 도시 길드의 규칙에서 농촌 가내공업, 선대제, 매뉴팩처로 옮아간 과정을 노동 통제, 거래비용, 규모의 경제로 읽는다.\par

\subsection{슬라이드 01 · 공업화 이전의 생산 조직}
공업의 역사는 기계 발명의 순서만으로 이해되지 않는다. 누가 원료와 도구를 소유하고, 노동을 어디서 조직하며, 완성품을 누가 팔았는가가 중요하다. 고대의 노예 수공업, 중세 도시 길드, 농촌의 선대제, 한 작업장의 매뉴팩처를 \textbf{소유·감독·시장 범위}라는 같은 축으로 비교하자.\par

\subsection{슬라이드 02 · 고대 그리스의 에르가스테리온}
노예 노동을 사용한 공방은 작업을 한곳에 모을 수 있었지만 자유 생산자의 시장 기회를 줄일 수 있었다. 고대 수공업의 “한계”를 시민이 생산을 경시했다는 문화 한 가지로만 설명하면 노예제의 법적 강제와 정치적 권리를 놓친다. 생산 효율과 시민 사이의 분배·지위 변화를 별도로 평가해야 한다.\par

\subsection{슬라이드 03 · 중세 도시의 등장 조건}
농업 잉여와 인구 증가가 상공인에게 식량을 공급하고, 교역로와 보호가 거래를 집중시켰다. “도시의 공기는 자유롭게 한다”는 격언은 도시가 모든 주민에게 완전한 자유를 주었다는 뜻이 아니다. 장원 밖의 상공업 기회와 일정 기간 거주 뒤의 법적 지위 변화라는 구체적 맥락에서 읽어야 한다.\par

\subsection{슬라이드 04 · 도시 기원설 세 가지}
로마의 기존 도시와 요새가 이어진 곳, 자연스러운 교역로에 시장이 생긴 곳, 영주가 세입을 위해 계획적으로 세운 곳이 모두 있었다. 한 이론만으로 모든 도시를 설명할 필요는 없다. 도시가 \textbf{교역 비용을 낮추는 곳}이면서 동시에 \textbf{보호와 과세의 정치 단위}였다는 공통 기제를 찾아보자.\par

\subsection{슬라이드 05 · 남유럽과 북서유럽의 도시}
장거리 무역 도시와 성·수도원 주변의 외곽 도시는 성장 출발점과 농촌 관계가 달랐다. 남유럽의 토지 소유 도시 귀족과 북서유럽의 상공인 자치 요구를 단순한 남북 문화 차이로 환원하지 말자. 어떤 재산권과 군사·사법 권한을 누가 보유했는지가 도시의 상업 규칙을 바꿨다.\par

\subsection{슬라이드 06 · 몽생미셸의 공간}
성벽과 수도원, 외곽 거주지는 방어·종교·교역의 기능이 결합한 중세 공간을 보여 준다. 사진의 현재 모습이 건립 당시와 완전히 같다고 볼 수는 없다. 사람이 모이면 거래 상대를 찾기 쉽지만 공간이 제한되어 혼잡과 지대 상승이 생긴다. 성벽 안팎의 위치가 법적 보호와 시장 접근에 차이를 만들었다.\par

\subsection{슬라이드 07 · 보마리스 성의 기능}
성은 외부 공격을 막는 군사 시설이면서 주변 지역을 통제하는 권력의 거점이다. 안전이 좋아지면 상인이 모일 유인이 생기지만, 성을 가진 권력자는 통행료와 세금을 거둘 수 있다. 보호와 착취의 가능성이 같은 권력에서 나온다. 도시의 경제사는 건축 사진 뒤의 세입·권리 구조까지 읽어야 한다.\par

\subsection{슬라이드 08 · 카르카손의 성벽}
이중 성벽과 좁은 출입구는 방어와 통행의 통제를 모두 가능하게 한다. 도시의 독점권이나 시장 규칙이 작동하려면 물리적 경계와 법적 집행이 필요했다. 현대의 열린 도시 시장을 과거의 도시로 그대로 투영하지 말고, 누가 언제 성문을 통과하고 거래할 수 있었는지 묻자.\par

\subsection{슬라이드 09 · 부르주아와 도시 자치}
성벽 바깥의 상공인 거주가 늘면서 새로운 방어벽과 시민권 요구가 생겼다. bourgeois는 특정 시기의 도시 거주·재산·권리와 연결된 말이며 현대의 “부자”나 “자본가”와 항상 같지 않다. 영주의 허가와 통행세에 의존하던 상인은 자치권을 얻어 규칙을 스스로 정하려 했다. 자치는 새로운 독점의 가능성도 만들었다.\par

\subsection{슬라이드 10 · 도시 경제정책의 두 얼굴}
대외적으로는 주변 시장을 묶어 독점 이익을 얻고, 대내적으로는 회원 사이의 가격·품질·거래 시간을 규제했다. 대내적 “평등”은 도시의 모든 주민에게 같은 권리가 있었다는 뜻이 아니라 길드 회원 사이의 경쟁 제한에 가까웠다. 소비자와 외부 상인의 입장에서는 비용이 높아질 수 있다.\par

\subsection{슬라이드 11 · 금제권과 도시의 배후지}
Bannmeile는 도시 주변에서 특정 거래를 도시에 집중시키는 권리다. 거래가 모이면 정보와 안전의 이익이 있지만, 농촌 생산자는 더 낮은 가격을 받거나 먼 시장에 접근하지 못할 수 있다. 도시가 교환의 중심이라는 사실과 도시의 독점권이 사회 전체 후생을 높였다는 주장은 다르다.\par

\subsection{슬라이드 12 · 공정 가격과 거래 규제}
매점매석 금지와 공정 가격은 식량 부족의 불안을 낮출 수 있지만, 가격을 억누르면 다른 지역에서 물건을 가져올 유인이 약해질 수 있다. 당시의 “공정”은 공동체가 인정한 질서와 생계 기준을 뜻하며 현대의 경쟁 균형가격과 다르다. 가격 규칙을 평가하려면 공급 안정, 품질과 배제된 거래의 비용을 함께 보아야 한다.\par

\subsection{슬라이드 13 · 길드의 양면성}
길드는 회원의 안전·정보·교섭을 도우면서 외부인의 진입을 막았다. 동료의 품질을 감시하고 계약 불이행을 제재하면 거래비용이 줄지만, 독점과 높은 가격은 소비자에게 비용을 줄 수 있다. 동일한 조직에 \textbf{보험·정보의 편익과 지대 추구의 비용}이 함께 있으므로 시대별 시장 조건을 비교해야 한다.\par

\subsection{슬라이드 14 · 길드의 보험 기능}
장거리 무역은 도난·난파·상대방 불이행의 위험을 안았다. 길드가 손실을 나누고 회원에게 신용을 제공하면 개별 상인이 감당하기 어려운 거래가 가능해진다. 반대로 위험 분담의 혜택을 회원만 받게 해 외부 상인을 배제할 수 있다. 생산량 통제라는 기능과 위험 공동부담이라는 기능을 함께 읽자.\par

\subsection{슬라이드 15 · 종합 길드에서 직종 길드로}
처음에는 상인과 수공업자가 한 공동체에 묶였지만 도시의 규모와 전문화가 커지며 직종별 이해관계가 달라졌다. 직종 분리는 품질 기준을 구체화할 수 있고 상인 길드의 권력에 저항하는 방식도 되었다. 동시에 직종마다 진입 장벽이 생겨 기술의 이동이 어려워질 수 있다.\par

\subsection{슬라이드 16 · 상인 길드의 자치와 제재}
국가 법원이 약한 곳에서 상인들은 회원의 평판을 공유하고 거래를 어긴 사람을 함께 제재했다. 반복 거래에서 속임수의 일회 이익 $G$가 미래 거래 손실의 현재가치보다 작으면 협력이 가능하다. 할인인자 $\delta$와 매기 협력 이익 $R$를 쓰면 단순 조건은 $G\leq \delta R/(1-\delta)$다. 길드의 제재가 유용할수록 동시에 회원 밖 사람을 배제할 힘도 커진다.\par

\paragraph{반복거래의 조건을 도출하기}

$G$는 \textbf{오늘 정직하게 거래하는 경우와 비교한} 일회 속임수의 추가 이익이다. 적발되면 다음 기부터 회원 거래를 잃고, 그 순이익이 매기 $R$이라고 하자. 내년 이익의 현재가치는 $\delta R$, 그다음 해는 $\delta^2R$다. 무한등비급수를 더하면 미래 손실은 $\delta R/(1-\delta)$이므로 오늘의 속임수를 포기할 조건이 위 식이다. 예를 들어 $\delta=0.8$, $R=10$이면 미래 손실은 40이다. 여기에는 위반을 관찰하고 배제 조치를 집행할 수 있다는 중요한 전제가 들어간다.\par

\subsection{슬라이드 17 · 수공업 길드의 독립}
수공업자는 품질과 훈련을 직접 통제하려 했고 상인 귀족과 정치적 갈등을 겪었다. 직종 길드의 성장은 도시 시민의 권리를 넓힌 측면이 있지만 견습공·여성·이주민을 모두 같은 권리로 받아들인 것은 아니다. 길드의 수와 직종의 세분화는 경제의 전문화를 보여 주되 독점 범위도 함께 늘린다.\par

\subsection{슬라이드 18 · 도제에서 장인으로}
도제는 낮은 임금을 받고 기술을 배우며, 직인은 다른 작업장에서 숙련을 쌓고, 장인은 독립 영업권을 얻는다. 이 경로는 인적자본 투자의 비용을 나누지만 장인 자격을 제한하면 노동 공급을 인위적으로 줄인다. “학위 과정”과의 비유는 단계적 훈련을 이해하는 데 도움이 되지만 법적 영업권과 현대 학위의 효과는 다르다.\par

\subsection{슬라이드 19 · 수업과 편력의 의미}
문학 작품의 수업시대·편력시대는 숙련자가 한 곳의 관습만 배우지 않고 다른 작업장과 도시를 경험했다는 문화를 비춘다. 실제 장인 이동에는 길드 규칙과 비용이 따랐다. 여행을 무조건 자유로운 관광처럼 읽지 말고, 기술 전파와 일자리 탐색, 자격 인정의 문제로 연결하자.\par

\subsection{슬라이드 20 · 이동의 여러 말}
Trip·travel·tour·journey는 기간과 목적이 다르다. 직인의 편력은 새로운 작업장과 기술을 찾아 다니는 노동 이동에 가깝다. 이동을 한 번의 출발과 도착으로만 측정하면 중간의 학습과 정보망을 놓친다. 공업 기술은 책과 도구뿐 아니라 숙련자가 이동하며 전파되었다.\par

\subsection{슬라이드 21 · 도제제도의 순효과}
훈련 규칙은 초보자의 품질을 보장하고 기술을 전수할 수 있다. 반면 장인 수를 제한하거나 새로운 공정을 금지하면 혁신과 진입을 막는다. 오늘의 자격증을 무조건 길드와 같다고 보지 말고 정보 비대칭의 정도, 소비자 피해, 경쟁 진입의 가능성을 비교해 평가해야 한다.\par

\subsection{슬라이드 22 · 중세 섬유공업 지도}
섬유 생산에는 원료 공급, 방적, 직조, 마감과 장거리 판매가 필요하다. 지도에서 생산 중심지가 항구·하천·양모 산지와 어떻게 연결되는지 보자. 한 도시가 모든 공정을 내부에서 수행했다는 뜻은 아니다. 작업이 분산될수록 상인은 품질과 납기를 조정할 능력이 필요했다.\par

\subsection{슬라이드 23 · 15세기의 위기와 기회}
중세 후기의 장원·도시 질서가 약해지면서 일부 자영농과 농촌 수공업자는 새로운 기회를 얻었다. “침체”는 모든 사람의 소비가 동시에 악화했다는 뜻이 아니다. 토지와 노동의 상대가격 변화는 영주 수입을 줄이면서 살아남은 농민의 임금을 올릴 수 있었다. 계층별 경험을 나누어 읽어야 한다.\par

\subsection{슬라이드 24 · 민부와 소상품 생산자}
자기 도구를 가진 농민·수공업자가 시장을 향해 생산하고 잉여를 축적하면 기존의 영주·길드 통제 밖에서 자본의 씨앗이 생긴다. “민부”는 모든 민중이 동일하게 부유해졌다는 뜻이 아니라 생산수단을 가진 소생산자 층의 성장이라는 맥락에서 읽자. 누가 원료와 판매망을 통제했는가에 따라 독립성은 달라진다.\par

\subsection{슬라이드 25 · 선대제의 계약 구조}
Putting-out system에서는 상인이 원료·선금을 주고 여러 가내 작업자에게 공정을 맡긴 뒤 완성품을 회수해 판매한다. 공장은 한곳에 없어도 \textbf{자금과 시장의 조정}이 상인에게 집중된다. 지역 노동이 풍부하고 건물·감독 비용이 클 때 분산 작업은 유리하지만 품질과 납기를 확인하는 거래비용이 생긴다.\par

\subsection{슬라이드 26 · 선대제의 현금과 물건 흐름}
상인 → 작업자에게 원료·도구·선금이 가고, 작업자 → 상인에게 완제품이 돌아온다. 성과급이 지급되면 작업자의 수입은 생산량에 연결되지만 원료와 판로를 상인이 장악하면 가격 교섭력은 약하다. 그림의 화살표를 돈·물건·노동 명령으로 구분하면 “독립된 작업장”과 “상인에 대한 경제적 종속”이 동시에 보인다.\par

\subsection{슬라이드 27 · 분산 공정의 그림}
작업이 여러 집에 흩어져 있으면 건물과 고정 설비 비용을 줄일 수 있지만 각 공정 사이의 이동과 검사 비용이 증가한다. 공정의 순서가 길어질수록 한 집의 지연이 전체 납품을 늦추는 병목이 된다. 그림의 각 작업자를 하나의 공장 노동자로 취급하지 말고, 누가 도구를 소유하고 어떤 계약으로 연결되는지 보자.\par

\subsection{슬라이드 28 · 길드 분화에서 상인 자본으로}
부유한 길드 구성원이 원료·판매망을 장악하고 가난한 구성원에게 주문을 주면 법적으로 독립된 장인도 경제적으로 종속될 수 있다. 성과급은 노동의 강도를 높이는 유인이지만 품질 저하와 위험 전가를 낳을 수 있다. 길드의 쇠퇴가 곧 자유로운 경쟁의 출현이라는 단순 그림은 피해야 한다.\par

\subsection{슬라이드 29 · 계절적 노동의 활용}
농업의 파종과 수확 사이 유휴 노동이 있으면 가구는 방적·직조를 겸업할 수 있다. 상인은 연중 공장 노동자를 고용하는 대신 필요한 공정을 맡기므로 고정 임금 부담을 줄인다. 그러나 농번기에는 노동 공급이 줄어 납기가 불안정해진다. 원료가 저렴하고 작업 도구가 간단할수록 이 방식이 잘 맞는다.\par

\subsection{슬라이드 30 · 영국의 클로디어}
Clothier는 양모를 사서 방적·직조·축융 등 공정을 연결하고 천을 판매한 조직자다. 같은 이름 아래에도 자신의 직기로 일부 생산한 소규모 경영과 작업을 거의 모두 외주한 대규모 경영이 있었다. 상인의 자본 역할과 생산 현장의 소유 구조를 지역별로 나누어 비교해야 한다.\par

\subsection{슬라이드 31 · 북부의 소규모 경영}
작은 클로디어가 가족 노동과 농업을 병행하면 수입원을 분산할 수 있다. 여성과 아동의 방적 노동은 가구 소득에 중요했지만 기록에 낮게 평가되거나 빠질 수 있다. 자기 직기를 가진 생산자가 이웃 작업을 외주하는 형태는 단순히 “고용주 대 노동자”의 이분법에 맞지 않는다.\par

\subsection{슬라이드 32 · 동부·서남부의 대규모 경영}
상인은 양모를 여러 작업자에게 차례로 넘기고 완성품을 회수한다. 직접 방적하지 않아도 원료 구입과 유통, 재고 위험을 부담하므로 조직 능력이 중요하다. 작업자가 많이 분산될수록 품질 표준과 검사, 신용이 핵심 비용이 된다. 지역 차이는 노동 공급과 시장 접근, 자본의 규모를 반영할 수 있다.\par

\subsection{슬라이드 33 · 직기와 여성 노동}
그림의 직기 기술은 산업혁명 이전에도 기계적 보조가 존재했음을 보여 준다. “공업화 이전”이 곧 도구 없는 손작업을 뜻하지 않는다. 방적·직조·마감의 각 단계에 여성과 아동의 노동이 어떻게 배치되었는지 살펴야 하며, 그림 속 한 작업장의 모습으로 전체 지역의 평균 생산성을 단정하지 말자.\par

\subsection{슬라이드 34 · 모직물 공정의 연결}
세척한 양모를 빗질하고 실을 뽑아 직조한 뒤, 축융·기모·전모·염색으로 품질을 만든다. 한 단계의 산출은 다음 단계의 투입이다. 각 공정의 생산속도가 다르면 느린 단계가 전체 산출을 제한한다. 작업의 장소를 한곳으로 모을지 분산할지는 병목·운송·감독 비용의 비교로 이해할 수 있다.\par

\subsection{슬라이드 35 · 프로토공업화란}
Protoindustrialization은 증기기관 이전에 농촌에서 원격지 시장을 위한 상품 생산이 크게 늘어난 현상을 가리킨다. 단순 자급 부업과 다른 점은 판매 시장의 범위와 상인의 조직 역할이다. 이를 산업혁명의 필연적 “전 단계”로 만들면 지역별 실패·중단 사례를 설명하기 어렵다.\par

\subsection{슬라이드 36 · 반농반공의 경제적 효과}
농가의 현금 소득이 늘면 토지 부족 속에서도 가구 형성과 소비가 가능해질 수 있다. 반면 더 많은 자녀와 노동 공급이 1인당 소득 상승을 상쇄할 수도 있다. 장기 시장을 향해 일하면 품질·납기·숙련을 배우지만, 모든 농촌 공업 지대가 기계 공장으로 이행한 것은 아니다.\par

\subsection{슬라이드 37 · 농촌 공업의 형성}
도시 길드의 진입 규칙을 피하고 저렴한 농촌 노동을 이용하면 상인은 생산을 분산할 유인을 갖는다. 농촌 작업자는 농업과 제조업을 결합해 계절 위험을 나눌 수 있다. 장원 제약의 완화, 시장 도로, 원료 공급이 함께 필요하며 “농촌”이라는 위치만으로 자본주의 여부가 정해지지 않는다.\par

\subsection{슬라이드 38 · 매뉴팩처의 위치}
Manufacture는 수작업을 중심으로 사람을 한 작업장에 모은 생산 방식이다. 가내수공업, 선대제, 공장제기계공업 사이의 구분은 기계의 유무만이 아니라 작업 장소와 노동 지휘의 차이를 포함한다. 산업혁명 전에도 규모가 큰 작업장이 있었고 기계 공장 이후에도 가내 작업은 남았다.\par

\subsection{슬라이드 39 · 분업에 기초한 협업}
한 노동자가 모든 공정을 하는 대신 공정별로 전문화하면 반복을 통해 숙련과 속도가 늘 수 있다. 그러나 노동을 한곳에 모으려면 건물, 감독자, 원료 재고와 안정적 주문이 필요하다. 소규모 매뉴팩처는 선대제와 경쟁하면서도 일부 공정을 외주하여 \textbf{혼합 조직}으로 운영될 수 있다.\par

\subsection{슬라이드 40 · 경영 주체의 차이}
농민적·상인적 매뉴팩처라는 구분은 누가 생산 기술과 유통 경로를 통제하는가에 초점을 맞춘다. 생산 이윤과 유통 이윤이 실제 사업에서 완전히 분리되는 것은 아니다. “진보적/보수적”이라는 평가는 특정 이론의 해석이므로, 두 유형의 투자·혁신·노동 조건을 실증적으로 비교해야 한다.\par

\subsection{슬라이드 41 · 숙련과 규율}
한 공간에서 반복 작업을 하면 도구와 동작을 표준화하고 감독하기 쉽다. 이것이 품질을 높일 수 있지만 작업 속도의 강제와 노동자의 자율성 축소도 가져온다. 노동 생산성의 증가는 근로시간, 숙련의 유지, 불량률을 함께 측정해야 한다. 매뉴팩처의 경험은 이후 기계 공장의 작업 조직에 영향을 주었다.\par

\subsection{슬라이드 42 · 국가가 매뉴팩처를 권장한 이유}
국가는 조세 수입, 군수품, 수입 대체와 일자리 확보를 위해 작업장을 장려했다. 하지만 구빈원·감옥의 강제노동과 자유로운 임금 고용을 같은 “고용 확대”로 평가할 수 없다. 콜베르의 정책을 이해하려면 보호·보조금의 비용을 누가 지고 제품을 누가 샀는지, 국제수지 목표와 실제 생산성 효과를 분리해야 한다.\par

\subsection{슬라이드 43 · 구빈원 이미지}
『올리버 트위스트』의 작업장 장면은 빈곤층에 대한 규율과 강제의 사회적 인식을 보여 준다. 문학적 장면은 당시 모든 구빈원의 정확한 노동 조건을 재는 통계가 아니다. 정책의 목표가 빈곤 구제, 노동 통제, 재정 절약 중 무엇이었는지 확인하고 빈곤층의 선택 가능성을 함께 살펴야 한다.\par

\subsection{슬라이드 44 · 매뉴팩처의 장단점}
집중 생산은 공정 간 이동 시간을 줄이고 표준화·감독·분업을 돕는다. 반면 건물과 고정 설비, 안정적 임금 지급이 필요하다. 두 방식을 비교할 때에는 먼저 일회성 고정비와 제품을 하나 더 만들 때 드는 가변비를 구별한다. 주문 규모가 커져 고정비를 더 많은 제품에 나눌수록 집중 방식의 단위 비용이 낮아질 수 있다.\par

\paragraph{생산 규모가 조직을 바꾸는 문턱}

같은 품질의 제품을 $Q$개 만든다고 하자. 집중 생산의 고정비는 $F$, 제품당 임금·내부 조정 비용을 합한 가변비는 $c_c$다. 분산 생산의 제품당 외주·운송·검사비 합계를 $c_d$라고 놓으면 $C_c(Q)=F+c_cQ$, $C_d(Q)=c_dQ$다. $c_c<c_d$라면 두 방식의 비용이 같아지는 주문량은 $Q^*=F/(c_d-c_c)$다. $Q>Q^*$에서 집중이 싸진다. 그러나 특정 공정의 숙련이 농촌 가구에 쌓여 있거나 수요가 계절마다 출렁이면 $c_d$가 낮아지거나 유휴 설비 부담이 커져 외주가 계속 유리할 수 있다.\par

\subsection{슬라이드 45 · 왜 매뉴팩처만으로 끝나지 않았나}
공정을 한곳에 모아도 단계별 처리 속도가 다르면 대기와 재고가 생긴다. 일정한 비율의 노동자를 모두 고용하기 어려우면 일부 공정을 선대제로 외주한다. 그래서 역사적 매뉴팩처와 가내공업은 경쟁하면서 공존했다. 다음 산업혁명에서 동력기계와 공장 조직이 결합할 때 어느 병목이 해소되는지 연결해 보자.\par

\end{document}
